\section{The corner state sum \texorpdfstring{$C$}{C}}\label{sec:statesum}

Except where an explicit larger domain is stated, $P$ is generic
throughout this section. The definitions follow the main
text's Section~VI and use the positive lifts of the actual carrier diagrams.

\subsection{Decompositions and subpolygons}

\begin{definition}[decomposition]\label{def:decomposition}
A \emph{decomposition} of $P$ is an independent set $S\in\Ind(G_P)$ of
crossings (Definition~\ref{def:interlace}): no two crossings of $S$ interlace.
\end{definition}

\begin{definition}[oriented smoothing and subpolygons]\label{def:smoothing}
Let $S\in\Ind(G_P)$. Cutting the traversal circle $\Gamma(P)$ at the two
visits of each $x\in S$ and reconnecting the arc arriving at one visit of $x$
to the arc leaving the other visit of $x$ (for both visits) partitions
$\Gamma(P)$ into closed cycles; the \emph{subpolygons} (or \emph{carriers})
of $S$ are the closed polygonal curves in $\RR^2$ traced by these cycles. Each
subpolygon $Q$ is a closed polygon whose corners are the vertices of $P$ it
passes through and the crossing points of $S$ it turns at (its
\emph{smoothing corners}); at a smoothing corner between $E_i$ and $E_j$ it
arrives along one of $\lt_i,\lt_j$ and leaves along the other. The
\emph{crossings of $Q$} are the crossings $x\in X(P)\setminus S$ both of whose
visits lie on $Q$; their number is $m_Q$. A crossing in $N(S)$ has its two
visits on different subpolygons and is a crossing of no subpolygon.
\end{definition}

\begin{convention}[ownership of a selected visit]\label{conv:selected-visits}
At a selected crossing with traversal visits $a,b$, reconnect the incoming
arc at $a$ to the outgoing arc at $b$, and the incoming arc at $b$ to the
outgoing arc at $a$. Assign the original visit $a$ to the resulting cycle
containing the incoming one-sided arc at $a$; assign $b$ by the analogous
rule. Unselected visits retain their ordinary traversal ownership.
Thus a selected visit is assigned to one abstract carrier, even though
the two carriers through that smoothing site have the same point of the
plane in their polygonal images. This is the \emph{incoming-visit convention}.
It specifies endpoint ownership only; the geometric reconnection is unchanged.
\end{convention}

\begin{remark}[why the convention exists]\label{rem:selected-visits}
Clause~(iv) of the lemma that follows quantifies over \emph{all} crossing
visits, including the two visits of a selected crossing, and both carriers
through a smoothing site pass through the same point of the plane; without
an ownership rule for those two visits the clause is not well posed. The
convention supplies that rule and nothing else: it changes no support, turn
or retained-crossing datum, and the auditor bench's referee of decision item
D-25-1 (phase ledger) reports that clause~(iv) holds under either ownership
rule on its finite test set, so the choice is a disambiguation, not a
mathematical commitment. The main text has no counterpart of this
convention.
\end{remark}

\begin{lemma}[subpolygons and their visit assignment]\label{lem:carriers}
Let $P=(\mu_1,\ldots,\mu_n)$, $n\geq3$, be a generic polygon: for every
three distinct indices the vertex determinant is nonzero, and no three
edge interiors concur.
Use one-based cyclic indices, directed edges
$E_i=[\mu_i,\mu_{i+1}]$, directions $d_i=\mu_{i+1}-\mu_i$, and turns
$\tau_i=\operatorname{sgn}\det(d_{i-1},d_i)$.
Let $X(P)$ be its set of transverse crossings of remote edges. The
traversal circle has two visits for each crossing, in the order of the
oriented traversal. Crossings interlace when their two pairs of visits
alternate on that circle. Let $S\subseteq X(P)$ be independent: no two
elements of $S$ interlace. Perform the above oriented reconnections,
tracing the resulting cycles by inherited straight subsegments, with
corners at original vertices and selected smoothing sites. Call these
polygonal cycles the carriers, and use the incoming-visit convention.
Then:
\begin{enumerate}
\item[(i)] There are exactly $|S|+1$ carriers, independent of the order
of reconnections. Each carrier traverses the visits assigned to it in
their cyclic order inherited from the original traversal circle.
The two visits of a selected crossing belong to different carriers.
\item[(ii)] Every carrier has nonzero edges, at least three corners,
and no antiparallel consecutive directions. All its corner turns are
nonzero. At an original vertex $i$ the turn sign is $\tau_i$.
At a selected crossing of $E_i,E_j$, the two smoothing corners have
signs $\operatorname{sgn}\det(d_i,d_j)$ and
$\operatorname{sgn}\det(d_j,d_i)$, one left and one right.
\item[(iii)] A carrier's self-intersections are exactly the unselected
crossings both of whose visits are assigned to it. They are transverse,
none is a corner, and no carrier has a triple point. An unselected
crossing interlacing some element of $S$ has its visits on different
carriers; an unselected crossing interlacing no element of $S$ has
both visits on one carrier.
\item[(iv)] The assignment of \emph{all} crossing visits is noncrossing.
There are no four distinct visits $u_1$, $u_2$, $u_3$, $u_4$ in that cyclic order
with $u_1,u_3$ on one carrier and $u_2,u_4$ on a different carrier.
This includes selected visits under the stated convention.
\end{enumerate}
\status{proved (refereed: bench A, 2026-09-05T19:48:09Z; transcribed from CV lem:carriers, lem:carrierword (d7\_anchors, d6\_vertexedge); RC def:canonical-total) for clauses (i)--(iii); clause (iv), ``including every selected visit'', is a new argument of round~1 with no frozen-source locator, cleared by bench~A on SM4 (2026-09-05T19:48:09Z) and listed in Section~\ref{sec:newlist} (F-25-133, F-25-134; with Convention~\ref{conv:selected-visits}, C005)}
\end{lemma}

\begin{proof}
\emph{Finite successor model.}
Mark the traversal circle at every original vertex and every crossing
visit. This is a nonempty finite cyclic list, even if there are no
crossings. Each consecutive pair bounds an oriented straight
subsegment. Write $\rho$ for the successor permutation of the list.
At the two distinct marked visits $a,b$ of a selected crossing, replace
the two arrows by
\begin{equation}\label{carrierpr:successors}
\rho_S(a)=\rho(b),\qquad \rho_S(b)=\rho(a).
\end{equation}
At an unselected marked point keep its outgoing arrow. Selected
crossings have disjoint pairs of marked visits. Thus these swaps
operate on disjoint outgoing slots, commute, and give a permutation
$\rho_S$ independent of their order. Its cycles are precisely the
oriented reconnections: the node $a$ retains its incoming arc and
leaves along the old outgoing arc at $b$. This also implements the
incoming-visit convention explicitly.

\emph{Splitting and induced order.}
We prove by induction as selected pairs are processed that the cycles
traverse their marked nodes in the cyclic order inherited from the
original list, and that the two endpoints of each unprocessed selected
pair lie on one current cycle. Initially there is the one original
cycle. Suppose a current selected pair has endpoints $a,b$ on a cycle
$L$. In its inherited cyclic order write that cycle as
\begin{equation}\label{carrierpr:split-list}
(a,A_1,\ldots,A_p,b,B_1,\ldots,B_q).
\end{equation}
Either of the lists between the endpoints may be empty. Swapping the
outgoing arrows gives exactly the two cycles
\begin{equation}\label{carrierpr:split-cycles}
(a,B_1,\ldots,B_q),\qquad (b,A_1,\ldots,A_p).
\end{equation}
Each retains the induced cyclic order. Every remaining selected pair
has both endpoints in one of the two open arcs between $a,b$ on the
original circle, since it does not interlace this selected pair.
If it was on $L$, its two endpoints therefore enter the same new cycle;
if it was on another cycle, it is unaffected. This proves the inductive
claims and shows that every selected reconnection splits one cycle
into two, never merging cycles. Hence the final count is $|S|+1$.
The endpoints $a,b$ enter different cycles in
\eqref{carrierpr:split-cycles}, and all later operations only split.
They remain on different final carriers. This proves (i).

\emph{Noncrossing away from selected endpoints.}
The preceding induction also gives noncrossing ownership of the
unmarked open traversal arcs. Here are the details, including arbitrary
points on those arcs. For no selected pairs there is one carrier, so
there is nothing to check. In a splitting step the one old carrier
$L$ is divided between the two open arcs cut out by $a,b$, with the
endpoints omitted for this paragraph. All other carrier sets are
unchanged. If two new carrier blocks alternated with four points,
and at most one of them came from splitting $L$, the same four points
would violate the previous partition's noncrossing property, using
$L$ for that new block. If both came from splitting $L$, their points
would alternate between the two open arcs of $a,b$. The four disjoint
cyclic transitions between these four points would then each contain
one of $a,b$. Two endpoints cannot lie in four disjoint transitions.
This is impossible. Induction proves noncrossing for all traversal
points away from preimages of $S$.

\emph{Including every selected visit.}
Assume four distinct crossing visits violated (iv). Move any of these
that are selected a sufficiently small distance backwards along the
original traversal circle; leave the unselected ones fixed. There are
finitely many selected visits and four distinct chosen points, so the
distances can be chosen positive and small enough that no selected
endpoint is passed and the cyclic order of the four points is unchanged.
By the incoming-visit convention, each moved point has exactly the
carrier assignment of its original selected visit. Now all four points
avoid preimages of $S$ and give the forbidden alternation of the
previous paragraph. This contradiction proves the all-visits assertion.

\emph{Ownership of unselected crossings.}
Let $x\notin S$. If it interlaces no selected pair, its two endpoints
start on the same original cycle and remain together after every split:
at a split whose carrier contains them, noninterlacement puts them in
one open arc of that selected pair. Thus they end on one carrier.
If $x$ interlaces a selected pair $s$, process $s$ first, which is
permitted by order independence. Its split puts the two visits of $x$
on different cycles. Later reconnections only split, so they can never
be reunited. This proves both ownership assertions in (iii).

\emph{Nonzero segments and corner signs.}
Genericity makes all original vertices distinct: coincident vertices
would give a zero triple with any third index. It also makes every
edge nonzero and every original turn nonzero. No vertex lies on a
nonincident closed edge: such an incidence would give a zero triple.
Two remote edges meeting at a point have neither endpoint there, and
their directions cannot be parallel, since coincident lines would
give collinear original vertices. Thus their intersection is an
interior transverse crossing. Distinct crossings have distinct points,
as a coincident pair of crossing points would lie on at least three
edge interiors, contrary to genericity.

Consequently consecutive marked points on an original edge are
distinct in the plane, with strictly increasing parameters; their
connecting subsegment has positive length. At a selected node $a$,
the point of the plane agrees with its twin $b$, but its new successor
is the old successor of $b$, which is at positive distance along
$b$'s outgoing subsegment. Equation~\eqref{carrierpr:successors}
does not insert a zero segment from $a$ to $b$.
Omitting unselected crossing marks joins successive positive subsegments
of the same original edge in the same direction. Hence every edge
between consecutive carrier corners is nonzero and inherits an
original oriented direction up to a positive scalar.

At an original vertex $i$ no reconnection occurs. The incoming and
outgoing carrier directions are positive multiples of $d_{i-1}$ and
$d_i$, so their determinant has sign $\tau_i$. At a selected crossing
of $E_i,E_j$, the node assigned by incoming $E_i$ leaves along $E_j$;
its determinant is a positive multiple of $\det(d_i,d_j)$. The twin
node has a positive multiple of the opposite determinant
$\det(d_j,d_i)=-\det(d_i,d_j)$. Transversality makes these nonzero,
with opposite signs. These are all corner types. In particular none
is antiparallel or flat.

Each cycle contains a corner: a cycle having only unselected visits
would follow the unchanged original successor around the full parent
traversal, which contains original vertices. Thus its nonzero segments
form a closed polygonal cycle with the indicated corners. One nonzero
segment cannot close. A closed polygon of exactly two nonzero segments
has opposite segment vectors and hence antiparallel corners, already
excluded. Every carrier therefore has at least three corners. This
proves (ii), including membership in the regular locus.

\emph{Exactly the retained self-intersections.}
Every traced carrier segment is an original straight subsegment with
its inherited parameter direction. Distinct traversal subsegments of
one original edge have disjoint interiors, so no overlap is introduced.
Original adjacent edges meet only at their common original vertex,
which has one traversal occurrence and lies on one carrier. All other
possible intersections are original crossings. At a selected crossing
the two incoming visits belong to different carriers by (i); a single
carrier therefore passes through that smoothing point only once.
An unselected crossing is visited twice by a carrier precisely when
both of its original visits are assigned to that carrier. Otherwise
it is a meeting of two distinct carriers, not a self-intersection of
either. Every retained self-crossing is still transverse with its
unchanged straight germs. It is not a corner because original vertices
and selected crossing sites are all distinct from it. No carrier can
have a triple point, since no point of the parent has three traversal
preimages and each selected site is used only once by a given carrier.
These statements prove the remaining assertions of (iii).
\end{proof}

\begin{definition}[uniform subpolygons; retained data]\label{def:uniform}
A subpolygon $Q$ is \emph{uniform} if all its turns have one sign and
\emph{mixed} otherwise. A decomposition $S$ is \emph{uniform} if all its
subpolygons are. For a subpolygon $Q$ write $r_Q=\rot(Q)$
(Lemma~\ref{lem:rot}, applicable by Lemma~\ref{lem:carriers}), $m_Q$ for its
number of crossings, and $\ell_Q$ for its number of left turns.
\end{definition}

\begin{lemma}[Carriers on the weakly generic locus]\label{lem:weak-carriers}
Let $P\in\mathcal W_n$, $n\geq3$, be weakly generic
(Definition~\ref{def:weak}). Form its Gauss word and interlacement graph
from its actual transverse crossings and traversal visits, and apply
Definitions~\ref{def:decomposition}, \ref{def:smoothing}
and~\ref{def:uniform} with Convention~\ref{conv:selected-visits}. These
constructions are well defined, and every conclusion (i)--(iv) of
Lemma~\ref{lem:carriers} holds for every independent set $S$ of crossings
of $P$. In particular every carrier lies in the regular locus, so $r_Q$,
$m_Q$ and the uniform/mixed classification of Definition~\ref{def:uniform}
are defined.
\status{new (round~3; adopted from the Codex proposal W-P1, P-25-13, lane sm2\_weak\_c\_v1, re-read by the author; extends Lemma~\ref{lem:carriers} to $\mathcal W_n$ under the explicit weak guards; F-25-90)}
\end{lemma}
\begin{proof}
Fix a labelled representative of $P$. By Lemma~\ref{lem:silentlocus} the
visible signature is locally constant on $\mathcal W_n$ and generic tuples
are dense there; choose a generic tuple $P'$ in a neighbourhood of $P$ on
which the signature is constant. Its crossings have the same pairs of
parent edges and the same order along each edge as those of $P$, so
inserting the original vertex labels gives the same finite marked cyclic
list, the same successor swaps, the same independent sets, the same cycle
assignments and the same incoming-visit ownership for $P$ and for $P'$.
Lemma~\ref{lem:carriers} for $P'$ therefore gives (i), (iv) and the two
interlacement and ownership assertions of (iii) for $P$, these being
statements about the finite marked list only. The geometric assertions are
verified directly at $P$.

All vertices of $P$ are distinct: coincident adjacent vertices would give a
zero edge, and if two nonadjacent vertices coincided, either edge incident
to one of them would be a nonincident closed edge containing the other,
contrary to Definition~\ref{def:weak}. All edges and all turn determinants
are nonzero by that definition. Two remote edges meet at most once and
transversally, and a meeting point is interior to both edges, since a vertex
on a nonincident closed edge is excluded. Distinct crossings have distinct
points, since two crossings at one point would put at least three edge
interiors there, contrary to (G2).

Consequently the marked points on each original edge have strictly
increasing parameters with positive separations. At a selected visit the
new successor is the old outgoing subsegment of its twin visit, of positive
length, so no zero segment is inserted between the coincident twin points;
omitting unselected marks concatenates positive subsegments of one directed
edge. Every carrier edge is therefore nonzero and is an original directed
edge vector times a positive scalar. At an original corner the two carrier
vectors are positive multiples of the parent's consecutive vectors, so the
turn keeps the parent's nonzero sign; at a selected crossing of $E_i,E_j$
the two smoothing corners have determinants with the signs of
$\det(d_i,d_j)$ and $\det(d_j,d_i)$, nonzero and opposite by
transversality. No corner is flat or antiparallel. Every cycle contains a
corner, since a cycle with only unselected visits would follow the
unchanged parent successor and reach an original vertex; one nonzero edge
cannot close, and two closing nonzero edges would be opposite and give an
antiparallel corner. Hence every carrier has at least three corners. This
proves (ii), including membership in the regular locus.

For the geometric part of (iii): distinct traced subsegments of one parent
edge have disjoint interiors; adjacent parent edges are not collinear,
their turn determinant being nonzero, and meet only at their common vertex;
every other meeting of two carrier segments is an original crossing. At a
selected crossing the two visits belong to different carriers by (i), so a
carrier passes through that site once. At an unselected crossing a carrier
passes twice exactly when it owns both visits; the two germs are the
parent's transverse germs, and the site is not a corner because it differs
from every original vertex and every selected site. No point has three
preimages on a carrier: that would give three parent traversal preimages,
whereas (G2) and the vertex exclusions allow at most two. These are the
remaining assertions of (iii).

Cyclic relabelling renames the marked list, the independent sets and the
cycles without changing the traced oriented curves or the ownership of
visits, so every construction and conclusion descends to the polygon.
\end{proof}

\subsection{The positive lift and the corner coefficient}

\begin{definition}[positive crossing; positive lift]\label{def:positive-lift}
An oriented link diagram in the plane is a finite collection of closed
oriented curves with finitely many transverse double points, no triple points,
and at each double point a choice of the \emph{over} strand. A double point
with over-strand direction $u_{\rm o}$ and under-strand direction $u_{\rm u}$
is \emph{positive} if $\det(u_{\rm o},u_{\rm u})>0$ and \emph{negative}
otherwise. The \emph{positive lift} of a subpolygon $Q$ is the oriented knot
diagram whose curve is $Q$, whose double points are the crossings of $Q$, and
in which at every double point the over strand is chosen so that the crossing
is positive. Its writhe (the sum of crossing signs) is $m_Q$.
\end{definition}

\subsection{An internal named-record bridge}

We use the plane orientation and crossing-sign convention of
Definition~\ref{def:positive-lift}. A diagram here is a finite polygonal
immersion, or a regular smooth immersion with finitely many transverse
double points and possible finitely many corners away from crossings.
The corners used below admit clean rounding discs; in particular all SM
subpolygon diagrams and their regular clean modifications are included.
The parametrizing circle is oriented. To retain the outer region in the
construction, compactify the oriented plane to an oriented sphere.
All topological extensions needed for the bridge are proved below.
The named record defined next (crossing occurrences, successor, pairing,
over/under bits and signs) is what this section and
Section~\ref{sec:knotlaws} call the \emph{crossing record} of a diagram.

\begin{definition}\label{def:gauss-record}
Let $\gamma:C\to S^2$ be an actual generic immersion of an oriented
circle in an oriented sphere, with over/under choices at its transverse
double points. Its finite crossing-occurrence set $M$ has forward successor
$s$, pairing involution $\tau$ without fixed points, an over/under bit at
each occurrence, and crossing signs
\begin{equation}\label{eq:gauss-cross-sign}
 \sigma(c)=\sgn\det(u_{c,O},u_{c,U}).
\end{equation}
A named record isomorphism is a bijection $\Phi:M\to M'$ preserving
successor, pairing, over/under bits and these signs. The parametrizing
circles are oriented; the finite bijection must preserve their cyclic
orders. It does not prescribe a map at every unmarked parameter and
does not permit traversal reversal. After a finite subdivision we choose
an orientation-preserving piecewise-linear circle map
$\overline\Phi:C\to C'$ extending $\Phi$: on each interval between
successive marked points use the positive affine map in oriented interval
coordinates. If $M$ is empty, choose any positive circle parametrization.
\end{definition}

\begin{lemma}[Finite planar models]\label{lem:gauss-pl-model}
The finite polygonal diagrams and regular smooth modifications used here
admit finite piecewise-linear shadow models by orientation-preserving
ambient isotopies of the page. The crossing names, the four-ray orders,
the traversal direction and the over/under designations are retained.
No preassigned homeomorphism at unmarked circle parameters is required.
\status{proved (refereed: bench A, 2026-09-05T15:51:04Z; transcribed from RC fd:pl-model (C016))}
\end{lemma}
\begin{proof}
Choose pairwise disjoint small crossing discs, each missing every remote
arc. At a smooth transverse double point, in local parameters the map
$(s,t)\mapsto\gamma_1(s)+\gamma_2(t)-p$ has invertible derivative.
After its derivative is used as the linear coordinate frame, its inverse
straightens the two branches and has derivative arbitrarily close to the
identity on a sufficiently small disc. Multiply its displacement by a
cutoff which is one on a smaller disc and zero near the outer boundary.
Choose the inner and outer radii with a fixed ratio; shrinking them makes
the displacement derivative, including the cutoff term, have norm less
than one. The interpolated maps $x\mapsto x+t v(x)$ are injective,
since $|v(x)-v(y)|<|x-y|$, and surjective, since
$x=y-tv(x)$ is solved uniquely by contraction. They are homeomorphisms
fixed outside the disc and give the required local isotopy.

After this operation the remaining edge pieces are compact embedded
regular arcs with disjoint thin strips. A strip can be constructed by
$(s,u)\mapsto\gamma(s)+u n(s)$: for a smooth arclength arc its
Jacobian is $1-u\kappa(s)$, nonzero for uniformly small $u$.
Local injectivity follows in a tangent chart; pairs of parameters outside
those local intervals have positive distance by compactness, giving
global injectivity after one further shrink. The strips can be chosen
disjoint outside the fixed crossing discs. Polygonal interpolation at
a sufficiently fine finite partition has tangent directions uniformly
close to the arc's tangent. Its strip projection is therefore strictly
increasing, so it is one continuous normal graph $u=f(s)$, with $f=0$
at fixed ends and arbitrarily small norm. In strip coordinates use
\begin{equation}\label{eq:gauss-strip-isotopy}
 (s,u)\longmapsto(s,u+t\chi(u)f(s)),\qquad 0\leq t\leq1.
\end{equation}
Take $\chi=1$ near zero, $\chi=0$ near the strip boundary, and
$\|\chi'\|\|f\|<1$. Each fibre map is strictly increasing,
onto and fixed at the strip ends; its inverse is continuous. This is an
ambient page isotopy. Finite polygonal corner pieces already need no
straightening; a smooth replacement in a clean corner disc is treated
by the same graph construction. Subdivide the fixed straight crossing
germs and all resulting polygon edges. This gives a finite PL graph
model. All maps preserve page orientation and the named finite data.
Their products with the height coordinate give ambient isotopies of the
diagram lifts. There is no application of a smooth theorem to the image
of an arbitrary circle homeomorphism.
\end{proof}

The same finite-disc and strip proof applies to finitely many components:
each mixed crossing still has two transverse local branches, and the
remaining compact arc collection is finite. All components are retained.


\begin{lemma}[Constructive polygonal two-disc theorem]
\label{lem:gauss-two-discs}
A simple polygonal circle in the oriented sphere has exactly two
complementary regions, and both closures are PL discs. Finite prescribed
positive PL boundary maps between such discs extend to positive PL disc
maps. A continuous positive boundary map also extends to a topological
disc map. The statement includes the exterior region.
\status{proved (refereed: bench A, 2026-09-05T15:51:04Z; transcribed from RC fd:pl-discs (C016))}
\end{lemma}
\begin{proof}
Here is a finite triangulation and disc construction, so no Jordan or
Schoenflies conclusion is being assumed. Remove one point off the curve
and put the polygon inside a large rectangle. Extend its finitely many
edge lines, and also lines through its vertices in a generic auxiliary
direction, across that rectangle. The arrangement has convex cells;
subdivide boundary segments at all intersections and triangulate each
convex cell by a fan from an interior point. Collinear redundant vertices
are simply retained as subdivisions, never used to claim a strict ear.
The circle is now a simple edge cycle $C$ in a finite triangulation of
the rectangle. Cap the rectangular boundary by a second triangulated disc.
This cap represents the actual one-point-compactified exterior: from
the rectangle's center, write its boundary as $r=R(u)$ and map an
exterior point $ru$ to $(R(u)/r)u$, with infinity sent to zero.
This has a continuous radial inverse and identifies the exterior with a
disc. The two-disc model is a triangulated sphere by radial identification
with the upper and lower hemispheres. The original exterior is part of this same
triangulation, not just a collar of $C$.

Work with chains over $\mathbb F_2$. If the sphere triangulation has
$V,E,F$ cells, then $V-E+F=2$. This follows first for the two-disc
rectangle model and is unchanged by subdividing an edge, adding a vertex
inside a face, or drawing a diagonal; these are precisely the finite
arrangement and fan subdivisions just used. The vertex-edge graph and
the face-adjacency graph are connected. The incidence map on edges has
rank $V-1$, since a spanning tree gives $V-1$ independent boundary
columns and every edge boundary has even coordinate sum. Hence its
cycle space has dimension $E-V+1=F-1$. The face-boundary map has
kernel exactly the two constant face assignments: across every edge a
boundary-zero face assignment has equal coefficients on its two adjacent
faces, and the face-adjacency graph is connected. Its rank is therefore
$F-1$. Face boundaries are cycles, so every cycle, in particular $C$,
is the boundary of a face assignment, unique up to complement.

Color faces by that assignment. Across an edge of $C$ the colors differ;
across any other edge they agree. At a vertex of $C$ its two cycle
edges divide the cyclic triangle fan into exactly two consecutive fans,
one of each color. At every other vertex all incident faces have one
color. Thus each color closure is a compact triangulated surface whose
boundary is $C$, with no pinch point. Following $C$ through those local
fans shows that all boundary-adjacent faces of a given color lie in one
connected face component. Any further face component would have no
boundary edge and hence no adjacency to its complement, contradicting
connectedness of the full dual graph. Thus there are exactly two
connected color regions, $R_0,R_1$, each with its single boundary $C$.

We next prove, rather than assume, that each is a disc. For any connected
triangulated surface $R$ with nonempty boundary, take a spanning tree
of its triangle-adjacency graph, rooted at a triangle having a boundary
edge. Remove that triangle across that boundary edge, and then remove
each triangle after its parent, across their common edge. At its turn
the chosen edge is free: the parent triangle is gone and no third
triangle is incident to that edge. A triangle and its free edge collapse
onto its other two edges. This preserves connectedness and Euler
characteristic and leaves, after finitely many steps, a connected graph
$K$. Consequently $\chi(R)=\chi(K)=1-\beta_1(K)\leq1$.
Since the two regions meet in $C$ and $\chi(C)=0$, cell counting gives
$\chi(R_0)+\chi(R_1)=2$. Both inequalities are therefore equalities,
so both residual graphs are trees.

For completeness these collapses also construct disc parametrizations,
not only a homotopy equivalence. Maintain a small closed regular
neighbourhood of the current collapsed complex, with vertex polygons
and edge rectangles. A free-triangle collapse changes its neighbourhood
by removing a boundary notch, which gives a boundary-compatible PL
homeomorphism. In the standard triangle
$x,y\geq0$, $x+y\leq1$, whose free edge is $x+y=1$, use the
neighbourhood of the other two edges
\begin{equation}\label{eq:gauss-triangle-notch}
 \{x,y\geq0, x+y\leq1,\ \min(x,y)\leq\varepsilon\},
 \qquad 0<\varepsilon<1/2.
\end{equation}
Its new boundary path runs from $(1,0)$ through
$(1-\varepsilon,\varepsilon),(\varepsilon,\varepsilon),
(\varepsilon,1-\varepsilon)$ to $(0,1)$.
Subdivide the old free edge at the intersections of the rays from
$(0,0)$ through those three corners. Map the resulting triangle fan
affinely to the corresponding fan of the notched region, fixing the two
retained edges pointwise. Every triangle map has positive determinant,
the fans meet on their common radial edges, and the images cover the
notched region. Thus the map is a PL homeomorphism, fixed where that
triangle joins the rest. In thin edge rectangles use the same product
map, with identity on the outer attaching boundary. This proves by
induction the stated regular-neighbourhood invariant for every free
collapse. Choose widths successively smaller than the positive distances
between disjoint cells; finiteness makes all choices compatible.

A regular neighbourhood of a tree is a disc by an explicit leaf
induction. One vertex gives a polygonal disc. Adding a leaf attaches
one edge rectangle and one leaf disc along a single boundary interval.
Triangulate that rectangle and disc into a fan along the interval; the
same positive fan maps identify their union with a longer boundary
collar. Thus each step is a disc and records its boundary parametrization.
Reverse the finite free-collapse maps to parametrize $R_0$ and $R_1$.
This completes both the interior and exterior constructions.

Finally identify two resulting discs with convex polygonal unit discs.
Subdivide their boundaries until a specified positive PL boundary map
is affine on each segment, choose one interior vertex in each, and map
the corresponding boundary fans affinely, interior vertex to interior
vertex. Their cyclic order makes all determinants positive. For a
continuous positive boundary map instead use radial coordinates and
$(r,u)\mapsto(r,b(u))$; the inverse uses $b^{-1}$ and continuity at
the center follows from the radial bound. This proves both extension
claims with the prescribed boundary, not just an unlabelled disc type.
\end{proof}

\begin{lemma}[PL sphere isotopy]
\label{lem:gauss-sphere-isotopy}
Every positive PL self-homeomorphism of the sphere is isotopic to the
identity through positive topological homeomorphisms. Only a relocation
of two sufficiently small triangles is needed in the construction.
\status{proved (refereed: bench A, 2026-09-05T15:51:04Z; transcribed from RC fd:sphere-isotopy (C016))}
\end{lemma}
\begin{proof}
We first give the required local motion. In a planar chart let $v(x)$
be supported in a ball and Lipschitz with constant less than one.
Then $x\mapsto x+t v(x)$ is a homeomorphism for $0\leq t\leq1$:
injectivity follows from the strict Lipschitz estimate and surjectivity
from contraction applied to $x=y-tv(x)$. Its inverse depends
continuously on $(y,t)$ by the same estimate. For a sufficiently small
increment $x\mapsto Mx+b$ of a positive affine motion, choose a cutoff
equal to one near the moving small triangle and zero near the chart
boundary, and set $v(x)=\psi(x)((M-I)x+b)$. A sufficiently fine time
subdivision makes its Lipschitz constant less than one. Thus any compact
path of small positive affine triangles, with a positive clearance from
the outer boundary, has an ambient extension fixed on that boundary.

Here a path between two interior triangles can be given explicitly:
contract each about its centroid to an arbitrarily small copy, change
the first shape to the second by a positive linear path, translate its
centroid along a polygonal path in the chart, and expand at the end.
The compact centroid path has positive boundary clearance, so a common
small size suffices. A positive linear path exists without a matrix
connectivity assumption: rotate the first column of a matrix to the
positive horizontal axis; the remaining upper-triangular matrix has
positive diagonal entries, and its straight interpolation to the
identity keeps both diagonal entries positive. This supplies the
shape path. All the local extensions just constructed are explicitly
isotopic to the identity and never require an extension on a region
with three boundary components.

Now let $h$ be a positive PL sphere map. If it is not the identity,
choose $p$ in the interior of a PL cell with $h(p)\ne p$.
Choose a PL disc chart containing $p,h(p)$ in its interior: remove a
small triangle about a third point and use Lemma~\ref{lem:gauss-two-discs}
on its complementary disc. Move $p$ slightly, if necessary, off the
finitely many subdivision edges and their $h$-preimages; the inequality
$h(p)\ne p$ persists. Refine so the chart and $h$ are affine
on a sufficiently small triangle $B$ about $p$ and on its image.
Take $B$ so small that $B$ and $h(B)$ are disjoint. Their two chart
images are positive triangles. The preceding local construction gives
a positive map $k$, supported in that chart interior and isotopic to
the identity, with $k(h(B))=B$. Thus $g=kh$ preserves $B$ setwise.
No claim about relocating arbitrary prescribed discs is needed.

Its boundary map $b$ is positive. In a positive circle coordinate, let
$\beta:\RR\to\RR$ be an increasing lift of $b^{-1}$ with
$\beta(x+1)=\beta(x)+1$. The maps
$\beta_t(x)=(1-t)x+t\beta(x)$ are strictly increasing with that
same periodicity and descend to a circle isotopy. The two discs bounded
by $\partial B$ have parametrizations from
Lemma~\ref{lem:gauss-two-discs}; extend $\beta_t$ radially over both and
glue their common boundary maps. Postcomposing $g$ by its last map
fixes $\partial B$ pointwise.

A boundary-fixed unit-disc homeomorphism $a$ has the explicit isotopy
\begin{equation}\label{eq:gauss-alexander}
 A_t(x)=\begin{cases}t\,a(x/t),&|x|\leq t,\\
 x,&t\leq|x|\leq1,\end{cases}\qquad A_0(x)=x.
\end{equation}
At $|x|=t$ the formulas agree. At $t=0$ both the moved point and its
image have norm at most $t$, proving continuity. Replacing $a$ by
$a^{-1}$ gives the inverse, also continuous. Apply this separately on
the two boundary-fixed discs and glue. Concatenating, and reversing
where necessary, the relocation, boundary-correction and two Alexander
isotopies gives an isotopy from the identity to $h$. Every constituent
preserves orientation.
\end{proof}

\begin{lemma}[Compatible heights have a fibrewise ambient isotopy]
\label{lem:gauss-height-isotopy}
For one actual generic finite diagram, two continuous height assignments
with the same strict over/under order at every double point represent
ambient-isotopic oriented links. This holds for any finite number of
components and requires no differentiability of the height functions.
\status{proved (refereed: bench A, 2026-09-05T15:51:04Z; transcribed from RC fd:height-isotopy (C016))}
\end{lemma}
\begin{proof}
Choose a height bound $H$ larger than the absolute values of both
assignments. In a small crossing chart each of the two branch heights
extends continuously to the chart by holding it constant in the transverse
coordinate of that branch. Shrink the chart so both source and target
extensions retain their strict lower/upper ordering. For each page point
$x$ define an increasing piecewise-linear map $f_x$ of the height axis
which sends, in order, the four source levels
$-H,h_{0,U}(x),h_{0,O}(x),H$ to
$-H,h_{1,U}(x),h_{1,O}(x),H$, and is identity outside $[-H,H]$.
All slopes are positive and depend continuously on $x$. On a strip
meeting only one branch use the same construction with just that one
intermediate source and target height. Take these strips outside smaller
crossing discs so that no strip meets an unaccounted second branch.
These finitely many local maps cover the shadow.

Choose nonnegative continuous weights subordinate to these chart
interiors and an identity map on their complement, whose sum is one
and whose identity weight is zero on the shadow. Explicitly use
distances to the complements of slightly smaller chart interiors,
together with the distance to the closed shadow, and divide by their
positive sum. The smaller charts still cover the shadow, so the sum
never vanishes; extending weighted displacements by zero across a
chart boundary is continuous since all displacements are bounded by
$2H$. Set $f(x,z)$ to the weighted average of the local $f_x(z)$
and the identity value $z$. A convex combination of strictly increasing
functions is strictly increasing, and all functions equal $z$ for
$|z|\geq H$. On a point of the shadow every participating local
map sends each actual source height there to its actual target height,
so their average does too.

Thus
\begin{equation}\label{eq:gauss-height-map}
 (x,z)\longmapsto
       \bigl(x,(1-t)z+t f(x,z)\bigr),\qquad 0\leq t\leq1,
\end{equation}
is a homeomorphism on every fibre and fixes the product boundary.
Continuity of its inverse follows directly: if $(x_j,t_j,w_j)$
converges, its inverse heights are bounded by $\max(H,|w_j|)$;
every convergent subsequence has the unique inverse limit by continuity
and strict monotonicity, hence the full sequence converges. The maps
extend by the identity to the ambient three-sphere and form an ambient
isotopy. They transport the original parameter orientation. No tubular
smooth-extension hypothesis is used.
\end{proof}

\begin{theorem}[Named-record bridge]\label{lit:gauss}
Let $D,D'$ be oriented one-component generic diagrams in the finite
polygonal or regular smooth class specified above, in the oriented plane.
Suppose a bijection of crossing visits preserves their oriented
cyclic order, crossing pairing, over/under designation and crossing sign
$\operatorname{sgn}\det(u_{\rm o},u_{\rm u})$.
Then the two diagrams present the same oriented knot: their compatible
lifts in $S^3$ are joined by an ambient isotopy preserving the traversal
orientation. The assertion includes diagrams with no crossings.
No preservation of the unbounded face or of the point at infinity is
assumed or asserted.
\status{proved (refereed: bench A, 2026-09-05T15:51:04Z; transcribed from RC fd:gauss, printed in place of registry L-2 under ruling R-25-14, flagged for source cross-check (C016))}
\end{theorem}
\begin{proof}
Add one point at infinity to each oriented plane. The resulting pages are
oriented spheres, and the diagrams miss the added point. View compatible
lifts in a thin product neighbourhood of the spherical page in $S^3$.
The hypothesis is exactly a named record isomorphism of
Definition~\ref{def:gauss-record}: its sign convention is that of
Definition~\ref{def:positive-lift}. Apply Lemma~\ref{lem:gauss-pl-model} to each diagram first, transporting its
finite record. Choose circle coordinates pulled back from the subdivided
PL shadow arcs; no original smooth parametrization at unmarked points
is prescribed. Subdivide these arcs and choose the positive PL extension
$\overline\Phi$ of the finite occurrence bijection as in
Definition~\ref{def:gauss-record}. The construction below concerns this chosen
extension, not an arbitrary prescribed map on every circle parameter.
For paired visits $v,w$, the determinant sign of the ordered tangents
$(u_v,u_w)$ is the recorded crossing sign if $v$ is the over visit, and its
negative if $v$ is the under visit. Thus the named bijection preserves this
determinant sign for either ordering. No reversal of traversal, page
orientation or crossing signs is allowed.

Write $v^+,v^-$ for the outgoing and incoming rays at a crossing, both
pointing away from it. For $w=\tau(v)$ the positive cyclic order is
\begin{equation}\label{eq:gauss-ray-order}
 \begin{cases}
 (v^+,w^+,v^-,w^-),&\det(u_v,u_w)>0,\\
 (v^+,w^-,v^-,w^+),&\det(u_v,u_w)<0.
 \end{cases}
\end{equation}
Indeed in the first case $u_w$ lies between $u_v$ and $-u_v$
counterclockwise, and in the second $-u_w$ does; the other two rays are
antipodes. The named map preserves these orders. It induces a graph map
\begin{equation}\label{eq:gauss-shadow-map}
 f_\Phi(\gamma(x))=\gamma'(\overline\Phi(x)).
\end{equation}
This is well defined because equal images of distinct parameters are exactly
paired crossing preimages. The inverse argument proves it a homeomorphism
of the directed shadow graphs, preserving the named rays and their rotation
systems.

Suppose first there is at least one crossing. The shadow is a connected
finite four-valent graph. Construct its closed regular neighbourhood $N$
from vertex discs and edge bands. For each boundary circle, the connected
interior of $N$ lies on one side; the other disc supplied by
Lemma~\ref{lem:gauss-two-discs} is disjoint from $N$.
These complementary discs cannot be nested, since an inner boundary belongs
to $N$ and would then lie in a disc disjoint from $N$. To see that their
interiors cover the complement of $N$, join any point outside $N$ by a path
to its connected interior and take the first boundary hit; until that hit
the path remains on the opposite side of that boundary circle, so the
initial point lies in its complementary disc. Adding the collars proves that the shadow fills the
sphere with disc faces. Connectedness of an abstract ribbon alone would not
prove this: if it has $q$ crossing vertices and $b$ boundary cycles, capping
gives genus $(q+2-b)/2$, which need not be zero without the actual sphere
embedding.

Attach abstract oriented discs and untwisted bands in the orders
\eqref{eq:gauss-ray-order}. Their positive boundary successor is determined by the recorded ray
orders and the directed edges joining successive visits. Map these labelled cells to the
actual vertex discs and edge bands. This gives an orientation-preserving
ribbon identification relative to the entire named spine, including each
boundary cycle. Applying it in both diagrams extends \eqref{eq:gauss-shadow-map}
to an oriented ribbon map. Each boundary has its corresponding complementary
disc. Subdivide its boundary map and extend by the positive triangle-fan
construction of Lemma~\ref{lem:gauss-two-discs}, then glue. The resulting
named PL sphere homeomorphism $h$ preserves ambient
orientation, the oriented traversal and the over/under designation
separately. These three kinds of orientation data are preserved separately by the
construction.

Identify the two oriented spheres. Lemma~\ref{lem:gauss-sphere-isotopy}
constructs a positive sphere isotopy $h_t$ from the identity to this
particular PL map $h$.

Lift each decorated diagram into a thin product neighbourhood of the sphere,
with the over height larger than the under height at each crossing.
$h\times\mathrm{id}$ carries the source lift to a compatible target lift.
Regard $S^3$ as the suspension of $S^2$; concretely
$[x,s]\mapsto(\sqrt{1-s^2}\,x,s)\in S^3\subset\RR^4$.
Then
\begin{equation}
 [x,s]\longmapsto[h_t(x),s]
\end{equation}
is an ambient isotopy of $S^3$, including the two suspension points, and
agrees with the product map near the equator. The height function of the
transported lift is continuous and has the target over/under order.
Lemma~\ref{lem:gauss-height-isotopy} gives an ambient isotopy to the chosen
target lift directly at that continuous regularity. Its parameter
orientation is unchanged.

If there are no crossings, both shadows are embedded oriented circles.
The map $\gamma'\overline\Phi\gamma^{-1}$ identifies their boundaries.
Name left and right discs by ambient and traversal orientation, using
Lemma~\ref{lem:gauss-two-discs}; map left to left and right to right with the
positive PL boundary extensions. This gives a PL oriented sphere map,
to which the same sphere isotopy and fibrewise lift argument applies.
No empty record is sent through the crossing-labelled ribbon construction.

The sphere map may send the original point at infinity into a different
complementary face. This is harmless: the isotopy constructed above lives
in $S^3$, not in the plane with infinity fixed. Therefore this proof gives
oriented knot equivalence without claiming a plane isotopy that preserves
the outer face.
\end{proof}


\begin{definition}[carriers of the flat configurations]\label{def:flat-carriers}
Under the hypotheses of Lemma~\ref{lem:flat-sides}, identify the crossing
visits of the two sides, of the centre $P(0)$ and of the deletion
$P(0)\setminus j$ through the common cyclic Gauss word of that lemma's
clauses (ii) and (iii), and fix an independent set $S$ in their common
interlacement graph. In each of these four configurations mark the
traversal circle at every original vertex and every crossing visit,
exchange the two outgoing successors at the two visits of every selected
crossing --- the reconnection of Definition~\ref{def:smoothing} with
Convention~\ref{conv:selected-visits} --- and trace the resulting oriented
closed cycles by the inherited straight subsegments. These oriented closed
polygonal cycles are the \emph{carriers} of $S$ in that configuration. On
the two generic sides they are the carriers of
Definition~\ref{def:smoothing}; at the centre and on the deletion the same
words define them, no genericity being assumed.
\end{definition}

\begin{corollary}[flat carrier data]\label{cor:flat-carriers}
Under the hypotheses of Lemma~\ref{lem:flat-sides}, fix an independent
set $S$ in the common interlacement graph and form the carriers of $S$ on
both sides, at the centre and on the deletion
(Definition~\ref{def:flat-carriers}). Then:
\begin{enumerate}
\item[(i)] The carriers correspond under their named traversal arcs.
Exactly one contains $\mu_j$. The central copy of that carrier differs
from its deletion copy only by the positive-flat subdivision at $\mu_j$;
every other central carrier is unchanged by deletion. Every carrier
has nonzero segments and no antiparallel corner; except for that one
central zero turn, all corner turns are nonzero.
\item[(ii)] Corresponding carriers have the same retained self-crossing
visits, pairing, signs and positive over/under bits. They have the same
signed rotation and hence the same absolute rotation. Every corresponding
corner away from $\mu_j$ has the same turn sign on both sides and in
the deletion. The extra corner at $\mu_j$ is right on the right side
and left on the left side.
\item[(iii)] Define a carrier's selector to be $1$ if all its turns are
right, $(-1)^c$ if all its $c$ turns are left, and $0$ if its turns are
mixed. For the carrier through $\mu_j$, let $W_{\rm right}$,
$W_{\rm left}$, and $W_{\rm del}$ be these selectors on the two sides
and on the deletion. Then
\begin{equation}\label{flatpr:selector-identity}
W_{\rm right}-W_{\rm left}=W_{\rm del}.
\end{equation}
All other corresponding carrier selectors agree.
\end{enumerate}
These assertions concern geometric and combinatorial carrier data.
They make no assignment of a state-sum value at the flat centre.
\status{proved (refereed: bench A, 2026-09-05T23:45:25Z; transcribed from CV lem:flatdata, carrier data (P-25-3, C007); round~3: the carriers are those of Definition~\ref{def:flat-carriers}, F-25-102; round~4: clause~(ii)'s rotation sentence covers the centre copy as well, which CV lem:flatdata(iv) omits --- a broadening of this document, refereed sound by bench~A (F-25-141))}
\end{corollary}

\begin{proof}
The common cyclic word identifies the visits and the successor
permutation before smoothing. Exchanging the same selected successors
in that permutation gives the same cycles afterwards. This is an
explicit correspondence and needs no ambient knot-equivalence theorem.
Noninterlacement ensures that each selected pair separates two carriers:
cut at one selected pair; each other selected pair lies entirely on one
of the two resulting arcs, since it does not interlace that pair.
Apply the same argument recursively on those arcs. It gives $|S|+1$
cycles, with the two reconnections at any selected crossing on distinct
cycles. Thus no carrier visits a selected smoothing point twice.

An original vertex appears once on the parent traversal; selected
crossing points are distinct and avoid all vertices. The subsegments
between successive vertices or crossing visits therefore have positive
length. The occurrence $\mu_j$ lies on exactly one cycle. Its incoming
and outgoing directions at the centre are positive multiples by
\eqref{flatpr:fusion}; removing it merely fuses those segments.
At another original corner the determinant is a positive multiple of
the nonzero original turn determinant. At a selected smoothing corner
it is a positive multiple of either $\det(d_e,d_f)$ or
$\det(d_f,d_e)$ at that transverse crossing. These exhaust the corners.
They prove the stated nonzero turns and absence of antiparallel corners,
including in the deletion. A deletion carrier has at least three
corners: one nonzero segment cannot close, and a closed two-segment
polygon has antiparallel directions. No such deletion carrier occurs.

An unselected crossing is a self-crossing of a carrier exactly when
both its visits belong to that cycle. The common successor permutation
preserves this condition and the two visits' inherited positive
over/under bits. Tracing inherited straight subsegments introduces no
other self-intersection: every intersection was already a parent
crossing, and a selected crossing lies on two different carriers, as
proved above. The nonzero determinants just listed have constant
signs on a smaller interval. Replacing either old incident direction
by the fused direction multiplies the relevant determinant by a
positive scalar. Hence all surviving corner signs and all retained
crossing data agree among the configurations in (ii).

For completeness, the rotation comparison uses limits of every affected
turn, not equality of finite-parameter angles. For a closed polygonal
carrier with nonzero segments and no antiparallel consecutive directions,
let $\vartheta_k\in(-\pi,\pi)$ be its principal turns. The successive
unit directions satisfy $q_{k+1}=e^{\mathrm i\vartheta_k}q_k$.
Multiplying around the closed cycle gives
\begin{equation}\label{flatpr:rotation-integer}
e^{\mathrm i\sum_k\vartheta_k}=1,\qquad
\frac1{2\pi}\sum_k\vartheta_k\in\mathbb Z.
\end{equation}
Principal turns are continuous when the two incident directions are
nonzero and not antiparallel. On the carrier through $\mu_j$ its
distinguished turn tends to zero, and each other principal turn tends
to the corresponding deletion turn. On every other carrier all turns
have their common deletion limits. Thus each side's signed rotation
tends to the deletion's integer rotation. An integer with distance
less than $1/2$ from that integer must equal it. There are finitely
many supports and carriers, so a common sufficiently small interval
works for all of them. This proves the signed and absolute rotation
equalities in (ii).

It remains to evaluate the selector. Let $c$ be the number of surviving
corners in the distinguished carrier. Their signs agree in all three
configurations. Its extra turn is right or left according to the named
side. The exhaustive table is
\begin{equation}\label{flatpr:selector-table}
\begin{array}{c|ccc}
\text{surviving turns}&W_{\rm right}&W_{\rm left}&W_{\rm del}\\\hline
\text{mixed}&0&0&0\\
\text{all right}&1&0&1\\
\text{all left}&0&(-1)^{c+1}&(-1)^c
\end{array}
\end{equation}
In the last row $0-(-1)^{c+1}=(-1)^c$; the other two rows have the
same first-minus-second identity immediately. All other carriers
retain all their corner signs and counts. This proves (iii).
\end{proof}


\begin{literature}[HOMFLY--PT polynomial]\label{lit:homfly}
There is a map $D\mapsto H_D(a,z)\in\ZZ[a^{\pm1},z^{\pm1}]$ on oriented
link diagrams that is invariant under the three Reidemeister moves and planar
isotopy, takes the value $1$ on the crossing-free circle, and satisfies
$aH_{D_+}-a^{-1}H_{D_-}=zH_{D_0}$ for every skein triple. Its value depends
only on the oriented link presented by $D$.
\status{lit: Registry~\ref{reg:homfly}; existence, the skein, the unknot value and Reidemeister invariance: Lickorish--Millett \cite[Theorem, pp.~112--113; proof pp.~113--120; isotopy-class descent p.~120]{LickorishMillett}, the coefficient ring $\ZZ[a^{\pm1},z^{\pm1}]$ under $l=ia$, $m=-iz$ from LM property~(1), p.~110, proved as their Proposition~22 \cite[statement p.~133, proof pp.~133--134]{LickorishMillett}, whose parity clause holds at every component count (seat L-1 EXIT; F-25-169), the descent through Reidemeister's theorem \cite{Reidemeister,Reidemeister1927}; the former uniqueness clause over $\ZZ[a^{\pm1},z^{\pm1}]$ was stated in no source, was consumed by no proof of this document from round~3 on, every consumer taking uniqueness from Lemma \texttt{lp:coefficient-transport}, and was struck from the statement in round~4 (R-25-33, F-25-131); F-25-110, F-25-54, F-25-55}
\end{literature}

\subsection{The local source construction and its algebraic consequences}

The source fact below is the literal construction used by
Lickorish--Millett. Its local conclusions suffice for the algebraic
proofs in this subsection. The broader source claim in
Literature input~\ref{lit:homfly} remains separately declared; no
arbitrary ambient-isotopy theorem is inferred from the local fact.
Diagrams are actual nonempty finite oriented diagrams, including all
crossing-free components. Algebraic empty products are not empty links.

\begin{literature}[Exact local LM construction projection]\label{lp:lm}\label{src:lm}
For every such diagram $D$, let $F_D(l,m)$ be the function constructed in
Section~1 of Lickorish--Millett~\cite{LickorishMillett}, not an arbitrary function with similar
values. The source constructs it in
$\mathbb Z[l^{\pm1},m^{\pm1}]$, independently of the auxiliary component
order, nonsingular basepoints, and order of required crossing switches.
It is defined modulo positive page isotopy, satisfies
\begin{equation}\label{lp:source-skein}
 l F_{D_+}+l^{-1} F_{D_-}+m F_{D_0}=0,
 \qquad \mu=-(l+l^{-1})m^{-1},
\end{equation}
and has the initialization $F_D=\mu^{c-1}$ for every based ordered
$c$-component UNDER-first diagram. UNDER-first means that, traversing
components in their order, from their basepoints and in their orientations,
the first encounter with each crossing is the underpass. In particular
every crossing-free $c$-component diagram has that value, regardless of
nesting or orientations. Here $c\geq1$.

The same function is invariant under each actual ordinary oriented
Reidemeister move. For a supplied finite sequence, take the maximum of the
finitely many crossing counts and apply source Proposition~4 at that level
one move at a time. No bound by the endpoint crossing counts is required.
The imported construction and its Propositions~1--6 stop before the final
p120 inference from arbitrary ambient isotopy to a finite move sequence.
\status{lit: Registry~\ref{reg:homfly}; local construction only}
\end{literature}

\begin{literature}[Uniqueness of the source polynomial over its own ring]\label{lp:lm-uniqueness}
Lickorish--Millett prove \cite[Theorem, pp.~112--113, with the uniqueness
argument on p.~120]{LickorishMillett} that the function $F_D$ of Literature
input~\ref{lp:lm} is the only map from oriented link diagrams (modulo plane
isotopy) to $T=\ZZ[l^{\pm1},m^{\pm1}]$ that depends only on the isotopy class
of the oriented link, takes the value $1$ on the unknot, and satisfies
$lF_{D_+}+l^{-1}F_{D_-}+mF_{D_0}=0$ on every skein triple; the same statement
is Theorem~15.2 of Lickorish~\cite{LickorishGTM}. Uniqueness is asserted
among all such maps, with no restriction on the support or on the
coefficients of a competitor. A map on diagrams that is invariant under
planar isotopy and the three Reidemeister moves depends only on the isotopy
class of the oriented link by Reidemeister's theorem~\cite{Reidemeister}, the
source's own input for its descent on p.~120; the uniqueness therefore
applies to every such invariant map with the source skein and the unknot
value. The ring is the source's: no statement over $\ZZ[a^{\pm1},z^{\pm1}]$
is imported here (registry L-1, ledger row F-25-54).
\status{lit: Registry~\ref{reg:homfly}; uniqueness over the source ring only, with Reidemeister's theorem}
\end{literature}

\begin{lemma}[Gaussian coefficient transport of skein uniqueness]
\label{lp:coefficient-transport}
Put $R=\ZZ[a^{\pm1},z^{\pm1}]$. Any two maps $D\mapsto Q_D\in R$ on oriented
link diagrams that are invariant under planar isotopy and the three
Reidemeister moves, take the value $1$ on the crossing-free circle, and
satisfy $aQ_{D_+}-a^{-1}Q_{D_-}=zQ_{D_0}$ on every skein triple coincide.
No restriction on the support or on the coefficients of the maps is imposed.
The lemma transports the uniqueness clause of Literature
input~\ref{lp:lm-uniqueness} to the ring $R$; it asserts no existence, no
ambient-isotopy invariance and no descent theorem.
\status{new (author, round~2; adapted from Codex S54-P1 = P-25-10, sealed lane candidate\_source\_findings\_51\_58\_v1; declared campaign step transporting Literature input~\ref{lp:lm-uniqueness} to $R$, flagged for source cross-check under rule~8; consumed in round~3 by Theorem \texttt{lp:core} (the identification $P=H$) and Theorem \texttt{cf:thm-carrierfloor} (clause~(R) and the mirror paragraph), and through clause~(R) by Proposition \texttt{prop:C-reversal}; F-25-54)}
\end{lemma}
\begin{proof}
Adjoin a formal $\mathrm i$ with $\mathrm i^2=-1$ and put
$T_G=\ZZ[\mathrm i][l^{\pm1},m^{\pm1}]$ and
$R_G=\ZZ[\mathrm i][a^{\pm1},z^{\pm1}]$. Since $1,\mathrm i$ is a free
integer basis of $\ZZ[\mathrm i]$, every element of $T_G$ is uniquely
$u+\mathrm iv$ with $u,v\in T$, and the inclusions $T\subset T_G$,
$R\subset R_G$ are injective.

\emph{Uniqueness over $T_G$.} Let $G$ be a map from oriented link diagrams
to $T_G$, invariant under planar isotopy and the three moves, with
$G_{\bigcirc}=1$ and $lG_{D_+}+l^{-1}G_{D_-}+mG_{D_0}=0$ on every skein
triple. Write $G_D=U_D+\mathrm iV_D$ with $U_D,V_D\in T$. Both coordinates
inherit the invariance, and since the three coefficients $l,l^{-1},m$ lie
in $T$, comparing coordinates gives the source skein for $U$ and for $V$
separately, with $U_{\bigcirc}=1$ and $V_{\bigcirc}=0$. The skein is linear,
so $U$ and $F+V$ are both invariant $T$-valued maps satisfying the source
skein with unknot value $1$. Literature input~\ref{lp:lm-uniqueness} gives
$U=F$ and $F+V=F$, hence $V=0$ and $G=F$. No division by $2$ is used.

\emph{Transport.} The assignments
\begin{equation}\label{lp:ring-transports}
 \phi:T_G\to R_G,\quad l\mapsto\mathrm ia,\ m\mapsto-\mathrm iz,
 \qquad
 \psi:R_G\to T_G,\quad a\mapsto-\mathrm il,\ z\mapsto\mathrm im,
\end{equation}
both fixing $\mathrm i$, send the Laurent generators to units and therefore
define ring homomorphisms. On generators,
$\psi(\phi(l))=\mathrm i(-\mathrm il)=l$,
$\psi(\phi(m))=-\mathrm i(\mathrm im)=m$,
$\phi(\psi(a))=-\mathrm i(\mathrm ia)=a$ and
$\phi(\psi(z))=\mathrm i(-\mathrm iz)=z$, so $\phi$ and $\psi$ are mutually
inverse isomorphisms. Let $Q$ be a map as in the statement, regarded as
$R_G$-valued, and put $\widetilde Q_D=\psi(Q_D)\in T_G$. It is invariant
under planar isotopy and the three moves, and
$\widetilde Q_{\bigcirc}=1$. Since $\psi(a^{-1})=(-\mathrm il)^{-1}
=\mathrm il^{-1}$, applying $\psi$ to
$aQ_{D_+}-a^{-1}Q_{D_-}-zQ_{D_0}=0$ gives
\begin{equation}\label{lp:transported-skein}
 (-\mathrm il)\widetilde Q_{D_+}-(\mathrm il^{-1})\widetilde Q_{D_-}
 -(\mathrm im)\widetilde Q_{D_0}=0,
\end{equation}
and multiplying by the unit $\mathrm i$ gives
$l\widetilde Q_{D_+}+l^{-1}\widetilde Q_{D_-}+m\widetilde Q_{D_0}=0$.
By the first step, $\widetilde Q=F$, hence $Q_D=\phi(F_D)$ for every
diagram $D$. Two maps as in the statement are therefore equal in $R_G$,
and by injectivity of $R\subset R_G$ equal in $R$.
\end{proof}

\begin{theorem}[Local campaign polynomial, algebraic part]\label{lp:core}
Set $R=\mathbb Z[a^{\pm1},z^{\pm1}]$. The same source construction has
an evaluation $P_D\in R$ with
\begin{equation}\label{lp:skein}
 aP_{D_+}-a^{-1}P_{D_-}=zP_{D_0},
 \qquad \delta=(a-a^{-1})z^{-1}.
\end{equation}
Its value on every UNDER-first $c$-component diagram is $\delta^{c-1}$.
It equals $H_D$ of Literature input~\ref{lit:homfly}. It is also the unique
function on this exact diagram domain satisfying this
skein and all these initialization values. This uniqueness statement
explicitly includes the initialization hypothesis; it is not a claim
proved from only an unknot value and unspecified ambient invariance.
For every $c$-component diagram,
\begin{equation}\label{lp:support}
 P_D\in z^{1-c}\mathbb Z[a^{\pm1},z^2],\qquad P_D\ne0.
\end{equation}
In particular $P_{\bigcirc}=1$, and knot evaluations are polynomials in
$z^2$, with no negative $z$ exponents. All local invariances in
Literature input~\ref{lp:lm} are retained.
\status{new (round~3: the identification paragraph of the proof consumes Lemma~\ref{lp:coefficient-transport}, F-25-54; the remainder of the proof is the round-1 transcription of RC lp:core, \path{d2b_local_polynomial} (C026), held by bench~A; consumes Literature inputs~\ref{lp:lm} and~\ref{lit:homfly})}
\end{theorem}
\begin{proof}
Start in $R_G=\mathbb Z[\mathrm i][a^{\pm1},z^{\pm1}]$ with
$\mathrm i^2=-1$. The elements $\mathrm ia$ and $-\mathrm iz$ are units,
so they define a Laurent-ring homomorphism and a first evaluation
\begin{equation}\label{lp:gaussian}
 \phi(l)=\mathrm ia,\qquad \phi(m)=-\mathrm iz,
 \qquad P_D^G=\phi(F_D).
\end{equation}
Inverting the first unit gives $(\mathrm ia)^{-1}=-\mathrm ia^{-1}$.
The image of the source skein is therefore
\begin{equation}\label{lp:gaussian-skein}
 \mathrm iaP_+^G-\mathrm ia^{-1}P_-^G-\mathrm izP_0^G=0.
\end{equation}
Multiplying by $-\mathrm i$ gives the campaign skein in~\eqref{lp:skein}.
For the source initialization, the numerator $l+l^{-1}$ maps to
$\mathrm i(a-a^{-1})$ and the denominator $m$ maps to $-\mathrm iz$.
Their quotient with its preceding minus sign is $(a-a^{-1})/z$.
Thus $P_D^G=\delta^{c-1}$ for every UNDER-first diagram. This calculation
does not yet assert integral coefficients for arbitrary diagrams.

Choose auxiliary basepoints and a component order. Let $N(D)$ be the
crossing count and let $b(D)$ count first encounters which are OVER.
The initialization is exactly $b=0$. A switch at the first bad crossing
keeps the parameter circles and their traversal order. It changes that
first encounter from OVER to UNDER and leaves every other first-encounter
designation unchanged. Thus $D^{\rm sw}$ has complexity $(N,b-1)$.
Its chosen crossing sign is negated, since the two ordered over/under
tangents are exchanged.

Oriented smoothing reconnects each incoming strand to the other strand's
outgoing end. It removes just the selected crossing and creates no other
crossing in its clean disc. A self crossing splits one parameter circle
into two; a mixed crossing joins two into one. This can be checked by
writing a self-crossing cycle as $(x A y B)$ and replacing it by the
two cycles $(x B),(y A)$ before erasing $x,y$. For a mixed crossing,
the cycles $(x A),(y B)$ become $(x B y A)$ before erasing the marks.
Empty resulting occurrence words are still circles. Thus the resulting
component counts are $c+1$, or $c-1\geq1$ in the mixed case, and the
smoothing always remains in the nonempty domain. Its crossing count is
$N-1$. New orders and basepoints can be chosen because the first
complexity coordinate has decreased. These facts justify lexicographic
induction on $(N,b)$, for every finite number of components.

At a positive chosen crossing the skein is
$aP_D^G-a^{-1}P_{D^{\rm sw}}^G=zP_{D^0}^G$.
Move the second term to the right and divide by the unit $a$ to obtain
\begin{equation}\label{lp:positive}
 P_D^G=a^{-2}P_{D^{\rm sw}}^G+a^{-1}zP_{D^0}^G.
\end{equation}
At a negative chosen crossing the positive diagram is $D^{\rm sw}$,
so the skein is $aP_{D^{\rm sw}}^G-a^{-1}P_D^G=zP_{D^0}^G$.
Move the latter two terms across and multiply by $a$. This gives
\begin{equation}\label{lp:negative}
 P_D^G=a^2P_{D^{\rm sw}}^G-azP_{D^0}^G.
\end{equation}

For integral descent let $M_c=z^{1-c}\mathbb Z[a^{\pm1},z^2]$,
viewed as an additive subgroup of $R_G$. The initialization belongs to
$M_c$, since $\delta^{c-1}=z^{1-c}(a-a^{-1})^{c-1}$.
The switched term has the same component count, so induction puts it in
$M_c$; multiplication by any $a$ unit preserves this subgroup. For a self
smoothing, induction puts the smoothed value in $M_{c+1}$, and
\begin{equation}\label{lp:self-support}
 zM_{c+1}=z\,z^{-c}\mathbb Z[a^{\pm1},z^2]=M_c.
\end{equation}
For a mixed smoothing, induction puts it in $M_{c-1}$, and
\begin{equation}\label{lp:mixed-support}
 zM_{c-1}=z\,z^{2-c}\mathbb Z[a^{\pm1},z^2]
             =z^2M_c\subseteq M_c.
\end{equation}
Both recurrences therefore put $P_D^G$ in $M_c$. The standard inclusion
$R\hookrightarrow R_G$ is injective (the Gaussian coefficient ring has
the free integer basis $1,\mathrm i$). There is consequently one $P_D\in R$
with image $P_D^G$. This proves support and integral descent simultaneously.
The local source equalities remain equalities after $\phi$ and, by
injectivity, are equalities in $R$.

For uniqueness, let $Q_D\in R$ satisfy the stated skein and every
UNDER-first initialization. Apply the same induction. At $b=0$ its value
is the specified $\delta^{c-1}=P_D$. Otherwise its value is given by the
same positive or negative solved recurrence. The switched and smoothed
values agree with $P$ by the inner and outer induction respectively.
Substitution gives $Q_D=P_D$. No existence or confluence is inferred from
this uniqueness argument; existence was the literal source construction.

Finally map $R$ to $\mathbb Q(a)$ by fixing the indeterminate $a$ and
sending $z$ to $a-a^{-1}$. This target is a field, and $a-a^{-1}$ is a
nonzero rational function, so the Laurent specialization is well-defined.
The initialized $\delta$ maps to $1$. The same induction shows that
every nonempty diagram specializes to $1$. In the positive case the two
inductively specialized smaller values are each $1$, and
\begin{equation}\label{lp:nonzero-positive}
 a^{-2}+a^{-1}(a-a^{-1})
       =a^{-2}+1-a^{-2}=1.
\end{equation}
In the negative case they are also each $1$, and
\begin{equation}\label{lp:nonzero-negative}
 a^2-a(a-a^{-1})=a^2-a^2+1=1.
\end{equation}
A zero element of $R$ would specialize to zero. This proves $P_D\ne0$.

Finally, $D\mapsto P_D$ and $D\mapsto H_D$ are maps from oriented link
diagrams into $R$ that take the value $1$ on the crossing-free circle,
satisfy the campaign skein~\eqref{lp:skein}, and are invariant under
planar isotopy and the three Reidemeister moves: for $P$ these properties
were proved above from Literature input~\ref{lp:lm}, and for $H$ they are
the existence and invariance clauses of Literature input~\ref{lit:homfly}.
Lemma~\ref{lp:coefficient-transport} identifies two such maps, so
$P_D=H_D$ on every actual diagram. This identification consumes the
uniqueness transported to $R$ by that lemma from Literature
input~\ref{lp:lm-uniqueness}, together with the existence clause of
Literature input~\ref{lit:homfly}; the uniqueness clause of the latter is
not used. It does not derive arbitrary ambient-isotopy invariance from the
local construction alone. The support proof itself used only the local
construction and the displayed induction.
\end{proof}

\begin{lemma}[A diagrammatically split circle]\label{lp:split-circle}
Let $D'$ be the actual diagram formed from $D$ by adding one simple
crossing-free component having no crossings with $D$. It may surround
some components of $D$; it need not lie in the unbounded complementary
face. Then
\begin{equation}\label{lp:split}
 P_{D'}=\delta P_D.
\end{equation}
In particular every crossing-free $c$-component diagram has value
$\delta^{c-1}$.
\status{proved (refereed: bench A, 2026-09-05T15:51:04Z; transcribed from RC d2b\_local\_polynomial, split circle (C026))}
\end{lemma}
\begin{proof}
Order and base $D$, and put the added component last with any basepoint.
This adds no crossing encounter, so $N$ and $b$ are unchanged. Induct on
$(N,b)$ for $D$. If $b=0$, the two source initialization values are
$\delta^{c-1}$ and $\delta^c$, which have the required relationship.
Otherwise choose the same first bad crossing in $D$ and $D'$. Switching
or smoothing does not touch the extra component. The switched pair and
the smoothed pair are therefore again related by adding exactly that
crossing-free component. In either solved recurrence the coefficients
are scalar elements of the commutative ring $R$. Substitute the inductive
factor $\delta$ for both smaller values of $D'$, and factor it out; the
remaining expression is the solved recurrence for $P_D$. This proves
the formula. Nesting never enters this induction. This lemma does not
claim that an arbitrary spatially split presentation already has this
diagram form; such a use requires its own certificate or product proof.
\end{proof}

\subsection{Scalar equality from named records}

We use the same source function $F_D$ from Literature input~\ref{lp:lm}
and its factor $\mu=-(l+l^{-1})m^{-1}$. Its UNDER-first initialization
includes every positive component count and every crossing-free circle.

\begin{lemma}[Polynomial equality from a named decorated record]
\label{rp:record-polynomial}
Let $D,D'$ be two actual nonempty finite generic oriented link diagrams.
Suppose there is a bijection of their components and crossing occurrences
which preserves oriented cyclic successor, crossing pairing, over/under
bits and crossing signs. The component bijection must also include all
components with no crossing occurrences. Then $F_D(l,m)=F_{D'}(l,m)$.
Consequently any common Laurent-ring substitution of these two source
values, and every coefficient of the substituted values, agrees.
This statement does not require or assert realizability of an arbitrary
abstract record or an ambient isotopy between the two diagrams.
\status{proved (refereed: bench A, 2026-09-05T15:51:04Z; transcribed from RC d2b/d2c, named-record polynomial equality (C026))}
\end{lemma}
\begin{proof}
Choose an order on the components of $D$, and one basepoint outside the
crossings on each component. Transfer the order to $D'$. On a component
with crossings, choose its new basepoint in the interval just preceding
the image of the old first crossing occurrence. On a crossing-free
component choose any point. These choices make the complete traversal
orders of crossing occurrences correspond. They are allowed because the
source construction is independent of component order and basepoint.

Call a crossing bad if its first encounter in this ordered traversal is
the overpass. Write $N$ for the total number of crossings and $b$ for the
number of bad crossings. Both numbers agree for the two based diagrams.
We use lexicographic induction on $(N,b)$. If $b=0$, both diagrams are
under-first ascending and have the same number $c\geq1$ of components.
Their two source values are therefore separately $\mu^{c-1}$ by the
initialization. This includes $N=0$ and any crossing-free components;
no value for an empty link is introduced.

For $b>0$, choose the first bad crossing in $D$ and its named image in
$D'$. Switching that crossing preserves each traversed component and its
basepoint. It swaps the two over/under bits at just that crossing.
Hence that crossing ceases to be bad and every other first-encounter
designation is unchanged: the switched diagrams have complexity
$(N,b-1)$ and still have isomorphic decorated records. Their values
are equal by the inner induction.

Oriented smoothing removes the two occurrences of the chosen crossing.
Cut each traversed strand there into an incoming and an outgoing end.
The oriented smoothing connects each incoming end to the outgoing end
of the other strand. The given bijection preserves these ends and their
orientation, so it preserves this precise recombination of successors.
All other paired crossings retain their over/under bits and signs.
The resulting actual smoothed diagrams thus have isomorphic decorated
records, including their component bijection. A self-crossing splits
one traversed component into two; a mixed crossing joins two into one.
Crossing-free resulting circles are retained as components rather than
discarded. The smoothed diagram is always nonempty: in the mixed case
the original component number is at least two, and in the self case
the number increases. The new crossing number is $N-1$.

Choose new component orders and basepoints on one smoothed diagram and
transfer them through its induced record bijection as above. The outer
induction applies regardless of their new bad-crossing count. It gives
equal values for the two smoothed diagrams. No relation between the old
and new component orders is needed, and no topological classification
is used to supply the record bijection.

It remains to insert these two equalities into the same scalar recursion.
The crossing signs are part of the record, so the chosen crossings in
$D,D'$ have the same sign. For a positive chosen crossing the source
skein reads
\begin{equation}\label{rp:positive-skein}
 lF_D+l^{-1}F_{D^{\rm sw}}+mF_{D^0}=0.
\end{equation}
Subtract the last two terms and divide by the unit $l$. This gives
\begin{equation}\label{rp:positive-recursion}
 F_D=-l^{-2}F_{D^{\rm sw}}-l^{-1}mF_{D^0}.
\end{equation}
For a negative chosen crossing the positive diagram is the switched one,
and the source skein instead reads
\begin{equation}\label{rp:negative-skein}
 lF_{D^{\rm sw}}+l^{-1}F_D+mF_{D^0}=0.
\end{equation}
Subtract the other two terms and multiply by $l$, obtaining
\begin{equation}\label{rp:negative-recursion}
 F_D=-l^2F_{D^{\rm sw}}-lmF_{D^0}.
\end{equation}
The identical formula with primes holds in the corresponding sign case.
The two switched terms agree by the inner induction and the two smoothed
terms by the outer induction. Substitution therefore proves $F_D=F_{D'}$.
This completes the induction, including both signs and both smoothing
component cases. No division by a polynomial value is performed.

A common Laurent-ring homomorphism preserves equality. Coefficient
extraction is an additive map defined also at zero and therefore preserves
that resulting equality. These give the last assertions.
\end{proof}

\begin{lemma}[Presentations of the same decorated record]
\label{lc:presentations}
Suppose two actual finite nonempty decorated diagrams have a specified
bijection of their oriented parameter circles and crossing occurrences,
preserving cyclic successor, pairing, over/under designations and signs.
Their evaluations by the same local LM construction agree. This includes
positive page-chart changes of an actual spherical diagram, clean
crossing-free replacements and different compatible height choices,
whenever the endpoint page diagrams belong to the finite-page domain.
All crossing-free circles must be included in the component bijection.
No arbitrary ambient isotopy is a hypothesis or conclusion of this scalar
statement.
\status{proved (refereed: bench A, 2026-09-05T19:35:00Z; transcribed from RC d2b\_local\_polynomial, lc:presentations (C026))}
\end{lemma}
\begin{proof}
Apply Lemma~\ref{rp:record-polynomial} and the common substitution
$l=\mathrm ia$, $m=-\mathrm iz$. For the stated presentation examples,
the bijections can be checked directly. A positive regular page chart
does not alter a parameter circle or its order; its positive differential
multiplies a crossing determinant by a positive number. An orientation
preserving finite piecewise-linear chart has the same cyclic ray order,
so after its regular finite-page model the signs are the same. The
specified crossing decorations retain the over/under bits. Choosing a
different point in a face to omit inserts no crossing occurrence.
These are comparisons of two actual diagrams, not a realization of an
arbitrary word or a claim that a continuous chart is differentiable.

A clean replacement of one crossing-free interval, fixed on its endpoint
collars and missing every other strand, inserts no visit. The original
circle parameter therefore supplies the same successor between its two
collars and identifies every old crossing germ. Each crossing-free circle
is explicitly retained. Finally, compatible heights mean precisely that
the prescribed over branch remains higher than the under branch in the
chosen positive projection frame. Changing those heights changes none
of the decorated record, even if the heights themselves are only
continuous. No smoothness of such a height change, or polynomial
invariance under its ambient extension, is inferred.
\end{proof}

\begin{lemma}[A single self crossing has scalar value one]
\label{lc:single-crossing}
An actual oriented one-circle diagram with just one self crossing has
local LM evaluation $P_D=1$, for either crossing sign. A crossing-free
one-circle diagram also has value one.
\status{proved (refereed: bench A, 2026-09-05T19:35:00Z; transcribed from RC d2b\_local\_polynomial, lc:single-crossing (C026))}
\end{lemma}
\begin{proof}
In the one-crossing case choose the nonsingular basepoint in the cyclic
interval immediately before its under occurrence. That occurrence is
first; there is no other crossing. The diagram is therefore UNDER-first
with one component and has source value $\mu^0=1$. The common Laurent
substitution keeps this value. The zero-crossing case is the same
one-component initialization. This does not assume that a selected
planar monogon is the bounded face, or discard another component.
\end{proof}


\subsection{Marked products and mixed rows}

We use the solved recurrences and nonempty smoothing domain of
Theorem~\ref{lp:core}, and the named-record and presentation equalities
of Lemmas~\ref{rp:record-polynomial} and~\ref{lc:presentations}.
The marked operations below do not assert a relative tangle classification.

\paragraph{Marked joining operation and its actual realizations.}
A marked diagram is a nonempty actual diagram with a specified closed
nonsingular oriented interval $I$ on one component; $I$ contains no
crossing and is contained in a clean disc. All crossing discs used in a
recursion are chosen disjoint from $I$. A marked join of $(A,I_A)$ and
$(B,I_B)$ cuts a smaller open interval inside each mark and joins the
oriented ends crosswise, preserving the surviving collars. There are no
additional crossings in the joining arcs. Its finite data are precise:
concatenate the marked component cycles at their marked gaps, retain
every other component, and keep exactly all old crossing pairings, bits
and signs. The component number is $c(A)+c(B)-1$. The two old marks in
the separate factors remain distinguished during recursion; a smoothing
selects the unique resulting component containing the relevant interval.
The operation is not defined as an unspecified ambient isotopy class.

These actual clean joins exist, including at an interior-face mark.
Here are sufficient planar constructions. Choose a point in a face
adjacent to $I_A$, off the diagram, and regard the page as an oriented
sphere. Use an orientation preserving chart omitting that point. The
marked interval is now adjacent to the unbounded face; the finite
decorated data have not changed, so Lemma~\ref{lc:presentations} preserves its
scalar. If a polygonal presentation is desired, the finite crossing-disc
and disjoint-strip construction of Lemma~\ref{lem:gauss-pl-model}
gives one without changing those data; its ambient conclusion is not
used here. Carry the interval and its endpoint collars with the local
page maps; this preliminary change of presentation is not claimed to
fix their original coordinates.

In the unbounded face choose a simple polygonal path from a point just
beside the interval to outside a rectangle containing the diagram.
Such a finite path exists: the points reachable from a given point by
finite polygonal paths in an open connected set form an open and
relatively closed subset. A polygonal path can be made simple by a small
perturbation and removal of its finitely many loops. At its initial end
use the clean interval disc. Its remaining compact image is disjoint
from the diagram and hence has positive clearance. A sufficiently thin
ribbon about the path therefore misses the diagram; at its finitely
many corners use disjoint small discs. Two parallel sides of this
ribbon extend the cut ends to the exterior rectangle without crossing
one another or the diagram. No old crossing is changed and no new one
is added. The remaining diagram is thus presented in a rectangle with
the two oriented cut ends accessible from its boundary.

Shrink this rectangle by a positive affine map into the clean disc of
$I_B$ after cutting its smaller subinterval. Match the two ends to the
two collars of $B$ with disjoint arcs, orienting the inserted traversal
from the incoming collar to the outgoing collar. There is no reversal
of $A$: label its incoming and outgoing cut ends accordingly, and use
an orientation preserving planar placement. The exterior of that clean
disc in $B$ is literally unchanged. This realizes the stated record
concatenation. The old $A$ presentation, the new one and any other
actual clean realization are compared only by Lemma~\ref{lc:presentations}.
In particular moving the omitted face has not been called a relative
ambient isotopy fixing an arbitrary exterior. Any finite list of other
marks can be kept out of the local crossing/join discs by shrinking
them first; their cyclic positions and collars survive.

\begin{theorem}[Simultaneous marked product]\label{mp:join}
For two nonempty actual marked diagrams and any actual clean marked
join $J(A,B)$ just specified,
\begin{equation}\label{mp:join-value}
 P_{J(A,B)}=P_AP_B.
\end{equation}
This includes smoothing-created multi-component factors and an arbitrary
crossing-free marked component. No assertion about a marked terminal
tangle being ambient-trivial is needed.
\status{proved (refereed: bench A, 2026-09-05T15:51:04Z; transcribed from RC d2c\_polynomial\_products, marked join (C026))}
\end{theorem}
\begin{proof}
Use induction on the sum $N=N(A)+N(B)$, with an inner induction on the
sum $b=b(A)+b(B)$ of bad crossings for chosen based orders. In EACH
factor put the marked component first and base it at its marked gap.
Choose arbitrary orders and bases for the remaining components.

If $b=0$, traverse the joined component starting just before the $A$
portion, then its $B$ portion. Traverse the remaining $A$ components
in their chosen order, then the remaining $B$ components in their
chosen order. For crossings of $A$, the relative order of all visits is
the same as in its factor traversal; visits of $B$ inserted between
some of them do not change which of the two visits to an $A$ crossing
comes first. The identical observation holds for $B$. There are no
crossings with one visit in each factor. Thus the entire joined diagram
is UNDER-first. Its source base gives
\begin{equation}\label{mp:join-base}
 P_{J(A,B)}=\delta^{c(A)+c(B)-2}
             =\delta^{c(A)-1}\delta^{c(B)-1}=P_AP_B.
\end{equation}
This is a scalar computation on an actual ascending diagram, not a
claim identifying its marked arcs with a trivial tangle.

If $b>0$, choose a bad crossing in one factor, say $A$. Its disc is
disjoint from the mark and from the joining operation. Switching it
gives $A^{\rm sw}$ with the same component number, same marked interval
and smaller $b$. In $J(A,B)$ it is precisely the corresponding crossing
switch. Smoothing gives $A^0$, of crossing number $N(A)-1$, with the
same interval on one unique component. In the joint diagram, smoothing
has exactly the record of $J(A^0,B)$: both operations are recombinations
at disjoint incoming/outgoing ends. All other successor relations and
all other crossings are unchanged. If the selected crossing involves the marked component, a self-smoothing
places the interval on exactly one of its two output components, and a
mixed smoothing joins that component to another while retaining the
interval. If the selected crossing involves no strand of the marked
component, that component and its interval are unchanged: a self-smoothing
splits one unmarked component into two, and a mixed smoothing merges two
unmarked components. In all four cases every output component is retained,
including all crossing-free components. These are the exhaustive
self/mixed and marked/unmarked cases; nonemptiness was checked in
Theorem~\ref{lp:core}.

Use the actual smoothed diagram as the realization of $J(A^0,B)$, or
apply Lemma~\ref{lc:presentations} to a newly chosen clean realization. Inner
induction gives $P_{J(A^{\rm sw},B)}=P_{A^{\rm sw}}P_B$; outer
induction gives $P_{J(A^0,B)}=P_{A^0}P_B$. Apply the appropriate solved
skein to the joint diagram and substitute these two equalities. Its
right side is respectively
$(a^{-2}P_{A^{\rm sw}}+a^{-1}zP_{A^0})P_B$ or
$(a^2P_{A^{\rm sw}}-azP_{A^0})P_B$. The factor in parentheses is
$P_A$ by the same solved skein in $A$. This proves the identity. If
the selected bad crossing is in $B$, the identical argument with the
two factor names interchanged applies. For each outer induction case
new auxiliary choices can again put its marked component first,
because the source value is independent of those choices.
\end{proof}

\begin{theorem}[Ordered stacked blocks, including split unions]
\label{mp:stack}
Let $D$ be an actual nonempty diagram whose components are partitioned
into $q\geq1$ nonempty tagged blocks. Assume that at every crossing
between different blocks the smaller-index block is under the larger
one. Let $D_i$ be the actual restriction retaining all components in
block $i$ and all crossings internal to it. Then
\begin{equation}\label{mp:stack-value}
 P_D=\delta^{q-1}\prod_{i=1}^qP_{D_i}.
\end{equation}
There is no bound on the number of components in a block, no restriction
on crossings inside a block, and no requirement that their page images
be separated. A crossing-free disjoint union is the special case with
no crossings between blocks. In particular
$P_{A\sqcup B}=\delta P_AP_B$ for two nonempty diagrams presented
without mixed crossings.
\status{proved (refereed: bench A, 2026-09-05T15:51:04Z; transcribed from RC d2c\_polynomial\_products, stack (C026))}
\end{theorem}
\begin{proof}
Order components block by block, choosing any based order within each.
All between-block crossings are first encountered under. Let $N$ be
the sum of the INTERNAL crossing numbers and $b$ the sum of their bad
crossing numbers. These need not be the total crossing number of $D$.

If $b=0$, the full diagram is ascending. With $c_i=c(D_i)$ its scalar
is $\delta^{\sum_i c_i-1}$. Each factor has scalar
$\delta^{c_i-1}$, so their product times $\delta^{q-1}$ has the same
exponent $q-1+\sum_i(c_i-1)=\sum_i c_i-1$.

Otherwise resolve a bad internal crossing. Switching it lowers $b$,
preserves the block assignment and the between-block hypothesis, and
is the same switch in its block restriction. Smoothing lowers $N$.
Every new component inherits the same block as the strands smoothed;
no two blocks merge, and no block becomes empty. All between-block
crossing visits persist on their original strands, so their under/over
relation still has the required block order. The restriction of the
smoothed diagram in that block is exactly the oriented smoothing of
the old block restriction; the other restrictions do not change.
Lexicographic induction and the appropriate solved skein now give
\eqref{mp:stack-value}, since its factors outside the resolved block
and the scalar $\delta^{q-1}$ are identical in all three terms.
This proof supplies the scalar without moving any component in space.
\end{proof}

\begin{lemma}[Mixed signed crossings in a stack]\label{mp:zero-link}
For two distinct components of an actual generic oriented plane diagram,
the sum of $\operatorname{sgn}\det(u_1,u_2)$ over their transverse
intersections, in this fixed component order, is zero. Consequently
the half-sum of decorated crossing signs is an integer, and it is zero
if one component is always over the other.
\status{proved (refereed: bench A, 2026-09-05T15:51:04Z; transcribed from RC d2c\_polynomial\_products (C026))}
\end{lemma}
\begin{proof}
Use finite polygonal page models preserving all transverse crossing
signs; the clean crossing-disc and strip construction cited above
does not change any signed intersection. Write the oriented edges of
component 2 as $[v_j,v_{j+1}]$. Choose an auxiliary point $o$ off the
finitely many lines and intersections that would make a fan triangle
degenerate or a radial edge nongeneric against component 1. A finite
union of proper lines and points cannot fill the plane, so such $o$
exists. Include further finite forbidden lines through a vertex of
component 1 and a vertex of component 2 to exclude radial-edge hits
at the former vertices. Orient the triangle chain
$[o,v_j,v_{j+1}]$ by this order. Its coefficient relative to the usual
positive plane orientation is its determinant sign.

For one positively oriented triangle, a closed transverse path has
as many entries into its interior as exits, counted with sign. At an
entry the ordered tangent pair (path, triangle boundary) has the
opposite determinant sign from an exit: the triangle lies consistently
on the left of its positive boundary. Thus its algebraic intersection
sum with that boundary is zero. Reversing triangle orientation negates
the sum and still gives zero. Summing over the finitely many fan
triangles, each radial edge occurs once in each orientation and
cancels. Their boundary chain is exactly component 2. Hence the
intersection sum of components 1 and 2 is zero. This remains valid for
self-intersecting components: the argument counts traversed occurrences,
not distinct points of an embedded boundary for either component.

At each mixed crossing the decorated sign is either the fixed-order
determinant sign or its negative, according to which component is
over. Their total differs from the zero fixed-order sum by twice an
integer. Its half is therefore integral. If the over-component is
constant, all signs use the same one of those two conventions and
the total is zero. This proves the assertions using the original mixed
crossings of this coambient pair, not a replacement split pair.
\end{proof}

\begin{theorem}[All-component lowest mixed row]\label{mp:lowest}
Let $D$ be an actual diagram with $c\geq1$ tagged oriented components,
and let $D_i$ be their actual knot restrictions. Define
$\ell_{ij}=\tfrac12\sum\sigma(x)$ using ALL mixed crossings between
the original components $i,j$, and put $\Lambda=\sum_{i<j}\ell_{ij}$.
Then
\begin{equation}\label{mp:lowest-value}
 [z^{1-c}]P_D
 =a^{-2\Lambda}(a-a^{-1})^{c-1}
       \prod_{i=1}^c[z^0]P_{D_i}.
\end{equation}
For $c=2$ this is the two-component mixed row; the two
tagged restrictions are intrinsic knot diagrams while $\ell_{12}$ is
computed in their original common diagram.
\status{proved (refereed: bench A, 2026-09-05T15:51:04Z; transcribed from RC d2c\_polynomial\_products, lowest mixed row (C026))}
\end{theorem}
\begin{proof}
Set $h(D)=[z^{1-c}]P_D$. At a mixed crossing the smoothed diagram has
$c-1\geq1$ components and support at least $z^{2-c}$ by
Theorem~\ref{lp:core}. Its multiplication by $z$ therefore has
support at least $z^{3-c}$, so contributes nothing to the $1-c$ row
of the skein. Hence $a h(D_+)=a^{-1}h(D_-)$ and
$h(D_-)=a^2h(D_+)$. The switch from positive to negative changes one
mixed sign from $+1$ to $-1$, and hence changes $\Lambda$ by $-1$.
It follows that $a^{2\Lambda}h$ is unchanged by that switch. Reversing
the equality proves the same for a negative-to-positive switch.

Switch exactly the mixed crossings necessary to put each smaller-index
component UNDER every larger-index component. These are finitely many
switches of an actual diagram, and no self-crossing or intrinsic
component restriction changes. By Lemma~\ref{mp:zero-link} every pair
of final components has linking number zero, hence final $\Lambda=0$.
Apply Theorem~\ref{mp:stack} with each component as one block. Its
value is $\delta^{c-1}\prod_iP_{D_i}$. Knot support is nonnegative
and even in $z$, so the $1-c$ coefficient of this expression is exactly
$(a-a^{-1})^{c-1}\prod_i[z^0]P_{D_i}$: any positive $z$ power in one
factor raises that exponent. Restoring the preserved weight
$a^{2\Lambda}$ proves the formula. For $c=1$ there are no switches,
$\Lambda=0$, and the assertion is the identity $[z^0]P_D=[z^0]P_D$.
No coefficient or whole polynomial has been assumed nonzero or divided out.
\end{proof}

\begin{lemma}[Nested connected interlacement blocks]\label{mp:blocks}
Let an actual oriented one-circle decorated record have a nonempty
crossing set, partitioned into the connected components of its
interlacement graph. Suppose that for every component $H$ an actual retained diagram
$C_H$ with exactly that restricted named cyclic record is supplied.
Then a finite succession of clean marked joins of these actual diagrams
realizes the full record, and every actual diagram with that full record
has polynomial $\prod_H P_{C_H}$. The joins preserve the sign of every
crossing and the writhe is the sum of the writhes of $C_H$.
\status{proved (refereed: bench A, 2026-09-05T15:51:04Z; transcribed from RC d2c\_polynomial\_products (C026))}
\end{lemma}
\begin{proof}
First fix a connected component $A$ of the interlacement graph and a
chord $b$ outside it. Cutting the circle at the endpoints of $b$ gives
two open intervals. Since $b$ interlaces no chord of $A$, each chord
of $A$ has both endpoints in one of them. If chords of $A$ occupied
both intervals, an interlacement path in $A$ would have an adjacent
pair in different intervals. Such chords cannot alternate. Thus all
endpoints of $A$ lie in just one interval, which means both endpoints
of $b$ lie in one cyclic gap between successive endpoints of $A$.

If $b,b'$ are interlacing chords outside $A$, they must occupy the same
gap of $A$: distinct gaps are disjoint intervals and cannot contain
alternating endpoints. Applying this along an interlacement path shows
that every other connected component $B$ has all its endpoints in one
gap of $A$. This argument needs connectedness of $A$ as well as $B$;
noninterlacement of individual chords alone would not suffice.

Choose a base gap for the full cyclic record and let $A$ be the block
containing its first occurrence. The other blocks lie in gaps of $A$.
Within a nonempty such linear gap take the block $B$ containing its
first occurrence. Every other block inside the interval lies either in
an internal gap of $B$ or after its last occurrence: it cannot precede
its first occurrence, and the one-gap property excludes a block
straddling two such gaps. Recurse on the internal gaps, and then on
the suffix after that last occurrence. The same description for the
outer cyclic gap is made by cutting at the chosen base gap. Each step
removes at least one nonempty connected block from the pending list,
so the resulting finite nested forest, attached to root $A$, terminates.

Start with the actual $C_A$. For each gap attach the recursively
constructed actual diagrams for that gap in their inherited order,
using small disjoint crossing-free intervals in that gap. Each clean
join concatenates just the chosen cyclic blocks and introduces no new
crossing. Multiple insertions can use successively smaller intervals;
there are only finitely many. The construction of marked joins above
supplies actual planar realizations at every step, including nested
and exterior gaps. Thus the final oriented named record is literally
the full given record, with no assumption that an arbitrary restricted
word is realizable. Every leaf was an already supplied actual diagram.

Repeated use of Theorem~\ref{mp:join} gives the product value for the
constructed diagram. Lemma~\ref{lc:presentations} equates it with any actual
diagram having the same full data. Every old crossing is present once,
with its old sign, and no joining arc has a crossing, so writhe adds.
In the canonical application each sign is positive; its sum is
$\sum_H|H|$. If there are no owned blocks, the grouped value remains
the separately defined algebraic empty product $1$ and empty sum $0$.
No empty link is introduced and no knot is asserted to represent that
empty product by this lemma.
\end{proof}



\begin{definition}[corner coefficient and the state sum]\label{def:C}
For a subpolygon $Q$ of a decomposition of the generic polygon $P$, let
$H^+_Q$ be the HOMFLY--PT polynomial of its positive lift and put
\[
d_Q=1-m_Q-|r_Q|,\qquad c(Q)=\bigl[a^{d_Q}z^0\bigr]H^+_Q(a,z)\in\ZZ,
\]
the coefficient of $a^{d_Q}z^0$ (zero if that monomial is absent). The
\emph{corner state sum} of $P$ is
\[
C(P)=(-1)^{\ell(P)}\sum_{\substack{S\in\Ind(G_P)\\ S\text{ uniform}}}
(-1)^{|S|}\prod_{Q\text{ subpolygon of }S}c(Q).
\]
\end{definition}

\begin{definition}[the state sum on the weakly generic locus]\label{def:C-weak}
For a weakly generic polygon $P\in\mathcal W_n$, form its decompositions
and carriers from its actual crossings as in
Lemma~\ref{lem:weak-carriers}; let $H^+_Q$ be the HOMFLY--PT polynomial of
the positive lift (Definition~\ref{def:positive-lift}) of a carrier $Q$,
and put $d_Q=1-m_Q-|r_Q|$ and $c(Q)=[a^{d_Q}z^0]H^+_Q(a,z)$, with $r_Q$,
$m_Q$ and uniformity the data of Definition~\ref{def:uniform}, defined by
that lemma. The \emph{corner state sum} of $P$ is
\[
C(P)=(-1)^{\ell(P)}\sum_{\substack{S\in\Ind(G_P)\\ S\text{ uniform}}}
(-1)^{|S|}\prod_{Q\text{ carrier of }S}c(Q),
\]
$\ell(P)$ the number of left turns of $P$. On $\mathcal U_n$ this is
Definition~\ref{def:C} term by term, and the same symbol $C$ is used. No
value of an amplitude or of a gate at zero enters this definition.
\end{definition}

\begin{lemma}[local constancy of the weak state sum]\label{lem:C-weak-local}
The function $C$ of Definition~\ref{def:C-weak} is invariant under the
cyclic shift and locally constant on $\mathcal W_n$; hence it is constant
on every visible chamber. It is the unique locally constant function on
$\mathcal W_n$ that agrees with Definition~\ref{def:C} on $\mathcal U_n$.
\status{new (round~3; adopted from the Codex proposal W-P2, P-25-13, lane sm2\_weak\_c\_v1, re-read by the author; the identification of the carrier polynomials with $H$ goes through Theorem~\ref{lp:core}, unrefereed; no Hypothesis~R; F-25-90)}
\end{lemma}
\begin{proof}
Fix a labelled weakly generic tuple $P$ and, by Lemma~\ref{lem:silentlocus},
a relative neighbourhood $N$ of $P$ in $\mathcal W_n$ on which the visible
signature is constant. On $N$ the original vertices and the crossings carry
fixed names. For a transverse crossing the two edge parameters solve a
two-by-two linear system with nonzero determinant and are continuous
functions of the parent vertices; after shrinking $N$ their strict order on
each edge and their membership in the open edge intervals persist. The
marked successor data, the independent sets and the cycle lists of
Lemma~\ref{lem:weak-carriers} are therefore the same for every tuple in
$N$, and every carrier has a fixed finite corner list --- original vertices
and selected crossing points --- whose coordinates depend continuously on
the parent.

At $P$ each such list lies in the regular locus of its own arity, by
Lemma~\ref{lem:weak-carriers}(ii). There are finitely many independent sets
and carriers; for each corner list choose a ball about its value at $P$
inside that regular locus, which is open, and shrink $N$ so that the
continuous image of $N$ lies in the ball. A ball is convex, so
Lemma~\ref{lem:rot}(ii) applied to straight paths in it makes the rotation
of every carrier constant on $N$. All corner determinants are nonzero at
$P$ by Lemma~\ref{lem:weak-carriers}(ii); shrink $N$ once more so that
their signs are constant. Hence the uniform/mixed classification, every
$r_Q$, every $m_Q$ (the number of unselected crossings with both visits on
the cycle, fixed with the cycle lists) and $\ell(P)$ are constant on $N$.

The retained double points of a carrier are exactly the unselected
crossings named by its cycle list, in a fixed traversal order. At each of
them the two strand vectors are positive multiples of the corresponding
parent edge vectors, whose nonzero determinant keeps its sign on $N$; so
the over strand of the positive lift is the same named strand throughout
$N$, and every crossing sign is $+1$. The named decorated record of each
actual positive carrier diagram --- oriented cyclic successor, crossing
pairing, over/under bits and signs, including the one-component record of
a carrier without crossings --- is thus the same at every tuple of $N$.
Lemma~\ref{rp:record-polynomial} gives equality of the source polynomials
of corresponding carrier diagrams, the common Laurent substitution of
Theorem~\ref{lp:core} gives equality of their $P$-polynomials, and that
theorem identifies these with their $H$-polynomials. Since $d_Q$ is
constant, every coefficient $c(Q)$ is constant on $N$, and so is the finite
sum of Definition~\ref{def:C-weak}. This proves local constancy on the
labelled locus, at simultaneous silent zeros as well as at simple ones: no
factorization of a path into wall crossings was used.

Cyclic relabelling renames the crossings, the independent sets and the
traced oriented carrier diagrams without changing their turns, rotations,
positive lifts or coefficients, or $\ell(P)$; so $C(\sigma P)=C(P)$. The
union of the cyclic translates of a neighbourhood on which $C$ is constant
is a saturated open set on which $C$ has that value, so $C$ is locally
constant on the quotient as well. A locally constant function is constant
on a connected set, since the preimage of a value and its complement are
both open; this gives constancy on every visible chamber.

On $\mathcal U_n$ the definition agrees with Definition~\ref{def:C} term by
term. If $C'$ is another locally constant function on $\mathcal W_n$ with
that restriction, then at any $P\in\mathcal W_n$ the intersection of
neighbourhoods on which $C$ and $C'$ are constant meets $\mathcal U_n$, by
the density assertion of Lemma~\ref{lem:silentlocus}, and at a generic
point of it $C=C'$; hence $C'(P)=C(P)$.
\end{proof}

\begin{lemma}[selector form of the carrier state sum]\label{lem:C-X1}
For a carrier $L$ of an independent support $S$, let $k(L)$ be its number
of corners. Put $\mathrm{wt}(L)=1$ if all its turns are right,
$(-1)^{k(L)}$ if all are left, and $0$ if its turns are mixed.
Set $\operatorname{wind}(S)=\prod_L\mathrm{wt}(L)$. Then
\begin{equation}\label{eq:C-selector-form}
C(P)=\sum_{S\in\Ind(G_P)}\operatorname{wind}(S)\prod_L c(L),
\end{equation}
where $c(L)$ remains the coefficient of the actual positive carrier lift
in Definition~\ref{def:C}.
\status{new (restricted in round 1 to the selector identity, C018; F-25-38)}
\end{lemma}
\begin{proof}
Every original vertex belongs to exactly one carrier and retains its turn.
Every selected crossing contributes one left and one right smoothing
corner by Lemma~\ref{lem:carriers}. Thus the total number of left carrier
corners is $\ell(P)+|S|$. If $S$ is uniform, precisely its all-left
carriers contribute to this count. Their selector product is therefore
$(-1)^{\ell(P)+|S|}$, the sign of this support in Definition~\ref{def:C}.
If $S$ is not uniform, a mixed carrier makes its selector product zero;
this support is absent from that definition. The coefficient factors
are unchanged in both cases. Summing these term identities proves
\eqref{eq:C-selector-form}.
\end{proof}

\begin{remark}[well-definedness]
The positive lift of a subpolygon is an honest knot diagram on an actual
closed plane curve (Lemma~\ref{lem:carriers}); no realization of an abstract
Gauss word is involved, and Literature input~\ref{lit:homfly} is consumed
only on such diagrams. The selector lemma is an identity for this
definition. It asserts no equivalence with a residual-piece polynomial
construction in a frozen source; any use of such an equivalence would
require a separate printed factorization proof.
\end{remark}

\subsection{The polynomial bound for an individual front}
\label{ng:front-section}

\begin{definition}[Finite oriented fronts]\label{ng:front-domain}
A front $F$ is an actual map of a nonempty finite union of parameter
circles to the oriented $(x,z)$ plane, with finitely many transverse
double points and ordinary semicubical cusps, no other singularities,
and no vertical tangencies on regular arcs. The limiting tangent at every
cusp is also nonvertical: in a semicubical parameter $u$ with cusp at
$u=0$, require $x''(0)\ne0$. Cusps meet no other strand or singularity.
At a crossing the branch with smaller $dz/dx$ is over. A downward cusp is
traversed from its locally upper arm to its locally lower arm. Write
$D(F)$ for the number of downward cusps, $w(F)$ for the sum of the
over-first tangent-determinant crossing signs, and $s(F)$ for the number
of crossings plus cusps.

In disjoint clean cusp discs replace each cusp by a simple regular arc
with the same oriented attachments and no crossing. The resulting
ordinary diagram is denoted $S(F)$.
\end{definition}

\begin{lemma}[independence of the rounding]\label{ng:smoothing-record}
Let $F$ be a front on the domain of Definition~\ref{ng:front-domain}. Any two
ordinary diagrams $S(F)$ obtained by the cusp replacement of that definition
have the same full named record: the same component circles, including
crossing-free ones, and the same crossing occurrences, cyclic orders,
over/under bits and signs. Consequently $P_{S(F)}$ does not depend on the
choice of rounding.
\status{new (round~5: the prose paragraph after Definition~\ref{ng:front-domain} wrapped as a statement so that the users of the symbol $P_{S(F)}$ cite it; F-25-136)}
\end{lemma}
\begin{proof}
A cusp replacement inside a clean cusp disc creates no crossing and changes
no successor of an old crossing visit along the oriented parameter circle:
the replacing arc has the same oriented attachments as the cusp it replaces.
Every crossing pairing, sign and over/under bit is determined by germs
outside the cusp discs, which are unchanged, and the component
correspondence is the identity on the parameter circles, including those
without crossings. Thus any two such diagrams have a named decorated-record
isomorphism, and Lemma~\ref{rp:record-polynomial} identifies their
polynomial values. This does not use an ambient identification with a
spatial cusp-shaped link.
\end{proof}

The polynomial used here is the same Lickorish--Millett construction in
the campaign variables. Its precise inputs, from
Literature input~\ref{lp:lm}, Theorem~\ref{lp:core}, and
Lemmas~\ref{lp:split-circle} and~\ref{rp:record-polynomial}, are:
\begin{enumerate}
\item For every actual nonempty finite oriented ordinary diagram $E$,
  $P_E$ is a nonzero element of $\mathbb Z[a^{\pm1},z^{\pm1}]$.
\item Positive page isotopy and each supplied ordinary oriented
  Reidemeister-I, -II or -III move preserve $P_E$.
\item Full oriented named records determine the value, with all component
  cycles retained.
\item An actual oriented skein triple satisfies
  $aP_+-a^{-1}P_-=zP_0$.
\item $P_U=1$ for a simple crossing-free circle, and
  $P_{E\sqcup U}=\delta P_E$, where
  $\delta=(a-a^{-1})/z$, if $E$ is nonempty and $U$ is a simple
  crossing-free component without mixed crossings. Nesting is allowed.
\end{enumerate}
No knot-type equality, arbitrary ambient isotopy, global Legendrian
completeness, mirror or reversal symmetry, braid representation theorem,
or contact-neighborhood theorem is an input to this argument.

\begin{definition}[degrees in $a$]\label{def:adeg}
For a nonzero Laurent polynomial $f$ in $a$ with coefficients in the
integral domain $\ZZ[z^{\pm1}]$, $\deg_af=\max\deg_af$ denotes the largest
$a$-exponent with nonzero coefficient and $\mindeg_af$ the smallest; both
are integers. The same symbols with the subscript $z$ denote the largest and
smallest $z$-exponents of a nonzero element of $\ZZ[a^{\pm1},z^{\pm1}]$.
\end{definition}

With Definition~\ref{def:adeg}, set
\begin{equation}\label{ng:defect}
 d(F)=\deg_a P_{S(F)},\qquad
 B(F)=D(F)-w(F)-d(F)-1.
\end{equation}
We will prove $B(F)\geq0$ for each front, not a statement about a maximum
over representatives.

\subsubsection{Ordinary certificates for the local front moves}
\label{ng:local-certificates}

Use Rutherford's elementary front words: $l_m$ inserts a left cusp in
positions $m,m+1$, $r_m$ joins those positions at a right cusp, and
$\sigma_m$ crosses them. Number cut positions from top to bottom, read
words from left to right, and type each operation by its incoming and
outgoing strand counts. In the standard crossing the incoming upper
strand travels down to the lower output and is over; the incoming lower
strand travels up to the upper output and is under. With both arrows
pointing right, the over-first tangent representatives $(1,-1),(1,1)$
have determinant $2$, so this crossing is positive.

Every operation below has an empty support rectangle, fixed ordered
boundary attachments, and an unchanged common exterior. All exterior
crossings, including mixed crossings, and all components are retained.
The permitted orientations need not produce a single component.

\begin{lemma}[Commutations and nonsingular deformation]
\label{ng:commutation}
Disjoint-gadget commutations and deformations through fronts without a
singular event preserve $D,w,d$, and hence $B$. Every supplied finite
front can be represented by a finite elementary front word.
\status{proved (refereed: bench A, 2026-09-05T15:51:04Z; transcribed from RC ng\_front (C029))}
\end{lemma}
\begin{proof}
Disjoint gadgets commute by sliding their empty rectangles. The index
changes track the same physical positions after insertion or removal of
two strands. After cusp rounding this is a positive page isotopy;
equivalently, the full named records are the same. The cusp directions
and crossing signs are unchanged, so $D$ and $w$ are unchanged as well.

A deformation without a singular event preserves the records, $D$ and
$w$: signs, cusp directions and cyclic attachments cannot change.
Lemma~\ref{rp:record-polynomial} gives scalar equality, without a global
move-generation theorem.

If singularities have equal $x$ coordinates, separate them by small local
$x$ translations, constant near the singularities and interpolated on
the intervening graph arcs. There are finitely many clean discs and
compact regular graph pieces. Their nonzero $x$ derivatives and
transverse crossings persist for sufficiently small translations, and
constant translations preserve the cusp germs. Reading successive
vertical cuts then gives the finite elementary word. No classification
of fronts is used.
\end{proof}

\begin{lemma}[Front type I]\label{ng:front-I}
The front type-I moves preserve $B$.
\status{proved (refereed: bench A, 2026-09-05T15:51:04Z; transcribed from RC ng\_front (C029))}
\end{lemma}
\begin{proof}
Each word $l_m\sigma_{m-1}r_m$ or $l_m\sigma_{m+1}r_m$ replaces one
through-strand. After rounding its two cusps, the unique crossing bounds
an empty ordinary monogon. In the first word, the strand crosses over
heading right, passes the right cusp and then the left cusp, and returns
under heading right. Both crossing arrows are therefore rightward, or
both are leftward when the entire strand is reversed. The other word
interchanges the over/under order along the through-strand but again
has simultaneously directed crossing tangents. Its crossing is also
positive. Exactly one of the two cusps is downward and one is upward.
Ordinary Reidemeister I gives, in the direction creating this front curl,
\begin{equation}\label{ng:type-I-counts}
 \Delta w=1,\qquad \Delta D=1,\qquad \Delta d=0.
\end{equation}
Equation~\eqref{ng:defect} now gives $\Delta B=0$.
An arbitrary front zigzag is not a type-I move.
\end{proof}

\begin{lemma}[Front type II]\label{ng:front-II}
The front type-II moves preserve $D,w,d$, and hence $B$.
\status{proved (refereed: bench A, 2026-09-05T15:51:04Z; transcribed from RC ng\_front (C029))}
\end{lemma}
\begin{proof}
The words $l_{m-1}\sigma_m\sigma_{m-1}$ and
$l_{m+1}\sigma_m\sigma_{m+1}$ replace $l_m$. In the first, the
through-strand is under at both crossings; in the second, it is over
at both. The other crossing branches are the oppositely directed arms
of the same cusp. Their signs are opposite for every permitted
through-strand and cusp orientation. After rounding the cusp, the arcs
bound an empty ordinary bigon with one common over-strand: an oriented
Reidemeister-II site. The same upper and lower cusp arms remain after
deletion. Thus $\Delta w=\Delta D=\Delta d=0$.
For right-cusp versions, follow these same local strands backwards;
no global reversal identity for $P$ is assumed.
\end{proof}

\begin{lemma}[Front type III]\label{ng:front-III}
The front type-III moves preserve $D,w,d$, and hence $B$.
\status{proved (refereed: bench A, 2026-09-05T15:51:04Z; transcribed from RC ng\_front (C029))}
\end{lemma}
\begin{proof}
The words $\sigma_{m+1}\sigma_m\sigma_{m+1}$ and
$\sigma_m\sigma_{m+1}\sigma_m$ form an actual ordinary
Reidemeister-III configuration. Label the three physical strands by
their initial top-to-bottom order. The three over/under choices give
one strict height order, not a cyclic order. Each physical pair crosses
on both sides with the same over/under bit and transported arrows.
An ordinary oriented Reidemeister-III move therefore applies, with
unchanged $D,w,d$ and $B$. In a reflected local template the strict
height order reverses, and an ordinary Reidemeister-III move still
applies. Checking this local template does not identify an entire
diagram with its reflection.
\end{proof}

\begin{lemma}[Zigzag and crossed-cusp deletion]
\label{ng:deletions}
Deleting an empty zigzag lowers $s$ by two and cannot increase $B$; applying
the crossed-cusp shortcut lowers $s$ by one and cannot increase $B$.
\status{proved (refereed: bench A, 2026-09-05T23:45:25Z and 2026-09-06T01:58:43Z; transcribed from RC ng\_front (C029); round~5: the $s$ consequence of both branches carried into the statement ($\Delta s=-2$ for the zigzag deletion, $-1$ for the crossed-cusp shortcut) and printed at the two displays of the proof, with the citation of Ng at its point of use, F-25-142)}
\end{lemma}
\begin{proof}
A zigzag consists of two consecutive cusps with no crossing in an empty
rectangle. Their directions are both downward or both upward: the
strand reverses its $x$ direction twice while continuing its vertical
passage. After rounding, it is a simple ordinary arc, positively page
isotopic relative to its endpoints to the straightened arc. Deletion
therefore gives
\begin{equation}\label{ng:zigzag-counts}
 \Delta w=\Delta d=0,\qquad \Delta D\in\{0,-2\}.
\end{equation}
Thus it cannot increase $B$; the two cusps are removed and no crossing is
touched, so $s$ decreases by two. This covers either oriented destabilization
without a knot-type assertion.

The word procedure also encounters $l_i\sigma_i$ or $\sigma_i r_i$,
a cusp whose own arms cross once. Replace it by the uncrossed cusp with
the same two oriented boundary attachments. After rounding, ordinary
Reidemeister I deletes its empty monogon. The arms are oppositely
directed, so the old crossing has sign $-1$. Deletion raises $w$ by one
and flips the cusp direction because the boundary arms exchange places.
If the old cusp was downward, $\Delta D=-1$; otherwise $\Delta D=1$.
Consequently
\begin{equation}\label{ng:crossed-cusp-counts}
 \Delta d=0,\qquad \Delta w=1,\qquad \Delta B\in\{-2,0\}.
\end{equation}
One crossing is removed and the cusp count is unchanged, so $s$ decreases
by one; this is the compressed dashed-arrow case of Ng's Figure~1
\cite{Ng}. It is not asserted to be a Legendrian isotopy. The direct
ordinary certificate avoids inserting a temporary higher-complexity
front merely to recognize its stabilization.
\end{proof}

\begin{lemma}[Standard front circles]\label{ng:circle}
Deleting a separated standard front circle with nonempty remainder
preserves $B$. A single standard front circle, and any union of such
circles, has $B=0$.
\status{proved (refereed: bench A, 2026-09-05T15:51:04Z; transcribed from RC ng\_front (C029))}
\end{lemma}
\begin{proof}
A simple crossing-free component with exactly one left and one right
cusp has $D=1$, $w=0$ and ordinary value one, in either orientation.
A standard circle separated by the word procedure has no mixed
crossings. Its smoothed component may lie in a bounded face of the
remainder; Lemma~\ref{lp:split-circle} permits this nesting.

For nonempty remainder, $P_{\rm before}=\delta P_{\rm after}$.
The leading $a$ coefficient of $\delta$ is $z^{-1}$, in degree one.
Its product with the leading coefficient of $P_{\rm after}$ is nonzero
because the coefficient ring is an integral domain. Therefore
\begin{equation}\label{ng:circle-counts}
 d_{\rm before}=d_{\rm after}+1,\qquad
 D_{\rm before}=D_{\rm after}+1.
\end{equation}
The writhe is unchanged, so $B_{\rm before}=B_{\rm after}$.
If this is the last component, stop instead: $D=1$ and $w=d=0$ give
$B=0$. Repeated deletion shows that a union of standard front circles
also has $B=0$. No empty-link polynomial or degree is used.
\end{proof}

\subsubsection{The two oriented cusp-skein branches}
\label{ng:cusp-skein-section}

\begin{lemma}[Cusp-skein inequality]\label{ng:cusp-skein}
For either principal direction of an oriented cusp-skein interchange,
$B$ of the earlier branch is at least the minimum of $B$ of the other
principal branch and $B$ of the unique compatible smoothing. The
smoothing has one fewer singularity; the principal branches have the
same singularity count.
\status{proved (refereed: bench A, 2026-09-05T15:51:04Z; transcribed from RC ng\_front (C029))}
\end{lemma}
\begin{proof}
Omitting spectator strands, label one left endpoint $L$ and three
right endpoints $1,2,3$ from top to bottom. Put
\begin{equation}\label{ng:cusp-words}
 A=l_2\sigma_1,\qquad A'=l_1\sigma_2,\qquad
 C_{\rm top}=l_1,\qquad C_{\rm bottom}=l_2.
\end{equation}
Both $A$ and $A'$ pair $L$ with $2$ and $1$ with $3$. The $L$--$2$
strand is the through-strand; the other contains the cusp. In $A$ the
through-strand is over, whereas in $A'$ the cusp strand is over.
Each strand encounters the one crossing once. Exchanging over and under
at that crossing makes their local ordered crossing records identical,
including the boundary attachments. After gluing any same actual
exterior, the full named records are identical. Thus their polynomial
values are the two crossing choices at one ordinary skein site. The
move before that over/under exchange is not asserted to be page isotopy.

Let $t=1$ for the orientation from $L$ to $2$, and $t=-1$ otherwise.
Let $u=1$ for the orientation from $1$ to $3$, a downward cusp, and
$u=-1$ otherwise. In $A$, the cusp's upper arm has $x$ direction $-u$;
in $A'$, the cusp's lower arm has $x$ direction $u$. The crossing
determinant rule gives signs $-tu$ and $tu$, respectively:
\begin{center}
\begin{tabular}{@{}ccccc@{}}
$(t,u)$ & $\operatorname{sign}(A)$ & $\operatorname{sign}(A')$
 & Compatible smoothing & Local $D$ in $A,A',C$\\
\hline
$(+,+)$ & $-1$ & $+1$ & $C_{\rm top}$    & $1$\\
$(+,-)$ & $+1$ & $-1$ & $C_{\rm bottom}$ & $0$\\
$(-,+)$ & $+1$ & $-1$ & $C_{\rm bottom}$ & $1$\\
$(-,-)$ & $-1$ & $+1$ & $C_{\rm top}$    & $0$
\end{tabular}
\end{center}
Oriented smoothing joins each incoming strand to the other outgoing
strand. If $t=u$, it pairs $1$ with $2$ and $L$ with $3$: the top cusp
and through-strand. If $t=-u$, it pairs $L$ with $1$ and $2$ with $3$:
the bottom cusp and through-strand. The new cusp direction is $u$ in
either case, from $1$ to $2$ or from $2$ to $3$. This proves the last
column without any hypothesis on how the exterior joins the two strands.
Smoothing may split or join parameter components, but cannot create an
empty diagram. Only the one compatible smoothing is used, never both
unoriented smoothings.

Let $E_+,E_-,E_0$ be these actual ordinary skein diagrams. Their writhes
are $w_0+1,w_0-1,w_0$, where $w_0$ is the common exterior writhe.
Solving the skein relation separately in the two directions gives
\begin{equation}\label{ng:skein-plus}
 P_+=a^{-2}P_-+a^{-1}zP_0
\end{equation}
and
\begin{equation}\label{ng:skein-minus}
 P_-=a^2P_+-azP_0.
\end{equation}
The degree of a nonzero sum is at most the maximum of the summand
degrees. Multiplication by $a^k$ shifts degree by $k$, and multiplication
by $z$ does not change it. Equation~\eqref{ng:skein-plus} yields
$d_+\leq\max\{d_--2,d_0-1\}$. Adding $w_0+2$ gives
\begin{equation}\label{ng:degree-plus}
 d_++1+(w_0+1)
 \leq\max\{d_-+1+(w_0-1),\ d_0+1+w_0\}.
\end{equation}
Equation~\eqref{ng:skein-minus} instead yields
$d_-\leq\max\{d_++2,d_0+1\}$. Adding $w_0$ gives
\begin{equation}\label{ng:degree-minus}
 d_-+1+(w_0-1)
 \leq\max\{d_++1+(w_0+1),\ d_0+1+w_0\}.
\end{equation}
The table and unchanged exterior give the same downward-cusp count
in all three diagrams. Subtracting~\eqref{ng:degree-plus} and,
respectively,~\eqref{ng:degree-minus} from this common number proves
\begin{equation}\label{ng:skein-defect}
 \begin{aligned}
 B(A)&\geq\min\{B(A'),B(C)\},\\
 B(A')&\geq\min\{B(A),B(C)\}.
 \end{aligned}
\end{equation}
Here $C$ is the compatible smoothing. It has one fewer crossing and
the same number of cusps, so $s(C)=s(A)-1=s(A')-1$; the principal
crossing-change branch preserves $s$. Reflected and right-cusp templates
follow by relabeling their actual attachments and arrows in this local
calculation, not from a global reflection equality for $P$.
\end{proof}

\subsubsection{The finite geometric reduction}
\label{ng:word-section}

\begin{literature}[Finite front-word reduction]\label{ng:finite-word}
For an actual finite front word on the domain of
Definition~\ref{ng:front-domain}, the procedure of Ng
\cite[pp.~6--8 and Figure~1]{Ng} and the elementary-word proof of
Rutherford \cite[Lemma~3.2, pp.~8--12 of the arXiv version]{Rutherford}
supplies a finite principal chain using:
\begin{enumerate}
\item typed disjoint-gadget commutations and the three local front moves
  in Lemmas~\ref{ng:front-I}, \ref{ng:front-II} and~\ref{ng:front-III};
\item the cusp-skein crossing interchange~\eqref{ng:cusp-words}, or its
  reflected pattern;
\item zigzag deletion, the crossed-cusp shortcut of
  Lemma~\ref{ng:deletions}, or deletion of a separated standard front
  circle.
\end{enumerate}
Stop at the first strict decrease of $s$ or at the standard-circle base.
Before that decrease the selected procedure does not increase $s$.
Secondary progress is measured at the boundaries of complete finite
blocks of the word procedure: a completed nonterminal block moves a
singularity to the left of the selected rightmost left cusp, and hence
decreases the number on its right. This is not a strict decrease at
every individual elementary move. Arm-string extension within a block
is bounded by the active strand count. Thus every nonbase front has a
finite principal chain to a front of smaller $s$. At each cusp-skein
interchange, the compatible smoothing branch has smaller $s$ than the
front at that stage.

This is the constructive word procedure, not a quantification over
arbitrary Legendrian isotopies. The imported content is the source's finite
descent, not its printed sentence: Ng's Lemma~1 gives one step, iterated
here until $s$ decreases; Rutherford's two ruling-polynomial terminal
branches are replaced here by the geometric shortcuts of
Lemma~\ref{ng:deletions}; and the finiteness of arm-string extension, which
Rutherford asserts, is given its reason here (seat L-3/L-4 on SM2, R6;
F-25-83).
\status{lit: finite Ng--Rutherford word procedure; Registry~\ref{reg:slbound}; round~3: the scope of the import stated, F-25-83}
\end{literature}

For precision, the finite cases and index transcription used in this
source input are as follows. Write the word at its rightmost left cusp
as $Xl_mY$, where $Y$ contains only crossings and right cusps. The two
arm strings are
$(\sigma_{m-1}\cdots\sigma_{m-N_1})
 (\sigma_{m+1}\cdots\sigma_{m+N_2})$.
They use disjoint crossing positions and commute. The complete cases
are listed below; each row includes both arm versions where applicable.

\begingroup
\small
\setlength{\tabcolsep}{4pt}
\begin{longtable}{@{}>{\raggedright\arraybackslash}p{.30\linewidth}>{\raggedright\arraybackslash}p{.43\linewidth}>{\raggedright\arraybackslash}p{.21\linewidth}@{}}
\textbf{First letter of $Y$} & \textbf{Finite procedure}
 & \textbf{Progress or terminal}\\
\hline
\endfirsthead
\textbf{First letter of $Y$} & \textbf{Finite procedure}
 & \textbf{Progress or terminal}\\
\hline
\endhead
Crossing outside the strings and their next neighbors
 & Commute into $X$
 & Fewer right-hand singularities\\[3pt]
Crossing at a next neighboring position
 & Extend that arm string
 & $N_1$ or $N_2$ increases, bounded by active strand count\\[3pt]
Crossing at an occupied outer string endpoint
 & Move the cusp along its string using skein interchanges
 & Expose type II and remove two crossings\\[3pt]
Crossing strictly inside a string
 & Type III, then commutations
 & Move the crossing into $X$\\[3pt]
Middle crossing, exactly one nonempty string
 & Expose type II
 & Remove two crossings\\[3pt]
Middle crossing, no strings
 & Apply the $l_i\sigma_i$ shortcut
 & Lower $s$ by Lemma~\ref{ng:deletions}\\[3pt]
Middle crossing, both strings nonempty
 & Cusp interchange, type III and commutations
 & Move the crossing into $X$\\[3pt]
Right cusp beyond an arm's next neighbor
 & Commute into $X$ with the correct two-strand index shift
 & Fewer right-hand singularities and active strands\\[3pt]
Right cusp at a next neighbor
 & Move the selected cusp along the arm string
 & Expose a zigzag, then delete it\\[3pt]
Right cusp at a string endpoint
 & Expose $\sigma_i r_i$
 & Crossed-cusp shortcut\\[3pt]
Right cusp strictly inside a string
 & Commute to expose type II
 & Remove two crossings\\[3pt]
Middle right cusp, neither string nonempty
 & Standard circle without mixed crossings
 & Delete, or stop at the last component\\[3pt]
Middle right cusp, exactly one string nonempty
 & Type I
 & Remove one crossing and two cusps\\[3pt]
Middle right cusp, both strings nonempty
 & Cusp interchange, then commuting
 & Expose type II and remove two crossings
\end{longtable}
\endgroup

Right-cusp absorption lowers active strand count by two; crossing
absorption lowers the right-hand crossing count. String extension is
bounded by active strand count and cannot continue without eventually
reaching another row. These are the source's nested finite descents.
Only deletion directions of types I and II occur here; type III,
commutations and cusp-skein changes preserve $s$. Stabilized cusp cases
use the direct shortcut of Lemma~\ref{ng:deletions}
(display~\eqref{ng:crossed-cusp-counts}; $\Delta s=-1$), so no temporary
increase of $s$ is needed.

Two impossible printed ranges $m-N_1>i>m$ on Rutherford's pp.~11--12
are transcribed here as the interior range $m-N_1<i<m$, as the
neighboring cases and actual strands require. The duplicated $m+2$
in the increasing string on p.~11 continues with $m+3$, with the
corresponding two-strand shift on p.~12. The malformed literal strings
are not imported as identities: the actual strand operations and Ng's
Figure~1 specify the moves used here. Ng cites the published numbering
Lemma~3.3; the retained Rutherford arXiv version numbers this proof
Lemma~3.2.

No Rutherford ruling or Kauffman value, maximizing representative, or
earlier general characterization of Legendrian equivalence is consumed.
Literature input~\ref{ng:finite-word} is precisely the local constructive
portion of those proofs. It is not an application of the literal
statement of Ng's Theorem~1 to an ambient link invariant. The local
polynomial/count certificates above and the following degree induction
are separate arguments.

\subsubsection{The degree induction}
\label{ng:induction-section}

\begin{theorem}[Individual-front polynomial bound]
\label{ng:local-front-bound}
For every front $F$ on the domain of Definition~\ref{ng:front-domain},
with the same polynomial evaluated on its actual ordinary cusp rounding,
\begin{equation}\label{ng:front-inequality}
 w(F)-D(F)\leq-\deg_a P_{S(F)}-1.
\end{equation}
\status{proved (refereed: bench A, 2026-09-05T15:51:04Z, 2026-09-06T00:12:27Z and 2026-09-06T01:50:42Z; transcribed from RC ng:local-front-bound (C029); consumes Literature input~\ref{ng:finite-word})}
\end{theorem}
\begin{proof}
Use strong induction on $s(F)$; the quantity $d(F)=\deg_aP_{S(F)}$ of
display~\eqref{ng:defect} does not depend on the rounding $S(F)$, by
Lemma~\ref{ng:smoothing-record}. Every nonempty component has at least
one left and one right cusp. Thus the base $s=2$ is a standard circle,
and Lemma~\ref{ng:circle} gives $B=0$. Any union of standard front
circles encountered as a terminal base also has $B=0$. An empty object
is never evaluated.

Suppose the assertion holds for all fronts with smaller $s$ than $F$.
Take the finite principal chain of Literature input~\ref{ng:finite-word},
stopped at its first strict decrease or at a standard-circle base.
Every preterminal principal front has $s$ at most $s(F)$. At each
cusp-skein step, its compatible smoothing therefore has strictly smaller
$s$ than $F$, and strong induction gives nonnegative $B$ for it. The
final principal front has nonnegative $B$ by the same induction or by
its explicit standard-circle base.

Traverse the chain backwards. Commutations and local front moves
preserve $B$ by Lemmas~\ref{ng:commutation}, \ref{ng:front-I}, \ref{ng:front-II}
and~\ref{ng:front-III}.
A deletion cannot increase $B$ in the forward direction, by
Lemmas~\ref{ng:deletions} and~\ref{ng:circle}; nonnegativity after
deletion therefore implies nonnegativity before it. At a cusp-skein
step, the later principal branch and its smaller smoothing branch both
have nonnegative $B$. The appropriate inequality
of Lemma~\ref{ng:cusp-skein}, display~\eqref{ng:skein-defect}, gives nonnegative $B$ for the earlier
principal branch. These are all possible operations. Finite backward
propagation yields $B(F)\geq0$, completing strong induction.
Substituting~\eqref{ng:defect} gives~\eqref{ng:front-inequality}.

Nonvanishing is used before every degree, including both skein branches
and the split product. Cancellation in sums was allowed. The proof
covers arbitrary positive component counts and all orientations.
\end{proof}

\begin{remark}[Boundary of the front bound]\label{ng:consumer-boundary}
Theorem~\ref{ng:local-front-bound} is only a polynomial inequality for
one actual front. Etnyre's separate local formula identifies
$w(F)-D(F)$ with the self-linking number of its chosen positive
transverse pushoff. Identifying $P_{S(F_T)}$ with the value on a
separately supplied transverse diagram requires a separate endpoint
comparison. The contact composition below supplies it using the
explicitly retained global HOMFLY source premise, after constructing
the actual ordinary rounding and ambient family. That premise is not
part of this individual-front argument; it is
Literature input~\ref{lit:homfly}, consumed as the registry records.
\end{remark}

\subsection{Contact conventions and a bound for each representative}

All Legendrian and transverse knots and isotopies in this subsection are
smooth embedded curves and smooth families in standard contact space.
The neighbourhood, generic-front, ambient-isotopy and linking/framing
constructions below are internal proofs, using smooth calculus, the inverse
function theorem and smooth ODE existence, uniqueness, continuation and
dependence.

\par\noindent\emph{Published proof comparison.}
The neighbourhood construction uses the relative Moser method of
Geiges~\cite{GeigesContact}, Theorem~2.32 and Example~2.33, with
the stability proof of Theorem~2.27 (statement p.~21, proof pp.~22--25 of
the retained arXiv version).
His form $dZ+X\,dY$ pulls back to $dz-y\,dx$ under
$(X,Y,Z)=(-y,x,z)$; this coordinate change has determinant $+1$.
His Hamiltonian proof, Theorem~2.22 and Corollary~2.23,
pp.~16--17, allows real-valued Hamiltonians; the $\RR^+$ in
Theorem~2.22's statement is inconsistent with its proof.
Here compact support supplies global flows on $\RR^3$.

Geiges's Lemma~2.45, pp.~38--39, gives an open-interval
front perturbation and the local cusp model. Its instruction to
perturb one coordinate and integrate the Legendrian equation does
not by itself preserve the period of a closed curve. The closed
generic-front construction here instead uses ambient contact flows
and carries the chosen pushoff. The linking integrand agrees with
the Euclidean formula in DeTurck--Gluck~\cite{DeTurckGluckLinking},
p.~3; the crossing-count and uniform-framing arguments are printed
here separately. These comparisons retain the distinction between
the published inputs and the constructions printed here, which await refereeing.

\begin{lemma}[A transverse neighbourhood and its Legendrian pushoff]
\label{fd:transverse-neighborhood}
Let $\alpha=dz-y\,dx$ and let
$T:\RR/(2\pi\mathbb Z)\to\RR^3$ be a smooth embedded oriented circle
with $\alpha(T')>0$. There are $\delta>0$ and a smooth embedding
$H:S^1\times D_\delta\to\RR^3$ fixing the parametrized core $T$ such that
\begin{equation}\label{fd:transverse-model}
 H^*\alpha=h\alpha_0,\qquad
 \alpha_0=d\theta+u\,dv-v\,du,\qquad h>0.
\end{equation}
There is an oriented Legendrian knot $L$ in this neighbourhood whose
positive transverse pushoff is transversely isotopic to $T$. An explicit
compactly supported ordinary ambient isotopy carries the parametrized
$L$ to the parametrized $T$, preserving their orientations.
\status{proved (refereed: bench A, 2026-09-05T15:51:04Z; transcribed from RC d2d\_contact\_constructions (C030))}
\end{lemma}
\begin{proof}
Write $T=(x,y,z)$, $a=\alpha(T')>0$, and $c=\sqrt{2a}$.
The global frame $E_1=\partial_x+y\partial_z$, $E_2=\partial_y$ spans
$\ker\alpha$, and $d\alpha(E_1,E_2)=1$. Set
\begin{equation}
 F(\theta,u,v)=T(\theta)+c(\theta)
       \bigl(uE_1(T(\theta))+vE_2(T(\theta))\bigr).
\end{equation}
At the core its derivative has independent columns $T',cE_1,cE_2$:
applying $\alpha$ to a relation first kills the $T'$ coefficient.
The inverse function theorem and compactness give a common radius on
which every derivative is invertible. This radius can be decreased so
that $F$ is injective: otherwise two distinct points with disc radii
tending to zero have equal images; limiting circle parameters coincide
because $T$ is embedded, placing both points in one local inverse chart,
a contradiction. Thus $F$ embeds a product neighbourhood.

Put $\beta=F^*\alpha/a(\theta)$. At the core $\beta=d\theta$ and
\begin{equation}
 d\beta(\partial_u,\partial_v)
 =a^{-1}d\alpha(cE_1,cE_2)=c^2/a=2.
\end{equation}
The derivative of $a^{-1}$ vanishes on this vertical pair. Consequently
$\alpha_t=(1-t)\alpha_0+t\beta$ satisfies
$(\alpha_t\wedge d\alpha_t)(\partial_\theta,\partial_u,\partial_v)=2$
on the core. Compactness of $[0,1]\times S^1$ supplies a common smaller
neighbourhood where every $\alpha_t$ is contact and
$\alpha_t(\partial_\theta)>0$.

Let $\nu=\beta-\alpha_0$ and define
\begin{align}
 X_1&=\partial_u-
   \frac{\alpha_t(\partial_u)}{\alpha_t(\partial_\theta)}\partial_\theta,
 &X_2&=\partial_v-
   \frac{\alpha_t(\partial_v)}{\alpha_t(\partial_\theta)}\partial_\theta,
 \\
 D_t&=d\alpha_t(X_1,X_2),
 &V_t&=\frac{\nu(X_1)X_2-\nu(X_2)X_1}{D_t}.
\end{align}
The $X_i$ form a basis of $\ker\alpha_t$; contactness makes $D_t\ne0$.
Evaluation on each $X_i$ gives
\begin{equation}
(\iota_{V_t}d\alpha_t)|_{\ker\alpha_t}=-\nu|_{\ker\alpha_t}
\end{equation}
The vector field $V_t$ vanishes on the core, since $\nu$ does.
Choose $\rho>0$ so that $S^1\times\overline D_\rho$ lies strictly
inside the common neighbourhood where $V_t$ is smooth for every
$t\in[0,1]$. On this closed product its norm is at most
$C\sqrt{u^2+v^2}$ uniformly in $t,\theta$. Along a trajectory,
$(u^2+v^2)'\leq2C(u^2+v^2)$, hence
$r(t)\leq e^{Ct}r(0)$. Choosing the initial disc radius below
$\rho e^{-C}/2$ keeps the flow in a compact inner product neighbourhood
for $0\leq t\leq1$. Smooth ODE existence, continuation and uniqueness
give a flow $\Phi_t$, fixing the core and invertible onto its image.

The covector $\nu+\iota_{V_t}d\alpha_t$ annihilates $\ker\alpha_t$,
so it is $\mu_t\alpha_t$ for a smooth scalar $\mu_t$.
Since $\alpha_t(V_t)=0$, Cartan's differentiation identity gives
\begin{align}
 \frac{d}{dt}\Phi_t^*\alpha_t
 &=\Phi_t^*(\nu+\mathcal L_{V_t}\alpha_t)\notag\\
 &= (\mu_t\circ\Phi_t)\Phi_t^*\alpha_t.
\end{align}
Solving this scalar equation gives
$\Phi_1^*\beta=\exp(\int_0^1\mu_t\circ\Phi_t\,dt)\alpha_0$.
For $H=F\circ\Phi_1$, the additional factor
$a\circ\operatorname{pr}_\theta\circ\Phi_1$ is also positive, proving
\eqref{fd:transverse-model}, including coorientation and core orientation.

Choose an integer $N>0$ with $b=N^{-1/2}<\delta$. The embedded circle
\begin{equation}
 L_0(\theta)=(\theta,b\cos N\theta,-b\sin N\theta)
\end{equation}
has $\alpha_0(L_0')=1-Nb^2=0$; its degree-one first coordinate makes it
an embedding. Hence $L=H\circ L_0$ is Legendrian. Fix $\kappa\ne0$ and
a small $\epsilon>0$ with $b+\epsilon<\delta$, and put
\begin{equation}
 B(\theta,s)=
 (\theta+\kappa s,(b-s)\cos N\theta,-(b-s)\sin N\theta),
 \qquad -\epsilon<s\leq b.
\end{equation}
Radius $b-s$ and the first coordinate prove injectivity. The radial
derivative proves immersion for $s<b$; at $s=b$ it remains independent
of $\partial_\theta$, and the Cartesian formula is smooth there.
Direct differentiation yields
\begin{equation}\label{fd:pushoff-annulus}
 B^*\alpha_0=\bigl(1-N(b-s)^2\bigr)d\theta+\kappa\,ds.
\end{equation}
Thus the annulus is transverse to the contact planes, even along $s=0$.
Its circles are negative transverse for $s<0$, Legendrian at $s=0$,
and positive transverse for $0<s\leq b$. A sufficiently small positive
circle is therefore the positive pushoff in the source convention.
The whole family from that circle to $s=b$ is a transverse isotopy to
$T(\theta+\kappa b)$, which is the same oriented transverse knot.
Finally $H\circ B$ on $S^1\times[0,b]$ is an embedded oriented annulus
with boundary $L-T$. To prove the stronger ambient-isotopy conclusion
needed by the polynomial, we give such an isotopy explicitly.
Write $v(\theta)=b(\cos N\theta,-\sin N\theta)$, choose
$\epsilon_0=(\delta-b)/4$, and put
$R_1=b+\epsilon_0$, $R_2=b+2\epsilon_0$.
Thus $b<R_1<R_2<\delta$. For
$\eta(s)=e^{-1/s}$ when $s>0$ and $\eta(s)=0$ when $s\leq0$, set
\begin{equation}
 \chi(s)=\frac{\eta(R_2^2-s)}
 {\eta(R_2^2-s)+\eta(s-R_1^2)}.
\end{equation}
Every positive-side derivative of $\eta$ is $e^{-1/s}$ times a
polynomial in $1/s$, tending to zero at zero; hence $\eta$ is smooth.
The denominator is positive because $R_1<R_2$. Thus $\chi$ is smooth,
equals one for $s\leq R_1^2$ and zero for $s\geq R_2^2$.
The smooth product vector field
\begin{equation}
 V(\theta,w)=\bigl(0,-\chi(|w|^2)v(\theta)\bigr)
\end{equation}
is well defined on the circle since $N$ is integral. Squared radius
makes it smooth at $w=0$; its support lies in
$S^1\times\overline D_{R_2}$. The embedding $H$ has an open image
$U$ and a smooth inverse there by the inverse function theorem.
Push $V$ forward by $H$ on $U$ and extend by zero outside $U$.
This field $X$ is globally smooth: its support lies in the compact
set $K=H(S^1\times\overline D_{R_2})\subset U$, and it vanishes
on a neighbourhood of every point outside $K$.

The field $X$ is bounded, so a trajectory moves a distance at most
$M|t|$ on any finite time interval. Smooth ODE continuation gives
its global flow $\Psi_t$ for every real $t$. Uniqueness gives the
smooth inverse $\Psi_{-t}$; thus these are ambient diffeomorphisms.
They fix every point outside $K$. Their derivative determinants
start positive and never vanish, so they preserve ambient orientation.
Extending them by the identity at infinity gives an ambient isotopy
of the oriented three-sphere, with the same compact support.

For $0\leq t\leq1$, the path $(\theta,(1-t)v(\theta))$ stays
at radius at most $b<R_1$, where $\chi=1$, and its time derivative
is exactly $V$. The chain rule and ODE uniqueness therefore give
\begin{equation}\label{fd:helix-ambient-isotopy}
 \Psi_t(L(\theta))=H(\theta,(1-t)v(\theta)),\qquad
 \Psi_1(L(\theta))=T(\theta).
\end{equation}
The parameter $\theta$ is unchanged, so the endpoint circle orientation
is exactly that of $T$. The core need not be fixed: its own track is
$\Psi_t(T(\theta))=H(\theta,-t v(\theta))$ on this interval.
There is thus no collapse of two distinct curves by a diffeomorphism.
This is an ordinary ambient isotopy, not a contact isotopy. Indeed
the model form on its tracked knot has value
$1-Nb^2(1-t)^2=2t-t^2>0$ for $0<t\leq1$, whereas the initial
knot is Legendrian. The preceding twisted annulus, not this ambient
extension, remains the positive-pushoff and transverse-isotopy producer.
\end{proof}

\begin{lemma}[Compact parameter avoidance]\label{fd:parameter-avoidance}
Let $K$ be a compact subset of a smooth $d$-dimensional coordinate
manifold and $B\subset\RR^m$ a closed parameter ball. Suppose that
$F$ is smooth on a neighbourhood of $K\times B$, takes values in
$\RR^q$, and has derivative of rank $q>d$ at every zero in that product.
The parameters $a$ for which $F(t,a)=0$ for some $t\in K$ form a closed
set with empty interior. The same assertion holds simultaneously for
a finite collection of such maps.
\status{proved (refereed: bench A, 2026-09-05T15:51:04Z; transcribed from RC d2d\_contact\_constructions (C030))}
\end{lemma}
\begin{proof}
The zero set is compact, so its parameter projection is closed.
At each zero a nonzero $q$-minor and the inverse function theorem express
the local zero set as a smooth graph of $e=d+m-q<m$ variables.
Finitely many such charts restricted to compact cubes cover the zero
set; their parameter projections are Lipschitz maps of those cubes.
Subdividing an $e$-cube into $O(n^e)$ pieces of side $O(n^{-1})$ covers
its projected image by $m$-boxes of total volume $O(n^{e-m})$.
This tends to zero. The image cannot contain a fixed open $m$-box,
because any finite box cover has total volume at least the covered
volume: partition at the finitely many box faces to prove this inequality.
A finite union of the chart images has the same estimate. This proves
empty interior and the finite-collection assertion without Sard's theorem.
\end{proof}

\begin{lemma}[Explicit contact motions]\label{fd:contact-motions}
For a smooth compactly supported $H:\RR^3\to\RR$, the vector field
\begin{equation}
 X_H=-H_y\partial_x+(H_x+yH_z)\partial_y+(H-yH_y)\partial_z
\end{equation}
has a global flow $\phi_H^s$ of coorientation-preserving contact
diffeomorphisms for $\alpha=dz-y\,dx$: $(\phi_H^s)^*\alpha=c^H_s\,\alpha$
with a smooth positive function $c^H_s$, its \emph{conformal factor}. Finite compositions of their
small-time flows depend smoothly on the times and are arbitrarily
$C^k$-close to the identity on compact sets for each fixed finite $k$.
\status{proved (refereed: bench A, 2026-09-05T15:51:04Z and 2026-09-06T01:50:42Z; transcribed from RC d2d\_contact\_constructions (C030))}
\end{lemma}
\begin{proof}
Substitution gives $\alpha(X_H)=H$ and
$\iota_{X_H}d\alpha=H_z\alpha-dH$; therefore
$\mathcal L_{X_H}\alpha=H_z\alpha$. Differentiating the pullback and
solving the resulting scalar ODE gives
\begin{equation}
 (\phi_H^s)^*\alpha
 =\exp\left(\int_0^s H_z\circ\phi_H^v\,dv\right)\alpha.
\end{equation}
This exponential is the conformal factor $c^H_s$; it is positive. The vector field is smooth, compactly supported
and bounded; trajectories cannot escape to infinity in finite time.
ODE continuation gives all finite times, and backwards uniqueness gives
the inverse. Smooth ODE dependence proves the last assertion.
\end{proof}

\begin{theorem}[A generic front preserving the chosen positive pushoff]
\label{fd:generic-front}
Every smooth oriented Legendrian embedding $L:S^1\to\RR^3$ admits
a smooth ambient coorientation-preserving contact isotopy to a
Legendrian embedding $L_g$ whose $xz$ front has only finitely many
semicubical cusps and transverse double points, with no triple point
and no cusp on another branch. At each cusp there is a smooth local
coordinate $u=y-y_0$ in which
\begin{equation}\label{fd:generic-exact-germ}
 x=x_0+Au^2,\qquad y=y_0+u,\qquad
 z=z_0+Ay_0u^2+\tfrac23 Au^3,\qquad A\ne0;
\end{equation}
the coordinate $u$ may increase or decrease along the prescribed
orientation. The isotopy carries a chosen thin
positive-pushoff annulus and preserves its positive-pushoff transverse
isotopy class, as well as the oriented topological knot type.
\status{proved (refereed: bench A, 2026-09-06T00:35:51Z and 2026-09-06T01:58:43Z; transcribed from RC d2d\_contact\_constructions (C030); round~3: the exact cusp germ its last proof step establishes exported to the conclusion, F-25-86 (Codex R-P5, P-25-11); round~5: the conformal factor of the concatenated contact isotopy named $c_s$ and derived from Lemma~\ref{fd:contact-motions}, the cusp-germ function renamed $\psi$ and the second-derivative matrix $M$, a departure from RC's letters, F-25-154)}
\end{theorem}
\begin{proof}
Write $L=(x,y,z)$, so $z'=yx'$. For finitely many compactly supported
Hamiltonians, let $\Phi_a$ be the composition of their time-$a_i$ flows.
At $a=0$, the $a_i$ derivative of $\Phi_a\circ L$ is $X_{H_i}\circ L$.
The path $s\mapsto\Phi_{sa}$ is an ambient contact isotopy by
Lemma~\ref{fd:contact-motions}; its images are exactly Legendrian
embeddings, without an approximation or an area-closing correction.

First remove simultaneous zeros of $(x',x'')$. At such a point
$\theta_0$, immersion and $z'=yx'=0$ imply $y'(\theta_0)\ne0$.
With $y_0=y(\theta_0)$, choose Hamiltonians equal near $L(\theta_0)$ to
\begin{equation}
 H_2=-\tfrac12(y-y_0)^2,\qquad H_3=-\tfrac16(y-y_0)^3,
\end{equation}
and multiply them by a smooth bump equal to one there. Their
infinitesimal $x$ changes are $y-y_0$ and $(y-y_0)^2/2$.
Thus their derivative columns for $(x',x'')$ at $\theta_0$ are
$(y',y'')^T$ and $(0,(y')^2)^T$, with determinant $(y')^3\ne0$.
Finitely many such minor neighbourhoods cover the initial compact zero
set. On their compact complement $|(x',x'')|$ has a positive minimum.
Continuity therefore supplies a small parameter ball in which every
possible zero retains one of the nonzero minors. Apply
Lemma~\ref{fd:parameter-avoidance} with $d=1,q=2$ to choose a small
parameter avoiding all these zeros. If the initial zero set is empty,
this step is unnecessary. Call the result $L_1$.
Every zero of $x_1'$ is simple, so compactness makes their number finite.
At each of them $y_1'\ne0$ by immersion and Legendrianity.

We next supply a uniform collar of the parameter diagonal. Cover $S^1$
by finitely many interiors of closed intervals of two kinds: regular
intervals where $x_1'$ has fixed nonzero sign, and critical intervals
where $x_1''$ and $y_1'$ have fixed nonzero signs and the endpoint signs
of $x_1'$ are opposite. On the first kind $x_1$ is strictly monotone.
On a critical interval there is one critical point $\theta_c$.
For two points on opposite sides with equal $x$ coordinate $X$,
the inverse branches and $dz=y\,dx$ give
\begin{equation}
 z_1(\theta_+)-z_1(\theta_-)
 =\int_{x_1(\theta_c)}^X
     \bigl[y_1(\theta_+(v))-y_1(\theta_-(v))\bigr]\,dv\ne0.
\end{equation}
The integrand has a fixed strict sign by monotonicity of $y_1$;
the integration interval has nonzero length. The formula also handles
a maximum by reversing the integration sign. The front is therefore
injective on each interval. Retaining the finitely many derivative and
endpoint-sign margins preserves this argument for every sufficiently
$C^2$-close Legendrian embedding. A Lebesgue number $\delta>0$ of the
finite interval cover gives uniform front injectivity for every pair
of distinct parameters at circular distance less than $\delta$.
This Lebesgue number follows from compactness: a sequence of pairs of
diameter tending to zero and contained in no cover interval would
converge to a point inside one of those open intervals.

Let $K_2$ be the compact ordered-pair set with distance at least $\delta$,
and $K_3$ the compact ordered-triple set with every pairwise distance at
least $\delta$. Write $p_a=(x_a,z_a)$. The bad configurations are zeros of
\begin{align}
 C(\theta,\eta,a)&=
  \bigl(x_a'(\theta),p_a(\theta)-p_a(\eta)\bigr)\in\RR^3,
 \\
 R(\theta,\eta,\tau,a)&=
  \bigl(p_a(\theta)-p_a(\eta),p_a(\theta)-p_a(\tau)\bigr)\in\RR^4.
\end{align}
At a zero of $C$ for $a=0$, use disjoint small balls about the two
distinct spatial points. The preceding $H_2$ at the first point
changes $x'(\theta)$ by $y_1'(\theta)\ne0$ but changes neither $x,z$
there. At the second point $q$, Hamiltonians equal locally to
$H^x=-(y-y_q)$ and $H^z=1$ have independent $(x,z)$ velocities
$(1,y_q)^T$ and $(0,1)^T$. Their supports avoid the first point.
Hence $C$ has parameter rank three. At a zero of $R$, choose disjoint
balls at the three distinct spatial points and use the same two
Hamiltonians at the second and third points only. The two independent
blocks give parameter rank four.

Choose finitely many such Hamiltonians covering both initial compact
zero sets. A zero-free initial map has a positive norm minimum and
needs no Hamiltonians. The nonzero minors persist near those zero
sets; the positive norm on their compact complements excludes new
zeros there for small parameters. Shrink the parameter ball also to
retain the preceding uniform local-injectivity margins. Apply
Lemma~\ref{fd:parameter-avoidance} simultaneously with $(d,q)=(2,3)$
and $(3,4)$, and choose a parameter outside both bad projections.
The resulting $L_2$ has no cusp on another branch and no triple point.

At any remaining double point both parameters are regular. Their
$y$ values differ, since otherwise all three spatial coordinates
coincide, contradicting embeddedness. The two front tangents have
determinant
\begin{equation}
 x'(\theta)x'(\eta)\bigl(y(\eta)-y(\theta)\bigr)\ne0.
\end{equation}
Thus every double point is transverse and isolated by the inverse
function theorem applied to $p(\theta)-p(\eta)$. Their ordered-pair
set is closed in compact $K_2$, so there are finitely many: an infinite
set would accumulate at another zero, contradicting isolation.
All properties so far persist under sufficiently small $C^2$
Legendrian perturbations. Indeed the local-injectivity margins persist,
and the now nowhere-zero maps $C,R,(x',x'')$ have positive norm minima
on their respective compact domains.

We finally make each cusp germ exactly semicubical, without a
singularity-normal-form assumption. At an $x$-critical point,
$y'\ne0$ allows $y$ as a local parameter. Write $x=f(y)$ and put
\begin{equation}
 u=y-y_0,\quad A=\tfrac12 f''(y_0)\ne0,\quad
 q(y)=f(y_0)+Au^2,\quad
 \psi(y)=\int_{y_0}^y(f(v)-q(v))\,dv.
\end{equation}
Here $\psi=O(u^4)$ and its first three derivatives have orders
$O(u^3),O(u^2),O(u)$. Where $H=\psi(y)$ the contact flow is exactly
\begin{equation}
 (x,y,z)\longmapsto
 (x-s\psi'(y),y,z+s(\psi(y)-y\psi'(y))).
\end{equation}
At time one, $x=q(y)$ and Legendrianity gives
\begin{equation}
 x=x_0+Au^2,\qquad z=z_0+Ay_0u^2+\tfrac23Au^3.
\end{equation}
The invertible affine page coordinates
$X=(x-x_0)/A$ and $Z=3(z-z_0-y_0(x-x_0))/(2A)$ give $(X,Z)=(u^2,u^3)$.

Multiply $\psi(y)$ by a smooth spatial cutoff equal to one on an inner
ball about the cusp and supported in a ball of radius $O(\epsilon)$
meeting no other portion of the embedded knot. Its derivatives of
order $k$ are $O(\epsilon^{-k})$, so the derivatives of $H$ through
order three are $O(\epsilon^{4-k})$. The vector field and its first
two derivatives are consequently $O(\epsilon)$ on the support.
The time-one flow is $C^2$-close to the identity: differentiating its
ODE gives $J'=DX_H J$ and
$M'=DX_H M+D^2X_H[J,J]$, with initial values $I,0$, and the elementary
exponential estimate bounds $J-I,M$ by $O(\epsilon)$.
An inner subarc has displacement $O(\epsilon^3)$ and a gap of order
$\epsilon$ to the boundary of the region where the cutoff equals one.
Its whole flow therefore uses the exact displayed formula.
Finitely many disjoint cusp balls and sufficiently small radii retain
all previous genericity margins, including those in the transition
regions. This gives the required exact cusp germs without new crossings.

Concatenate the constructed contact isotopies, using smooth time
reparametrizations flat at the joins. Each of them is a composition of
flows of Lemma~\ref{fd:contact-motions}, so $\Phi_s^*\alpha=c_s\,\alpha$ with
$c_s>0$ the product of their conformal factors. If $B$ is the chosen thin
pushoff annulus and $K$ its positive parallel circle, then
$\Phi_s\circ B$ remains a transverse annulus and
\begin{equation}
 \alpha((\Phi_s\circ K)')=(c_s\circ K)\,\alpha(K')>0.
\end{equation}
Its central circle is the Legendrian isotope and its positive parallel
circle is a positive pushoff thereof. Thus the chosen pushoff moves
through positive transverse embeddings. The ambient isotopy also
preserves the oriented topological knot type. No contact-isotopy
extension or approximation-uniqueness theorem is used.
\end{proof}

\begin{lemma}[Linking calculus and uniform transverse framing]
\label{fd:linking-calculus}
For two disjoint smooth oriented parametrized circles $C_1,C_2$ in
oriented $\RR^3$, use the normalized linking pairing
\begin{equation}\label{fd:gauss-linking}
 \ell(C_1,C_2)=\frac1{4\pi}\int_{S^1\times S^1}
       G\cdot(G_u\times G_v)\,du\,dv,
 \qquad G(u,v)=\frac{C_2(v)-C_1(u)}{|C_2(v)-C_1(u)|}.
\end{equation}
It is symmetric and is constant under smooth families of disjoint
oriented pairs. Call a direction $\nu\in S^2$ \emph{generic for the pair} if at
every parameter pair $(u,v)$ at which $C_2(v)-C_1(u)$ is parallel to $\nu$
the tangents of the two circles projected along $\nu$ are linearly
independent; the projection along such a $\nu$ displays the pair with
finitely many transverse mixed crossings. Let $\nu$ point toward the
observer and orient the projection plane so that its positive basis followed
by $\nu$ is positive in $\RR^3$; at a mixed crossing the strand nearer the
observer is over. Then $\ell(C_1,C_2)$ equals one half the sum over the
mixed crossings of the overpass-first sign $\sgn\det(u_{\rm o},u_{\rm u})$
of the projected over and under tangents.
Thus it has exactly the linking normalization used in the source's
front calculations.

More generally, let $C_s$, $0\leq s\leq1$, be a smooth family of embedded oriented
circles and let $v_s$ be a smooth vector field along them, everywhere
linearly independent of $\partial_uC_s$. One common sufficiently small
positive $\epsilon$ gives disjoint framed pairs
$(C_s,C_s+\epsilon v_s)$ throughout the family. Their pairing is
independent of that radius and of $s$. This includes a homotopy of
normal framings on a fixed curve.

If $T_s:S^1\to\RR^3$, $0\leq s\leq1$, is a smooth family of
embeddings with $(dz-y\,dx)(\partial_uT_s)>0$, there is one
$\epsilon_0>0$ such that $T_s$ and $T_s+\epsilon\partial_y$ are
disjoint positive transverse embeddings for every $s$ and
$0<\epsilon\leq\epsilon_0$. The number
\begin{equation}\label{fd:framed-linking}
            sl(T_s)=\ell(T_s,T_s+\epsilon\partial_y)
\end{equation}
is independent of such $\epsilon$ and of $s$.
\status{proved (refereed: bench A, 2026-09-06T00:35:51Z and 2026-09-06T01:58:43Z; transcribed from RC d2d\_contact\_constructions (C030); round~3: the convention sentence moved to Remark~\ref{rem:sl-convention}, F-25-81; round~5: the projection-genericity class of the mixed-crossing clause defined in the statement (a direction regular for the pair) and the observer direction and plane orientation that fix its sign stated there, as its proof uses them, F-25-155)}
\end{lemma}
\begin{proof}
The denominator in \eqref{fd:gauss-linking} is bounded away from zero
on the compact parameter torus. It remains so uniformly in any compact
smooth family of disjoint pairs. All derivatives and differentiations
under the integral are therefore ordinary smooth calculus on a compact
domain. The integrand can equivalently be written
\begin{equation}\label{fd:gauss-integrand}
 \frac{(C_1(u)-C_2(v))\cdot(C'_1(u)\times C'_2(v))}
 {|C_1(u)-C_2(v)|^3}.
\end{equation}
Indeed radial derivative terms disappear in the determinant, and the
derivative in $u$ has the sign from $-C'_1$. Exchanging the two curves
negates both the difference vector and the cross product, leaving this
scalar unchanged after exchanging the integration variables. This proves
symmetry, including the source's possible ordering of the pushoff first.

For a smooth family $G(u,v,s)$ into the unit sphere set
\begin{equation}
 A=G\cdot(G_u\times G_v),\qquad
 B=G\cdot(G_s\times G_v),\qquad
 C=G\cdot(G_u\times G_s).
\end{equation}
The vectors $G_u,G_v,G_s$ all lie in the two-dimensional plane
orthogonal to $G$. Hence every scalar triple product of these three
vectors vanishes. On expanding the three derivatives, the two terms
containing $G_{uv}$ cancel because the cross product is antisymmetric;
the remaining terms match by equality of mixed derivatives. Explicitly,
\begin{align}
 \partial_s A
 &=G\cdot(G_{us}\times G_v+G_u\times G_{vs}),\notag\\
 \partial_u B+\partial_v C
 &=G\cdot(G_{su}\times G_v+G_s\times G_{vu}
                   +G_{uv}\times G_s+G_u\times G_{sv})\notag\\
 &=\partial_s A.
 \label{fd:gauss-divergence}
\end{align}
Integrating a full period in $u$ kills $\partial_u B$ by the
fundamental theorem of calculus; integrating a full period in $v$
kills $\partial_v C$. Thus the derivative of
\eqref{fd:gauss-linking} is zero. This is a direct calculation, not
an appeal to invariance of degree or to an ambient-isotopy theorem.

We verify the normalization used in a generic diagram. Let $\omega$
be the outward sphere area form divided by $4\pi$. In positively
oriented Cartesian coordinates $(X,Y,Z)$ on the unit sphere, put
\begin{equation}\label{fd:sphere-primitive}
 \lambda_N=-\frac{X\,dY-Y\,dX}{4\pi(1-Z)}
          \quad\hbox{on }S^2\setminus\{N\},\qquad N=(0,0,1).
\end{equation}
This is smooth also at the south pole. In polar coordinates about
$N$, it is $-(1+\cos\theta)d\phi/(4\pi)$; differentiation
gives $\sin\theta\,d\theta\wedge d\phi/(4\pi)=\omega$.
On the positively oriented boundary of a cap of angular radius $h$,
its integral is $-(1+\cos h)/2$, tending to $-1$ as $h\to0$.

Suppose $N$ is a regular value of the particular map $G$ under
consideration. Its preimage is a finite set: the inverse function
theorem makes it discrete, and it is closed in a compact torus.
Choose disjoint inverse charts there. Compactness of their complement
then gives a small cap whose full inverse image lies in those charts.
Delete the inverse cap discs from the torus. On what remains,
$G^*\omega=d(G^*\lambda_N)$. The boundary integral formula gives
the integral as the sum of the integrals around the deleted discs
with negative boundary orientation. A disc whose local orientation
sign is $\sigma$ contributes $\sigma(1+\cos h)/2$.
The omitted disc integrals tend to zero: the pullback form is smooth
and their areas tend to zero in the fixed inverse charts. Consequently
\begin{equation}\label{fd:regular-pole-count}
                  \int_{S^1\times S^1}G^*\omega
                     =\sum_{p\in G^{-1}(N)}\sigma(p).
\end{equation}
The empty-preimage case uses the primitive on the entire torus and
has both sides zero. No assertion that an arbitrary value is regular
was used. A positive rotation of the coordinates proves the same
formula for any specified regular pole, including the opposite pole.

For completeness, the boundary integral formula just used requires
only planar Green's formula. Cover the compact punctured torus by
finitely many interior or boundary coordinate patches and choose
smooth bump functions supported in the patches whose sum is positive
there; division by that sum gives a finite partition of unity.
Apply the planar formula to each supported one-form. For an interior
patch the two integrals of partial derivatives vanish by the
fundamental theorem and Fubini. In a boundary patch straighten the
smooth boundary to a coordinate graph; the same one-variable
integration leaves exactly its oriented boundary integral.
The derivatives of the partition functions cancel because their sum
is one. This proves the required formula without a topological
degree theorem, general transversality theorem or source assumption.

Now let $\nu$ point toward the observer, and orient the projection
plane so that its positive basis followed by $\nu$ is positive in
$\RR^3$. In a generic projection, a crossing between the components
is exactly a preimage of $\nu$ or $-\nu$ under $G$; the projected
tangents are independent, so both poles are regular values, with
the empty-preimage case allowed. At $G=\nu$, the second component
is over. Up to a positive common factor the projected derivatives
of $G$ are $-C'_1,C'_2$. Their determinant in the tangent plane
oriented by $\nu$ is
\begin{equation}
       \det_\nu(-C'_1,C'_2)=\det_\nu(C'_2,C'_1),
\end{equation}
the overpass-first crossing sign. At $G=-\nu$, the first component
is over and the sphere tangent orientation is reversed. Its sign is
therefore
\begin{equation}
       \det_{-\nu}(-C'_1,C'_2)=\det_\nu(C'_1,C'_2),
\end{equation}
again the overpass-first sign. Formula~\eqref{fd:regular-pole-count}
at each pole says that each of these two signed crossing sums equals
$\ell$. Adding them and dividing by two proves the claimed mixed
crossing formula. The argument concerns the generic diagrams actually
used in the source computations, and does not invoke a theorem asserting
genericity of every projection.

It remains to supply a uniform framed pair, rather than assume that
pointwise small pushoffs assemble into a family. For the general
framing assertion put $n_s=\partial_uC_s\times v_s$ and
\begin{equation}\label{fd:family-normal-chart}
 F(s,u,r,t)=\bigl(s,C_s(u)+r v_s(u)+t n_s(u)\bigr).
\end{equation}
At $(r,t)=(0,0)$ its derivative is invertible: the last three
columns have determinant $|\partial_uC_s\times v_s|^2>0$.
The inverse function theorem, continuity and compactness give a
common small closed normal disc on which every derivative remains
invertible. For the interval endpoints use a smooth local extension
in $s$; the same estimates restrict back to $[0,1]$.
There is a still smaller common radius on which every fixed-$s$
normal chart is injective. Otherwise choose collisions at normal
radii tending to zero. Compactness gives subsequences with
$s_j\to s_*$, $u_j\to u_*$, and $u'_j\to u'_*$.
Equality of their limits gives $C_{s_*}(u_*)=C_{s_*}(u'_*)$;
embeddedness gives $u_*=u'_*$ on the circle. Both colliding points
then lie in one inverse chart of the four-dimensional map $F$
near $(s_*,u_*,0,0)$, a contradiction. This argument includes
both nearby and remotely separated parameter pairs.

Thus for a common $\epsilon_0>0$, the normal charts embed the
closed disc of radius $2\epsilon_0$ at every $s$. Restriction to
the circles with normal coordinates $(0,0)$ and $(\epsilon,0)$,
$0<\epsilon\leq\epsilon_0$, gives disjoint oriented embedded
pushoff pairs varying smoothly in $s$. Interpolating between two
positive radii stays in the same disc. The already proved paired
invariance therefore gives both family and radius independence.
This proves the general framing assertion, including homotopies
of normal fields without any ambient extension theorem.

In particular, for a Legendrian knot choose a smooth contact field
$w$ independent of its tangent. Such a field exists explicitly:
if $L'=A(\partial_x+y\partial_z)+B\partial_y$, take
$w=-B(\partial_x+y\partial_z)+A\partial_y$; $A^2+B^2>0$.
The homotopy $v_\theta=\cos\theta\,w+\sin\theta\,\partial_z$,
$0\leq\theta\leq\pi/2$, never becomes tangent to $L$.
For $\theta>0$ its contact-form evaluation is $\sin\theta>0$;
at zero independence follows from the choice of $w$.
The general framing assertion therefore justifies the replacement
by the $z$-direction pushoff in the source's front calculation of
$tb$, without an assumption about homotopic framings.

For $C_s=T_s$ and $v_s=\partial_y$, independence of the two
columns follows by applying $dz-y\,dx$ to a linear relation:
it is positive on $\partial_uT_s$ and zero on $\partial_y$.
Write $T_s=(x_s,y_s,z_s)$. Let $a_0>0$ be the minimum of
$z'_s-y_sx'_s$ and let $M$ bound $|x'_s|$ on the same compact
cylinder. Decrease $\epsilon_0$ until
$\epsilon_0M<a_0/2$ (no restriction from this inequality if $M=0$).
Then
\begin{equation}\label{fd:pushoff-positive-clearance}
 (dz-y\,dx)_{T_s+\epsilon\partial_y}
                   (\partial_uT_s)
       =(z'_s-y_sx'_s)-\epsilon x'_s>a_0/2>0.
\end{equation}
The uniform normal-chart construction has already proved that
these are disjoint embedded pairs with the required orientations.
Fixing any $0<\epsilon\leq\epsilon_0$ now gives the smooth family
to which \eqref{fd:gauss-divergence} applies. Interpolating between
two such positive radii stays inside the same normal chart and
proves radius independence as well. Finally $\partial_y$ is a
globally nonzero section of $\ker(dz-y\,dx)$, exactly the source's
standard-space framing. Its local linking computations use the
mixed-crossing formula proved above. This fixes the convention in
\eqref{fd:framed-linking} and proves the required smooth-family
self-linking invariance without importing it as a separate theorem.
\end{proof}

\begin{remark}[the self-linking convention]\label{rem:sl-convention}
In~\eqref{fd:framed-linking} the global nonzero contact section
$\partial_y$ is the section in Etnyre's standard-space self-linking
convention (Etnyre~\cite[Sections~2.6.3--2.6.4, p.~15]{Etnyre};
independence of the chosen nowhere-zero section of the contact structure:
Geiges~\cite[Definition~3.2 and the following paragraph, p.~45]{GeigesContact},
whose $\partial_X$ is $-\partial_y$ under the coordinate change
$(X,Y,Z)=(-y,x,z)$); no surface-dependent choice or extension is being
made. This identifies the number of Lemma~\ref{fd:linking-calculus} with
the literature's self-linking number; the lemma itself proves only the
invariance and the normalization stated there (F-25-81).
\end{remark}


\subsection{Exact regular rounding of the cusp endpoint}
\label{ce:geometry}

\begin{lemma}[Ordinary rounding of the exact cusp germs]
\label{ce:rounding}
Let $L$ be a smooth oriented spatial embedding of a finite nonempty union
of parameter circles in $\mathbb R^3$, and write $p=(x,z)$ for its
projection. Its only failures of regularity are finitely many isolated
cusps; all other projected coincidences are finitely many transverse
double points. Assume there are no triple points or cusps on another
branch, and the two $y$ heights at every double point are distinct.
At every cusp assume the exact germ, on a parameter interval with smooth
coordinate $u=y-y_0$,
\begin{equation}\label{ce:exact-germ}
 x=x_0+Au^2,\qquad y=y_0+u,\qquad
 z=z_0+Ay_0u^2+\frac{2A}{3}u^3,\qquad A\ne0.
\end{equation}
The coordinate $u$ may increase or decrease along the prescribed component
orientation, which is not changed. The formula is a hypothesis, not an
appeal to a classification of singularities.

There is a jointly smooth family $L_\lambda$, $0\leq\lambda\leq1$, of
oriented spatial embeddings with $L_0=L$, fixed outside disjoint cusp
parameter intervals, such that for every $\lambda>0$ its $xz$ projection
is an ordinary finite regular generic diagram. It cleanly smooths the
cusps, creates no crossing, and retains every original crossing with
its oriented decorated data. All original parameter circles, component
labels and traversal orientations are retained. With no cusps take the
constant family.
This is an ordinary spatial deformation, not asserted Legendrian or
positive transverse, and it carries no self-linking transport statement.
\status{proved (refereed: bench A, 2026-09-05T15:51:04Z; transcribed from RC ce\_rounding (C030))}
\end{lemma}
\begin{proof}
\par\noindent\emph{Disjoint clean supports and positive charts.} 
Restrict one exact-germ interval so that $u$ is a coordinate throughout.
The affine page chart
\begin{equation}\label{ce:positive-chart}
 X=\frac{x-x_0}{A},\qquad
 Z=\frac{3\bigl(z-z_0-y_0(x-x_0)\bigr)}{2A}
\end{equation}
sends the front to $(u^2,u^3)$. Its linear determinant is
$3/(2A^2)>0$ for either sign of $A$, so it preserves page orientation.
The cubic coordinate is injective, hence the full local arc is injective
even at its cusp.

Choose a smaller closed parameter interval about the cusp inside the
chart. The cusp image is absent from the compact image of the complement
of that interval's interior: local injectivity and the no-other-branch
hypothesis exclude
every other preimage, including on different components. Its distance
from that complement is positive. Thus a sufficiently small open page
neighbourhood $V$ meets only the exact-chart arc. The finitely many cusp
images are distinct; shrink their neighbourhoods to make them pairwise
disjoint and disjoint from every double point. Choose $b>0$ so small
that the closed subarc $|u|\leq b$ lies inside $V$ and the exact chart.

Let $\rho$ be a smooth real cutoff equal to one near zero, with support
strictly inside $(-b,b)$. The image of its compact support has positive
distance $d$ from the closed complement of $V$. Put
$M=\max|u\rho(u)|$ on that support; $M>0$ because $\rho=1$ near zero.
Choose $\epsilon>0$ with
\begin{equation}\label{ce:support-clearance}
 |A|\epsilon M\sqrt{1+y_0^2}<d/2.
\end{equation}
Make these choices at each of the finitely many cusps. No continuous
selection over a further family or generic projection theorem is used.

\par\noindent\emph{The complete rounding family.} 
In the normalized chart set
\begin{equation}\label{ce:rounding-formula}
 X_\lambda(u)=u^2+\lambda\epsilon u\rho(u),\qquad
 Z_\lambda(u)=u^3,\qquad y_\lambda(u)=y_0+u.
\end{equation}
Inverting \eqref{ce:positive-chart} shows that the physical changes are
\begin{equation}\label{ce:physical-displacement}
 \Delta x=A\lambda\epsilon u\rho(u),\qquad \Delta y=0,\qquad
 \Delta z=Ay_0\lambda\epsilon u\rho(u).
\end{equation}
The cutoff vanishes on collars of the chart ends, so every derivative
of the change vanishes there. Pasting to the unchanged $L$ therefore
gives a jointly smooth family on all parameter circles and the closed
$\lambda$ interval. The formula even extends smoothly as a map to a
slightly larger time interval; embeddedness outside $[0,1]$ is not claimed.

For $u\ne0$ one has $dZ_\lambda/du=3u^2>0$. At $u=0$,
$dX_\lambda/du=\lambda\epsilon$, since $\rho(0)=1$. Thus the two
projected derivatives are never simultaneously zero for $\lambda>0$,
regardless of the size or sign of a cutoff derivative. Composing with
the smooth coordinate $u$ preserves this when $du/d\theta<0$ as well.
At $\lambda=0$ the only local projected singularity is the original
cusp. Spatial immersion holds for every $\lambda$ because
$dy_\lambda/du=1$ in each modified chart and the rest of $L$ is unchanged.

Throughout the whole exact chart, including cutoff collars and unchanged
tails, the coordinate $Z_\lambda$ is exactly $u^3$. For $u<v$,
\begin{equation}\label{ce:cubic-difference}
 v^3-u^3=(v-u)(u^2+uv+v^2)>0.
\end{equation}
Indeed the second factor is $(u+v/2)^2+3v^2/4$ and can vanish only when
$u=v=0$, excluded here. Hence distinct parameters anywhere in that
chart have different projected points at every $\lambda$. This proves
injectivity, not merely a tangent test at the centre or at an endpoint.

Equations~\eqref{ce:support-clearance} and~\eqref{ce:physical-displacement}
keep every moved point inside $V$, since its projected displacement norm
is at most $|A|\epsilon M\sqrt{1+y_0^2}<d/2$. No parameter outside the
exact chart projects into $V$. Therefore a moved point cannot meet a
remote parameter. The different cusp neighbourhoods are disjoint, so
different modifications cannot meet each other. Together with strict
local injectivity these facts exclude every new projected crossing for
every $\lambda\in[0,1]$, including pairs on different components.
All old double points are outside the chosen supports. Their exact
branches, tangent determinants and $y$ heights remain unchanged, so
transversality, crossing signs and O/U choices persist. There are no
other coincidences by hypothesis. For positive $\lambda$ the result is
therefore precisely a finite regular generic diagram.

Any spatial coincidence would project to a new coincidence or to an
old double point. The former were excluded; at the latter the two
unchanged $y$ heights are distinct. Spatial injectivity follows for the
whole family. Combined with the immersion already proved, compactness
of the finite parameter union and the Hausdorff target give an embedding:
the continuous bijection to its image has continuous inverse.
Parameters and component labels have never changed.

Each modified portion is one regular embedded oriented arc in a clean
cusp neighbourhood, agreeing with the old arc in its endpoint collars
and containing no crossing. It is exactly a clean cusp smoothing.
There is no assertion that it stays in the contact plane. For example,
where $\rho=1$, differentiating the physical formulas gives
\begin{align}
 x'_\lambda&=2Au+A\lambda\epsilon,
 &y_\lambda&=y_0+u,\label{ce:contact-derivatives}\\
 z'_\lambda&=2Ay_0u+2Au^2+Ay_0\lambda\epsilon.
 \label{ce:contact-z-derivative}
\end{align}
Subtracting $y_\lambda x'_\lambda$ from the second line yields
\begin{equation}\label{ce:ordinary-not-contact}
 z'_\lambda-y_\lambda x'_\lambda=-A\lambda\epsilon u.
\end{equation}
This can have either sign and vanishes at the centre. It cannot justify
Legendrian or positive-transverse invariance, nor transport self-linking.
\end{proof}

\begin{corollary}[Scalar independence of clean cusp smoothing]
\label{ce:smoothing-record}
For a front satisfying Lemma~\ref{ce:rounding}, define a permitted
smoothed-front diagram by replacing each cusp with one regular embedded
oriented arc in a clean cusp neighbourhood, agreeing with the old germs
in endpoint collars and creating no crossing. Retain every double point
with its original height choice. Every such actual diagram has the same
named decorated record as the diagram $D_\epsilon=p(L_1)$ of the lemma,
including crossing-free components. Consequently their source values
$F_D(l,m)$ and campaign polynomials $P_D(a,z)$ agree.
\status{proved (refereed: bench A, 2026-09-05T15:51:04Z; transcribed from RC ce\_rounding (C030))}
\end{corollary}
\begin{proof}
A cusp replacement introduces no crossing visit and changes no successor
of an old visit along the oriented parameter circle. All crossing
pairings, signs and O/U bits remain unchanged because their germs lie
outside the cusp neighbourhoods. The component correspondence is the
identity on the original circles, including those without crossings.
Thus the two actual diagrams have a named decorated-record isomorphism.
Lemma~\ref{rp:record-polynomial}, with the same source construction of
Literature input~\ref{lp:lm}, gives equality of $F_D$; the common substitution
of Theorem~\ref{lp:core} preserves it. This is a scalar statement, not
a classification of arbitrary relative arc embeddings or an ambient
completeness assertion.
\end{proof}

\begin{remark}\label{ce:scope}
The exact cusp formula is supplied in the actual use by
Theorem~\ref{fd:generic-front}. The geometric lemma alone neither proves
a per-front polynomial inequality nor asserts that an arbitrary smooth
transverse knot has a finite generic specified projection. Polynomial
transport to a supplied ordinary endpoint uses the distinct compact
smooth-family endpoint comparison supplied below.
Self-linking uses the separate genuine transverse-pushoff family, never
the ordinary rounding above.
\end{remark}

\subsection{Contact-path polynomial transport and the representative bound}

\begin{lemma}[The supplied contact-path diagram endpoints]
\label{cp:finite-contact-path}
Let $C$ be a finite nonempty disjoint union of oriented parameter circles,
and let $L:C\to\RR^3$ be a smooth embedding with nonvanishing parameter
derivative. Assume its $xz$ projection $F$ is regular except at finitely
many cusps, each with the exact germ
\begin{equation}\label{cp:exact-cusp}
 x=x_0+Au^2,\qquad y=y_0+u,\qquad
 z=z_0+Ay_0u^2+\frac23Au^3,\qquad A\ne0.
\end{equation}
Assume the only multiple points of $F$ are finitely many transverse
double points, none at a cusp, and the two $y$ values differ at every
double point. Suppose a supplied jointly
smooth family of oriented spatial embeddings $G_t$ starts at this
parametrized $L$ and ends at $T$, whose specified $xz$ projection $D_T$
is an ordinary finite regular generic diagram. Then every clean ordinary
cusp smoothing $S(F)$, with smaller $y$ over at the unchanged crossings,
has the same campaign polynomial as $D_T$:
\begin{equation}\label{cp:endpoint-polynomial}
                    P_{S(F)}=P_{D_T}.
\end{equation}
No contact condition is imposed on $G_t$ or on the rounding family.
\status{new (round~5: the uniform normal-injectivity radius of the ambient extension supplied, by the argument of the uniform normal chart in the proof of Lemma~\ref{fd:linking-calculus} applied to the family $H_s$ of embeddings of the finite union $C$, F-25-151; Corollary~\ref{ce:smoothing-record} cited for the smoothing-independence of $P_{S(F)}$, F-25-136; the remainder is the round-1 transcription of RC d2e\_contact\_bound (C030), held by bench~A; round~9: the status word at the L1 premise removed, F-25-180)}
\end{lemma}
\begin{proof}
Use Lemma~\ref{ce:rounding} on the exact germs
\eqref{cp:exact-cusp}. It supplies a jointly smooth family $L_\lambda$
of spatial embeddings, $L_0=L$, fixed outside finitely many clean cusp
neighbourhoods. Its endpoint $L_1$ has an ordinary finite regular generic
diagram $D_\epsilon$, with all old crossings and every component retained.
The construction keeps $y=y_0+u$ and, in the positive affine page chart,
uses $X_\lambda=u^2+\lambda\epsilon u\rho(u)$, $Z_\lambda=u^3$.
The determinant of this chart is $3/(2A^2)>0$, so it has not changed the
crossing convention. Lemma~\ref{ce:rounding} proves injectivity on the
whole chart, including cutoff collars, and excludes remote intersections
for every parameter value; an endpoint tangent check alone would not do so.

The two actual diagrams $S(F)$ and $D_\epsilon$ have identical named
decorated records. Each clean cusp replacement inserts no crossing visit
and keeps the same oriented attachments. All crossing pairings, cyclic
successors, signs, O/U bits and crossing-free component circles agree.
Corollary~\ref{ce:smoothing-record} (through
Lemma~\ref{rp:record-polynomial}) therefore gives
\begin{equation}\label{cp:smoothing-record}
                    P_{S(F)}=P_{D_\epsilon}.
\end{equation}
This is scalar record equality, not a claimed classification of arbitrary
relative arc embeddings.

Let $\eta:[0,1]\to[0,1]$ be smooth and increasing, with its endpoint
values $0,1$ and every positive-order derivative zero at both endpoints.
For example normalize the integral of the bump
$\exp(-1/(s(1-s)))$ on $(0,1)$, extended by zero. Put
\begin{equation}\label{cp:flattened-family}
 H_s=\begin{cases}
       L_{1-\eta(2s)},&0\leq s\leq\tfrac12,\\
       G_{\eta(2s-1)},&\tfrac12\leq s\leq1.
      \end{cases}
\end{equation}
Both branches equal $L$ at the join. Every positive-order time derivative
vanishes there; mixed parameter derivatives have the same property,
while pure circle derivatives agree with those of $L$. Thus the chain
rule gives joint smoothness, including at the join. Every slice is one
of the supplied or proved embeddings. The literal endpoints are $L_1,T$,
with their given parameter orientations, not unidentified nearby knots.

Literature input~\ref{lit:homfly} retains the original full global
HOMFLY interface. The supplied smooth family $H_s$ is a smooth isotopy
of the actual oriented embedded circles. To make the ambient premise
explicit, extend its velocity to a compactly supported ambient field
as follows. At each $(s,u)$ choose a local normal chart for the embedded
circle, smoothly also in $s$, by the inverse function theorem; finitely
many charts cover the compact parameter cylinder. Restrict every chart to
one common normal radius on which a spatial point at time $s$ has at most
one normal-bundle preimage. Such a radius exists by the argument of the
uniform normal chart in the proof of Lemma~\ref{fd:linking-calculus},
applied to the normal bundle of the jointly smooth family $H_s$ of
embeddings of the finite union $C$: two distinct normal-bundle points at a
common time with a common image and normal lengths tending to zero would
accumulate, by compactness of the cylinder, at one point of the cylinder
--- the slice $H_{s_*}$ being an embedding of $C$ --- near which a single
normal chart is injective, a contradiction. Below that radius a point of
the tracked curve lies in a chart only as its own foot point. In each
chart extend
$\partial_sH_s(u)$ constantly in its two normal coordinates. Choose
smooth spatial-time cutoffs supported in those charts whose sum is
positive on the tracked curve, and divide by their sum near that curve.
Their weighted sum equals $\partial_sH_s(u)$ on the curve, because every
local extension has that same value there. Multiply by a further cutoff
equal to one near the whole compact tracked cylinder and zero outside a
compact neighbourhood. This gives a jointly smooth bounded spatial
vector field, with bounded derivatives on the compact time interval.
At the time endpoints extend the given flattened family smoothly to a
small interval; the same construction restricts to $[0,1]$.

Smooth ODE existence and continuation give its flow for the full compact
time interval. Uniqueness and reverse-time evolution give an ambient
diffeomorphism at each time. The chain rule and uniqueness imply that
the flow carries the parametrized $H_0$ to $H_s$. It fixes the complement
of a compact set and preserves ambient orientation, since its derivative
determinant starts positive and never vanishes. Thus its endpoint gives
an actual ambient isotopy between the two oriented links. The retained
global source premise yields $H_{D_\epsilon}=H_{D_T}$.
Theorem~\ref{lp:core} identifies these original values with $P$ on both
actual diagrams, proving $P_{D_\epsilon}=P_{D_T}$. Combining this equality
with~\eqref{cp:smoothing-record} proves~\eqref{cp:endpoint-polynomial}.
This step explicitly consumes the global L1 premise,
Literature input~\ref{lit:homfly}; it is not a source-free replacement
or an application of the local LM theorem beyond its stated domain.

Only the endpoints were assumed to have ordinary generic diagrams. The
cusped diagram at the join causes no problem because $H_s$ is a spatial
embedding throughout. The rounding is ordinary, not Legendrian or
positive transverse: where $\rho=1$ its contact-form evaluation in the
$u$ coordinate is $-A\lambda\epsilon u$. No self-linking number is
transported along this family.
\end{proof}

\begin{definition}[generic positive transverse front]\label{def:transverse-front}
In $(\RR^3,\ker(dz-y\,dx))$, a \emph{generic positive transverse front} is
the oriented knot diagram in the $(x,z)$ plane obtained from a smooth oriented
embedded knot $T$ with $z'-yx'>0$ whose $xz$ projection is an immersion of
the parameter circle with finitely many transverse double points and no
triple point, the over strand at each double point being the branch of
smaller $y$. It has no cusp. At a vertical tangent, $x'=0$, the inequality
gives $z'>0$: every vertical tangent of a generic positive transverse front
points upward, a consequence and not a hypothesis. The diagram determines
its writhe and its over/under counts; the knot $T$ is named separately where
it is used.
\end{definition}

\begin{literature}[Front and positive-pushoff inputs]\label{src:contact}
Use standard contact space
$(\RR^3,\ker(dz-y\,dx))$, with positive transverse orientation
$z'-yx'>0$ (Etnyre~\cite[Sections~2.1 and~2.4]{Etnyre} for the
conventions). Etnyre's front and pushoff
statements~\cite[Section~2.6.2, equations~(5) and~(7); Section~2.6.4,
equation~(9); Section~2.9, Lemma~2.22, equation~(17)]{Etnyre} give, for
an oriented Legendrian front with downward and upward cusp counts $D,U$
--- every cusp of an oriented front is traversed either downward or
upward, so the total cusp count in his $tb$ formula is $D+U$ ---
\begin{equation}\label{fd:contact-inputs}
 r=\frac{D-U}{2},\qquad tb=w-\frac{D+U}{2},\qquad
 sl(T_+(L))=tb(L)-r(L).
\end{equation}
(For the second formula, the Legendrian $tb$, the printed proof used is
Geiges~\cite[Proposition~3.5.9, p.~117]{Geiges}; Etnyre's equation~(7)
reduces to his Remark~2.14, which is stated there without proof.)
For a generic positive transverse front
(Definition~\ref{def:transverse-front}), self-linking equals its front
writhe: Geiges~\cite[Section~3.1, Lemma~3.3, p.~46]{GeigesContact} (printed
proof) and Geiges~\cite[Proposition~3.5.32, p.~127]{Geiges}; Etnyre's
equation~(9), Section~2.6.4, states it. These are the local front and
pushoff source formulas only.
\status{lit: Etnyre and Geiges local front and pushoff formulas; Registry~\ref{reg:etnyre}; round~3: page-precise locators, the cusp count $D+U$ and the front class named, F-25-82, F-25-83, F-25-84}
\end{literature}

Lemma~\ref{fd:transverse-neighborhood} constructs an oriented Legendrian
knot near each smooth positive transverse knot $T$, a chosen positive
pushoff transversely isotopic to $T$, and an explicit ordinary ambient
flow to $T$. Theorem~\ref{fd:generic-front} supplies a finite generic
front with the exact germs \eqref{cp:exact-cusp}, by a smooth ambient
contact isotopy carrying the chosen pushoff. These are internal
constructions. The linking normalization and invariance along smooth
positive transverse families are proved in Lemma~\ref{fd:linking-calculus}.
The local diagram proof in Theorem~\ref{ng:local-front-bound}, using
the finite source reduction \ref{ng:finite-word}, gives the following
bound, restated here as a lemma so that its consumer below cites it.

\begin{lemma}[the individual-front bound in self-linking form]\label{fd:ng-bound}
Let $F$ be a front on the domain of Definition~\ref{ng:front-domain}, with
$c_\downarrow(F)=D(F)$ its number of downward cusps, and let $S(F)$ be an
actual clean ordinary cusp smoothing of $F$ as in that definition. Then
\begin{equation}\label{fd:ng-input}
 sl_{\rm Ng}(F):=w(F)-c_\downarrow(F)
          \leq-\max\deg_a P_{S(F)}(a,z)-1.
\end{equation}
This is not an application of an ambient-invariant hypothesis to the
local diagram construction, nor a maximum over front representatives.
\status{new (round~3; Theorem~\ref{ng:local-front-bound} restated in Ng's self-linking form, so that the closure of Theorem \texttt{fd:contact} contains it; consumes Literature input~\ref{ng:finite-word} through that theorem; F-25-107; round~5: Lemma~\ref{ng:smoothing-record} cited for the symbol $P_{S(F)}$, F-25-136)}
\end{lemma}
\begin{proof}
Theorem~\ref{lp:core} gives $P_{S(F)}\neq0$, and $P_{S(F)}$ does not depend
on the rounding by Lemma~\ref{ng:smoothing-record}, so $\deg_a P_{S(F)}$ of
display~\eqref{ng:defect} is the largest $a$ exponent with nonzero
coefficient, $\max\deg_a P_{S(F)}$. Inequality~\eqref{ng:front-inequality}
of Theorem~\ref{ng:local-front-bound} reads
$w(F)-D(F)\leq-\deg_aP_{S(F)}-1$, which is the display with
$c_\downarrow(F)=D(F)$.
\end{proof}

Lemma~\ref{cp:finite-contact-path} supplies the separate polynomial
transport to the particular ordinary transverse diagram when it is used.

\begin{theorem}[The contact dictionary and finite-diagram representative bound]
\label{fd:contact}
In the $xz$ front page the smaller-$y$ branch is over, and its crossing
sign is $\sgn\det_{xz}(u_O,u_U)$. For a generic positive transverse front
(Definition~\ref{def:transverse-front}) one has
\begin{equation}\label{fd:front-writhe}
 sl(T)=\sum_q\sgn\det_{xz}(u_O(q),u_U(q)).
\end{equation}
Let $T$ be an individual smooth positive transverse knot whose specified
$xz$ projection $D_T$ is an ordinary finite regular generic diagram.
With $P_T$ denoting the original campaign polynomial of this actual
diagram, one has
\begin{equation}\label{fd:representative-bound}
 sl(T)\leq-\max\deg_a P_T(a,z)-1.
\end{equation}
The finite specified-diagram hypothesis is part of this auxiliary bound;
no assertion that every smooth knot has a generic specified projection
is required. The carrier-floor construction supplies such a diagram.
\status{new (round~6: the crossing from the source's self-linking number to the number of \eqref{fd:framed-linking} cites Remark~\ref{rem:sl-convention}, F-25-168; round~5: the domain hypothesis of Lemma~\ref{fd:ng-bound} discharged for $F_T$ from Theorem~\ref{fd:generic-front} and the lemma cited by label, F-25-153; the sentence after \eqref{fd:front-writhe} restored to RC's wording, F-25-152; the remainder is the round-1 transcription of RC fd:contact, d2e\_contact\_bound (C030), held by bench~A; consumes Literature inputs~\ref{src:contact} and~\ref{lit:homfly}; round~3: the front class of \eqref{fd:front-writhe} named, F-25-82)}
\end{theorem}
\begin{proof}
Writing $\alpha=dz-y\,dx$ gives $d\alpha=dx\wedge dy$ and
\begin{equation}
 \alpha\wedge d\alpha=dz\wedge dx\wedge dy=dx\wedge dy\wedge dz.
\end{equation}
The source observer is on the negative-$y$ side looking in the positive
$y$ direction, so smaller $y$ is over. For a Legendrian front this is also
the smaller-slope rule $y=dz/dx$. The normal toward the observer is
$\nu=-\partial_y$ and
$\det_{xyz}(\partial_x,\partial_z,\nu)=1$.
For projected over and under tangents $(a,c)$ and $(b,d)$,
\begin{equation}
 \det\begin{pmatrix}a&b&0\\0&0&-1\\c&d&0\end{pmatrix}
 =ad-bc=\det_{xz}((a,c),(b,d)).
\end{equation}
This identifies each source crossing sign with the campaign sign.
Summing the source transverse-front writhe formula, whose $sl$ is the
number of \eqref{fd:framed-linking} by Remark~\ref{rem:sl-convention},
proves \eqref{fd:front-writhe}. Absence of downward tangencies alone is not a
converse transverse-lift theorem: the crossing-height inequalities will be
checked separately in the carrier-floor proof.

With the same vertical coordinate, Ng's downward cusp count is $D$.
Substitute the first two formulas of \eqref{fd:contact-inputs} into the third:
\begin{align}
 tb-r&=w-\frac{D+U}{2}-\frac{D-U}{2}\notag\\
     &=w-\frac{2D}{2}=w-D=sl_{\rm Ng}(F).
\end{align}
This identifies Ng's pushoff as Etnyre's positive pushoff, fixing the
possible exchange of names in other conventions.

Fix $T$ with the stated finite ordinary generic diagram $D_T$.
Lemma~\ref{fd:transverse-neighborhood} gives the original Legendrian
helix $L$ and a chosen positive pushoff transversely isotopic to $T$.
Let $\Phi_s$ be the ambient contact isotopy of
Theorem~\ref{fd:generic-front}; write $L_T=\Phi_1\circ L$ and let
$F_T$ be its finite generic front, with the exact cusp germs
\eqref{cp:exact-cusp}. The isotopy carries the chosen pushoff through
positive transverse embeddings, so $T_+(L_T)$ is transversely isotopic
to the same $T$. Lemma~\ref{fd:linking-calculus} preserves self-linking
along this smooth positive transverse family. The preceding cusp
calculation identifies its value with $sl_{\rm Ng}(F_T)$. The front
$F_T$ lies on the domain of Definition~\ref{ng:front-domain}:
Theorem~\ref{fd:generic-front} gives finitely many semicubical cusps with
the germ \eqref{cp:exact-cusp}, whose $x''(0)=2A\neq0$, transverse double
points, no triple point and no cusp on another branch, and Legendrianity
$z'=yx'$ makes the front tangent vanish wherever $x'=0$, so that every such
point is one of these cusps and no regular arc has a vertical tangency.
Lemma~\ref{fd:ng-bound}, applied to $F_T$, bounds it by
$-\max\deg_a P_{S(F_T)}-1$.

Let $\Psi_t$ be the ordinary ambient flow of
\eqref{fd:helix-ambient-isotopy}. The maps
\begin{equation}\label{fd:contact-ambient-composition}
 \Theta_t=\Psi_t\circ\Phi_{1-t}\circ\Phi_1^{-1},
 \qquad 0\leq t\leq1,
\end{equation}
give the supplied family $G_t=\Theta_t\circ L_T$.
Indeed $\Phi_1^{-1}\circ L_T=L$, so
$G_t=\Psi_t\circ\Phi_{1-t}\circ L$; at zero it equals $L_T$ and
at one it equals the parametrized $T$. The two smooth compactly
supported ambient flows make $G$ jointly smooth, and each slice is an
embedding with nonzero parameter derivative. They transport the given
circle orientation. The generic front theorem supplies all other cusp
and crossing hypotheses of Lemma~\ref{cp:finite-contact-path}; the
specified ordinary diagram $D_T$ is the endpoint hypothesis already
included in the present statement. Consequently that lemma gives
$P_{S(F_T)}=P_{D_T}=P_T$. Substitution proves
\eqref{fd:representative-bound}.

The polynomial equality explicitly uses the original global
HOMFLY source premise in Literature input~\ref{lit:homfly}. The positive transverse family alone
transports self-linking; the distinct ordinary flow and cusp rounding
are used only for the polynomial. No maximum over representatives,
mirror, variable change or new owner assumption enters this composition.
\end{proof}

\subsection{Tangent lifts and the full corner construction}

For a local polygon $L=(q_1,\ldots,q_c)$ put
$\delta_i=q_{i+1}-q_i$. Its principal angle is written $\vartheta_i$, as
in Definition~\ref{def:regular}; $\tau_i$ elsewhere in this manuscript
remains the sign of that angle. The rotation of a direction loop and of a
smooth regular curve is Definition~\ref{cf:def-turning}; the lemma that
follows it constructs the lifts and proves independence. The notation
$P_D$ is the local polynomial, identified with the original $H_D$ by
Theorem~\ref{lp:core}. Its global reversal and mirror deductions below
take the uniqueness they need from Lemma~\ref{lp:coefficient-transport}
and retain the existence clause of Literature input~\ref{lit:homfly}.

\begin{definition}[direction loops and smooth rotation]\label{cf:def-turning}
A \emph{direction loop} is a continuous map $T\colon\RR/\ZZ\to S^1$ into
the unit circle of the oriented plane. A \emph{tangent-angle lift} of $T$
at a seam is a continuous $\theta\colon[0,1]\to\RR$ with
$T(s)=(\cos\theta(s),\sin\theta(s))$ for $0\leq s\leq1$; put
\[
   \operatorname{tw}(T)=\frac{\theta(1)-\theta(0)}{2\pi}.
\]
For a closed $C^1$ regular oriented curve $\gamma\colon\RR/\ZZ\to\RR^2$,
that is, with $\gamma'\neq0$ everywhere, put
\[
   T_\gamma=\frac{\gamma'}{|\gamma'|},\qquad
   \rot(\gamma)=\operatorname{tw}(T_\gamma).
\]
For a polygon in the regular locus, $\rot$ remains that of
Lemma~\ref{lem:rot}. For either a polygon or a $C^1$ regular closed curve
put $R(L)=|\rot(L)|$. The lemma that follows constructs the lift and
proves that these values do not depend on its choice, on the seam, or on
an orientation-preserving regular reparametrisation. (Transcribed from CV
def:rot, paragraph ``Direction loops and smooth curves'', d1\_setup;
F-25-88.)
\end{definition}

The geometric constructions are transcribed from the pinned CV floor
proof. Its recorded counterpart attribution for the rational curl,
affine insertion and dual-vector argument is retained; printing these
proofs claims responsibility for the translation, not their authorship.

\begin{lemma}[tangent lifts and principal turns]\label{cf:lem-turnlift}
Here $\rot$ is as in Lemma~\ref{lem:rot} for polygons and
Definition~\ref{cf:def-turning} for direction loops and closed $C^1$
regular curves.
\begin{enumerate}
\item[(i)] Every continuous direction loop has a tangent-angle lift.
The integer $\operatorname{tw}(T)$ is independent of the lift and seam,
is unchanged by an orientation-preserving reparametrisation, and is constant
under a continuous homotopy through direction loops. Consequently the rotation
of a closed $C^1$ regular curve is invariant under regular homotopy, and
reversing its orientation negates it.
\item[(ii)] If $L\in\mathcal R_c$ has principal turns $\vartheta_i$, then
\[
   2\pi\rot(L)=\sum_i\vartheta_i.
\]
Consequently rotation is constant along every path in $\mathcal R_c$, is
unchanged by a positive flat subdivision, and is negated by orientation
reversal. Every regular three-corner polygon has rotation $+1$ or $-1$
according to the common sign of its three turns.
\item[(iii)] Suppose a closed $C^1$ regular curve is obtained from $L$ by
replacing every corner by a regular arc whose compatible tangent-angle lift
has increment $\vartheta_i$. Its rotation equals $\rot(L)$.

More generally, let $b,c$ be regular oriented arcs having the same endpoints
and the same oriented tangent rays at both endpoints, and suppose replacing
$b$ by $c$ in a closed curve gives closed $C^1$ regular curves $F_b,F_c$.
For compatible tangent lifts with equal initial values, write their increments
as $\Delta_b,\Delta_c$. Then
\[
   \rot(F_c)-\rot(F_b)=\frac{\Delta_c-\Delta_b}{2\pi}.
\]
If $A_t$, $0\leq t\leq1$, is a supplied continuous path in
$\mathrm{GL}^+(2,\mathbb R)$ with $A_0=I$ and $A_1=A$, applying $A$ to both
arcs leaves $\Delta_c-\Delta_b$ unchanged.
\end{enumerate}
\status{proved (refereed: bench A, 2026-09-05T18:22:57Z and 2026-09-06T04:22:56Z; transcribed from CV lem:turnlift (C032); round~3: the scoping sentence covers the smooth case, F-25-88)}
\end{lemma}
\begin{proof}
For~(i), uniform continuity gives
$0=s_0<\cdots<s_m=1$ such that
$T(s)\cdot T(s_{j-1})>0$ on every
$[s_{j-1},s_j]$. After choosing $\theta(0)$, continue it there by
\[
 \theta(s)=\theta(s_{j-1})+
 \operatorname{atan2}\!\left(
   \det(T(s_{j-1}),T(s)),\,
   T(s_{j-1})\cdot T(s)\right).
\]
The second argument is positive, so the added angle lies in
$(-\pi/2,\pi/2)$; the formulas match at the subdivision points and construct a
continuous lift without a covering-space theorem. Since $T(1)=T(0)$, its
endpoint increment belongs to $2\pi\mathbb Z$. Two lifts differ continuously
by a value in $2\pi\mathbb Z$, hence by one constant.

If the increment is $2\pi k$, extend the lift by
$\theta(s+n)=\theta(s)+2\pi nk$. Every interval of length one then has the same
increment, proving seam independence. An orientation-preserving
reparametrisation of the circle has an increasing representative
$\phi\colon\mathbb R\to\mathbb R$ with
$\phi(s+1)=\phi(s)+1$; the lift $\theta\circ\phi$ has the same increment.

For a homotopy $T(s,t)$ fix $t_0$. Uniform continuity gives a neighbourhood of
$t_0$ on which
$T(s,t)\cdot T(s,t_0)>0$ for every $s$. There is therefore
a unique continuous relative angle
\[
 \alpha(s,t)=\operatorname{atan2}\!\left(
 \det(T(s,t_0),T(s,t)),\,
 T(s,t_0)\cdot T(s,t)\right)
 \in(-\pi/2,\pi/2).
\]
If $\theta_0$ lifts $T(\,\cdot\,,t_0)$, then
$\theta_0+\alpha(\,\cdot\,,t)$ lifts $T(\,\cdot\,,t)$.
Periodicity gives $\alpha(1,t)=\alpha(0,t)$, so the endpoint increment is
locally, hence globally, constant in $t$. A regular homotopy supplies such a
homotopy of normalised tangents. For the oppositely oriented curve,
$-T_\gamma(1-s)$ has lift $\theta(1-s)+\pi$, whose increment is the negative
of the original one.

For~(ii), the identity and all its polygonal consequences are exactly
Lemma~\ref{lem:rot}: polygonal rotation was defined there by the sum of
principal angles, whose integrality and continuity were proved there.

For~(iii), concatenate constant lifts on the straight pieces with the
prescribed corner lifts. Their total increment is $\sum_i\vartheta_i$, so~(ii)
proves the rounding assertion.

For the replacement assertion, use the same lift on the unchanged
complementary arc. Shifting that complementary lift at a join changes neither
its increment nor the computation; subtraction cancels it and leaves
$\Delta_c-\Delta_b$.

Finally follow the tangent path of $c$ and then the tangent path of $b$
backwards. This is a closed direction loop of increment
$\Delta_c-\Delta_b$. The maps
\[
   \nu_t(z)=\frac{A_tz}{|A_tz|}
\]
give a homotopy of that direction loop. Clause~(i) keeps its increment
constant, proving the last assertion without invoking degree.
\end{proof}

\begin{lemma}[rounding]\label{cf:lem-rounding}
Here $\rot$ is as in Lemma~\ref{lem:rot} for polygons and
Definition~\ref{cf:def-turning} for closed $C^1$ regular curves.
Let $L$ be a closed polygon with corners $q_1,\dots,q_c$, and let $D$ be an
oriented diagram whose underlying plane curve is $L$ --- $L$ together with an
over/under assignment at each of its double points. Assume the principal turns
of $L$ all exist and are \emph{nonzero}; that $L$ has finitely many double
points, all transversal, none of them a corner; and that no corner of $L$ lies
on an edge of $L$ other than the two incident to it. Then there is a \emph{clearance} $\varepsilon_0(L)>0$ such that for every
$\varepsilon\in\bigl(0,\varepsilon_0(L)\bigr)$ there is a $C^\infty$ regular
closed plane curve $L_\varepsilon$, and a diagram $D_\varepsilon$ carried by it, with
these properties.
\begin{enumerate}
\item[(a)] $L_\varepsilon$ coincides with $L$ outside the union of the discs of
radius $\varepsilon$ about the corners.
\item[(b)] Inside the disc about $q_i$ the unit tangent moves \emph{strictly}
monotonically, in the sense of $\sgn\vartheta_i$, from $\delta_{i-1}/|\delta_{i-1}|$
to $\delta_i/|\delta_i|$, sweeping an arc of length exactly $|\vartheta_i|$ and no
more, and attaining each direction of that arc at exactly one parameter.
Moreover the unit-tangent map is an \emph{immersion} on the open junction arc:
parametrized by arclength, its angular derivative is nonzero at every interior
parameter, while at the two ends it and all its derivatives vanish, the
junction meeting the straight edges flat to infinite order.

Both clauses are exported because both are read downstream, and the second does
not follow from the first: a strictly monotone angle may still have a vanishing
derivative at a parameter, and the direct count below needs the derivative
itself to be nonzero, not merely the monotonicity. "Moves monotonically" alone
is weaker still --- it is
satisfied by a junction that pauses on one direction for a whole sub-arc, and
then a count of the parameters carrying a prescribed direction is infinite. Both
hold because the profile fixed in the proof below is strictly increasing on
$(0,1)$ with $\varphi'>0$ there, and flat to infinite order at $0$ and $1$: the
tangent angle at arclength $\sigma$ along the junction of length $\ell$ is
$\theta_u+\vartheta_i\varphi(\sigma/\ell)$, whose derivative is
$\vartheta_i\varphi'(\sigma/\ell)/\ell\neq0$ for $0<\sigma<\ell$.
\item[(c)] $L_\varepsilon$ has the same double points as $L$, with the same
strands, the same over/under assignment and hence the same crossing signs and
the same writhe.
\item[(d)] $\rot(L_\varepsilon)=\rot(L)$.
\item[(e)] \emph{The disc package is returned.} There are pairwise disjoint
closed discs $D_1,\dots,D_c$, one about each corner $q_i$, such that each $D_i$
meets no edge of $L$ that is not incident to $q_i$, contains no double point of
$L$, meets the two incident edges exactly in the two sub-segments of length
$\varepsilon$ at $q_i$, and contains the whole of the modification made at
$q_i$. This clause is a conclusion of the lemma, not a package handed over by a
paragraph of its proof; its consumer asks for a disc meeting the diagram in one
embedded arc and reads it here.
\end{enumerate}
\status{proved (refereed: bench A, 2026-09-05T15:51:04Z and 2026-09-06T00:12:27Z; transcribed from CV lem:rounding (C032))}
\end{lemma}
\begin{proof}
\emph{The junction, constructed.} Since $0<|\vartheta_i|<\pi$, the unit directions
$u=\delta_{i-1}/|\delta_{i-1}|$ and $v=\delta_i/|\delta_i|$ are neither parallel
nor antiparallel. Write $A_0=q_i-\varepsilon u$ and $A_1=q_i+\varepsilon v$ for
the two points at distance $\varepsilon$ from $q_i$ along the incident edges, so
that the displacement to be realized is
\[
   A_1-A_0=\varepsilon\,(u+v),
\]
a nonzero vector along the bisector of the convex angle at $q_i$. Fix once and
for all the \emph{transition profile}
\[
   \varphi(t)=\frac{f(t)}{f(t)+f(1-t)},
   \qquad
   f(t)=\begin{cases} e^{-1/t}, & t>0,\\ 0,& t\leq0,\end{cases}
\]
which is exhibited rather than posited, and which has the four properties every
use below reads. It is $C^\infty$ on $[0,1]$: $f$ is $C^\infty$ on $\mathbb R$
with all derivatives vanishing at $0$, and the denominator
$f(t)+f(1-t)$ is positive on $[0,1]$, being $f(1)>0$ at $t=0$ and $f(1)>0$ at
$t=1$ and a sum of two nonnegative terms not both zero in between. It has
$\varphi(0)=0$ and $\varphi(1)=1$. It satisfies $\varphi(1-t)=1-\varphi(t)$, by
inspection of the formula. And it is \emph{strictly increasing on $(0,1)$}:
\[
   \varphi'(t)=\frac{f'(t)f(1-t)+f(t)f'(1-t)}{\bigl(f(t)+f(1-t)\bigr)^2}>0
   \qquad (0<t<1),
\]
both terms of the numerator being positive there, since $f>0$ and
$f'(t)=t^{-2}e^{-1/t}>0$ on $(0,1)$. Every derivative of $\varphi$ vanishes at
$0$ and at $1$, because every derivative of $f$ vanishes at $0$ and the
denominator is smooth and nonvanishing.

Strict increase is not a convenience. An earlier revision asked only that
$\varphi$ be nondecreasing, and a nondecreasing profile may be constant on a
subinterval; the tangent direction would then be constant on a whole sub-arc of
the junction, so a direction attained there would be attained at a continuum of
parameters rather than at isolated ones, and every count of the points carrying
a prescribed tangent direction --- the counts this section takes later --- would
be infinite. With $\varphi$ strictly increasing on $(0,1)$ the tangent angle
$\theta_u+\vartheta_i\varphi(t/\ell)$ is strictly monotone across the junction, and
each direction in the swept arc is attained exactly once. Let
$\theta_u$ be the angle of $u$ and, for a length $\ell>0$, let
\[
   \gamma_\ell(s)=A_0+\int_0^{s}
   \bigl(\cos(\theta_u+\vartheta_i\varphi(t/\ell)),\,
         \sin(\theta_u+\vartheta_i\varphi(t/\ell))\bigr)\,\mathrm dt,
   \qquad s\in[0,\ell].
\]
Then $\gamma_\ell$ is a unit-speed $C^\infty$ arc. Its tangent equals $u$
\emph{at} $s=0$ and $v$ \emph{at} $s=\ell$, and agrees with those constants
\emph{to infinite order} at those two parameters, every derivative of $\varphi$
vanishing there; it does not equal them on a neighbourhood, and cannot, since
the angle is strictly monotone on $(0,\ell)$.
Infinite-order agreement at the endpoint is what the junction needs: it is
exactly the condition under which replacing the corner by the arc leaves a
$C^\infty$ regular curve, since all one-sided derivatives match those of the
straight edge. The tangent angle moves strictly monotonically through exactly
$\vartheta_i$ and no more, and its derivative
$\vartheta_i\varphi'(s/\ell)/\ell$ is nonzero on $(0,\ell)$, which is (b).

Its endpoint is $A_0+\ell\,m(\varphi)$ where
$m(\varphi)=\int_0^1(\cos(\theta_u+\vartheta_i\varphi),\sin(\theta_u+\vartheta_i\varphi))\,
\mathrm dt$. The symmetry $\varphi(1-t)=1-\varphi(t)$ makes the angles
$\theta_u+\vartheta_i\varphi(t)$ and $\theta_u+\vartheta_i\varphi(1-t)$ reflections of
each other in the bisector, so $m(\varphi)$ points along $u+v$; and
$|m(\varphi)|>0$ because all the directions integrated lie in a closed angular
interval of width $|\vartheta_i|<\pi$. Taking
$\ell=\varepsilon\,|u+v|/|m(\varphi)|$ makes the endpoint exactly $A_1$.

\emph{The arc lies in the triangle $A_0q_iA_1$, in coordinates.} Put
$\alpha=|\vartheta_i|/2\in(0,\pi/2)$ and $s=\sgn\vartheta_i$, let $w$ be the unit vector
along $u+v$ and $z$ the unit vector with $\det(w,z)=1$, and take $q_i$ as
origin, writing a point as $(x_w,x_z)$ in that frame. Then
\[
   u=\cos\alpha\,w-s\sin\alpha\,z,
   \qquad
   v=\cos\alpha\,w+s\sin\alpha\,z,
\]
because $u$ and $v$ are unit vectors whose sum is along $w$ and whose angle is
$|\vartheta_i|=2\alpha$, turning from $u$ to $v$ by $\vartheta_i$. Hence
$A_0=-\varepsilon u$ has coordinates $(-\varepsilon\cos\alpha,\ s\varepsilon
\sin\alpha)$ and $A_1=\varepsilon v$ has $(\varepsilon\cos\alpha,\
s\varepsilon\sin\alpha)$, and the closed triangle with vertices $A_0,q_i,A_1$
is exactly
\begin{equation}
   T_i=\bigl\{\,(x_w,x_z):\ |x_w|\leq\cot\alpha\cdot(s\,x_z),
   \quad 0\leq s\,x_z\leq\varepsilon\sin\alpha\,\bigr\},
\label{cf:eq-triangle}
\end{equation}
the first inequality being the pair of cone conditions at the apex $q_i$ and the
second the chord through $A_0$ and $A_1$: a point $\lambda(-u)+\mu v$ with
$\lambda,\mu\geq0$ has $x_w=(\mu-\lambda)\cos\alpha$ and
$s\,x_z=(\lambda+\mu)\sin\alpha$, so $|x_w|\leq(\lambda+\mu)\cos\alpha=\cot\alpha
\cdot(s\,x_z)$ with equality exactly on the two edges, and $\lambda+\mu\leq
\varepsilon$ is the chord.

The junction arc satisfies all three inequalities of \eqref{cf:eq-triangle}. Write
its unit tangent as $T(\sigma)=\cos\beta(\sigma)\,w+s\sin\beta(\sigma)\,z$ with
$\beta(\sigma)=-\alpha+2\alpha\,\varphi(\sigma/\ell)$, which is the tangent
angle constructed above read in this frame: at $\sigma=0$ it is $u$
($\beta=-\alpha$), at $\sigma=\ell$ it is $v$ ($\beta=\alpha$), and $\beta$ is
strictly increasing with values in $[-\alpha,\alpha]$.
\emph{The two cone inequalities.} Let $n_1$ be the unit normal to $u$ with
$\langle n_1,v\rangle>0$. The line $\mathbb Ru$ contains both $q_i$ and $A_0$,
and $T_i$ lies in $\{\langle\cdot,n_1\rangle\geq0\}$. Along the arc,
$\tfrac{d}{d\sigma}\langle\gamma,n_1\rangle=\langle T,n_1\rangle\geq0$, because
$T=au+bv$ with $a,b\geq0$ --- a direction at angle $\beta\in[-\alpha,\alpha]$
from $w$ lies in the closed cone of $u$ and $v$ --- so
$\langle T,n_1\rangle=b\langle v,n_1\rangle\geq0$; and
$\langle\gamma(0),n_1\rangle=\langle-\varepsilon u,n_1\rangle=0$. Hence
$\langle\gamma,n_1\rangle\geq0$ throughout. Symmetrically, with $n_2$ the unit
normal to $v$ with $\langle n_2,A_0\rangle>0$, one has $\langle
T,n_2\rangle=a\langle u,n_2\rangle\leq0$, so $\langle\gamma,n_2\rangle$ is
nonincreasing and equals $\langle\varepsilon v,n_2\rangle=0$ at $\sigma=\ell$;
hence it is $\geq0$ throughout. Those two are the first inequality of
\eqref{cf:eq-triangle}.
\emph{The chord inequality.} $s\,x_z(\sigma)=\varepsilon\sin\alpha+\int_0^\sigma
\sin\beta(t)\,dt$, and $\int_0^\sigma\sin\beta\leq0$ for every
$\sigma\in[0,\ell]$: $\beta$ is negative on $(0,\ell/2)$ and positive on
$(\ell/2,\ell)$, and the symmetry $\varphi(1-t)=1-\varphi(t)$ makes
$\beta(\ell-t)=-\beta(t)$, so the integral decreases to its minimum at
$\ell/2$ and increases back to $0$ at $\ell$. Hence $s\,x_z\leq\varepsilon
\sin\alpha$, which is the second. The lower bound $s\,x_z\geq0$ is implied by
the two cone inequalities, the apex being the origin.

Its bisector coordinate has derivative
$dx_w/d\sigma=\cos\beta(\sigma)\geq\cos\alpha>0$.
Thus that coordinate is strictly increasing and the junction arc is
injective; no self-intersection can be hidden inside its rounding disc.

So the arc lies in $T_i$, and $T_i$ lies in the closed disc $D_i$ of radius
$\varepsilon$ about $q_i$: the disc is convex and contains the three vertices,
$|A_0-q_i|=|A_1-q_i|=\varepsilon$. An earlier revision argued this by saying the
arc ``cannot leave the region cut off by the two tangent lines'', which names
the conclusion rather than proving it and leaves the chord side unaddressed.

For (a) and (c), let $\mathcal E_i$ be the union of the closed edges
not incident to $q_i$, and define the four clearances separately:
\[
\begin{aligned}
 \eta_v&=\min_i\min_{k\ne i}|q_k-q_i|,
 &\eta_e&=\min_i\operatorname{dist}(q_i,\mathcal E_i),\\
 \eta_\ell&=\min_i|\delta_i|,
 &\eta_X&=\min_{i,\,x\text{ a double point}}|x-q_i|.
\end{aligned}
\]
Set
\[
 \varepsilon_0(L)=\tfrac13\min\{\eta_v,\eta_e,\eta_\ell,\eta_X\},
\]
and require $\varepsilon<\varepsilon_0(L)$; the index $i$ is bound by the outer
minimum in every term.
Each minimum over an empty index set is $+\infty$ --- which is how the last
one is read when $L$ has no double point. Every listed quantity is positive: the
corners are finitely many and distinct; a corner is at positive distance from a
non-incident closed edge because the edges are compact and a corner on one would
contradict the hypothesis that no corner lies on a non-incident edge; the edges have positive length; and a double point lies
in the relative interiors of two edges, so it is distinct from every corner.
With this $\varepsilon$ the discs $D_i$ are pairwise disjoint, each $D_i$ meets
no edge not incident to $q_i$, each $D_i$ contains no double point, and each
$D_i$ meets the two incident edges in the two sub-segments of length
$\varepsilon$ at $q_i$, and each contains the whole modification made at its
corner, the replacement arc lying within distance $\varepsilon$ of $q_i$. That
is clause~(e). Hence the rounding changes the curve only inside those
discs, where it meets nothing else, so no double point is created or destroyed
and the strands through each surviving double point are unchanged as arcs, with
their orientations; the over/under assignment is inherited, so the crossing signs
and the writhe are unchanged. That is (a) and (c). For (d), the total turning of
$L_\varepsilon$ is the sum of the turnings of its arcs, which is
$\sum_i\vartheta_i$ because the tangent is constant along the straight parts; and the
rotation number of a $C^1$ regular closed curve is $\frac1{2\pi}$ times its
total turning (Lemma~\ref{cf:lem-turnlift}(iii)), which is also the value assigned to $L$
by Lemma~\ref{lem:rot}.
\end{proof}

\begin{lemma}[exact negative-curl replacement]\label{cf:lem-curl}
Here $\rot$ is as in Lemma~\ref{lem:rot} for polygons and
Definition~\ref{cf:def-turning} for closed $C^1$ regular curves.
Let $F$ be a connected $C^\infty$ immersed circle in the plane --- one
component, with finitely many transverse double points and no triple points ---
given with an oriented diagram, and let $p$ be a point of $F$ at which the
tangent points in a fixed direction $u$, isolated among such points, lying in
an embedded arc of $F$ that contains no double point and along which the tangent
turns strictly positively. Then $F$ may be modified inside a disc $\Delta$
meeting the rest of the diagram only in that arc, so that the resulting
diagram $F'$
\begin{enumerate}
\item[(i)] is again such an oriented diagram, satisfies
$P_{F'}(a,z)=P_F(a,z)$, and has the same double points outside $\Delta$, with
the same signs;
\item[(ii)] has no point of $\Delta$ at which the tangent equals $u$, and exactly
one at which it equals $-u$;
\item[(iii)] has exactly one double point inside $\Delta$, and it is negative;
\item[(iv)] satisfies $\rot(F')=\rot(F)-1$ and $w(F')=w(F)-1$.
\end{enumerate}
\status{proved (refereed: bench A, 2026-09-05T15:51:04Z and 2026-09-06T00:32:58Z; transcribed from CV lem:curl (C032))}
\end{lemma}
\begin{proof}
Work first in the rational model on $-2\leq t\leq2$,
\[
b(t)=\Bigl(\tfrac{3(t^2+7)}{11},\,-3t\Bigr),\qquad
c(t)=\bigl(t^2-1,\;t-t^3\bigr).
\]
The endpoints and endpoint tangent rays agree up to positive scale:
$b(\pm2)=c(\pm2)=(3,\mp6)$ and $b'(\pm2)=\frac3{11}c'(\pm2)$. The old arc turns
strictly positively and has exactly one tangent pointing straight down, at
$t=0$:
\[
\det\bigl(b'(t),b''(t)\bigr)=\tfrac{18}{11}>0,\qquad b'(0)=(0,-3).
\]
The replacement turns strictly negatively and has exactly one tangent pointing
straight up, at $t=0$:
\[
\det\bigl(c'(t),c''(t)\bigr)=-2(1+3t^2)<0,\qquad c'(0)=(0,1).
\]
In each arc the first derivative has vanishing first component only at $t=0$, so
neither arc has any other vertical tangent; this is (ii) in the model. The
replacement runs between the same endpoint rays along the complementary,
clockwise tangent path, so its compatible lift increment is the old one minus
$2\pi$; Lemma~\ref{cf:lem-turnlift}(iii) gives (iv) for $\rot$. It has exactly one double point,
\[
c(-1)=c(1)=(0,0),\qquad c'(-1)=(-2,-2),\qquad c'(1)=(2,-2),
\]
and putting the later branch over the earlier one makes it negative, since a
crossing sign is the sign of the determinant of the overpassing tangent followed
by the underpassing tangent and
$\det\bigl(c'(1),c'(-1)\bigr)=-8<0$; that is (iii), and with it (iv) for $w$.
Being a single Reidemeister-I monogon, the replacement does not change the knot
type, which is (i).

To insert the model at $p$ we first record the projective fact it needs, in the
form the insertion argument requires.

\emph{Ordered-ray lemma.} Let $(r_1,r_2,r_3)$ and $(s_1,s_2,s_3)$ be triples of
rays from the origin such that $r_2$ lies in the open convex cone spanned by
$r_1,r_3$ and $s_2$ lies in the open convex cone spanned by $s_1,s_3$, with
$r_1,r_3$ and $s_1,s_3$ independent, and suppose in addition that for positive
representatives
\begin{equation}
   \sgn\det(\rho_1,\rho_3)=\sgn\det(\sigma_1,\sigma_3),
\label{cf:eq-orderedray}
\end{equation}
a condition independent of the representatives chosen. Then there is a linear
map $A$ with $\det A>0$ carrying each $r_k$ into $s_k$.
\emph{Proof:} if both determinants in \eqref{cf:eq-orderedray} are negative,
exchange the names $1$ and $3$ \emph{in both triples simultaneously}; this
preserves the correspondence $r_k\mapsto s_k$ and makes both determinants
positive. The linear map $A_0$ with $A_0\rho_1=\sigma_1$, $A_0\rho_3=\sigma_3$
then has positive determinant and carries $r_1,r_3$ to $s_1,s_3$. It carries the
open cone of $r_1,r_3$ onto that of $s_1,s_3$, so $A_0r_2$ is a ray $\sigma$ in
the latter cone. Writing $s_2=\mathbb R_{>0}(x\sigma_1+y\sigma_3)$ and
$\sigma=\mathbb R_{>0}(x'\sigma_1+y'\sigma_3)$ with $x,y,x',y'>0$, the map
$D=\mathrm{diag}(x/x',\,y/y')$ in the basis $(\sigma_1,\sigma_3)$ has positive
determinant, fixes $s_1$ and $s_3$, and carries $\sigma$ to $s_2$. Take
$A=DA_0$. $\square$

An earlier revision of this lemma omitted \eqref{cf:eq-orderedray} and said the
positive determinants could be arranged ``after swapping the names once if
necessary''. A swap performed in one triple only destroys the correspondence it
is supposed to produce, and the peer refuted that step; the hypothesis is now
stated and the swap is simultaneous. In the application below the condition is
not inferred from the two cones: both determinants are computed and are
positive.

\emph{A positive-turn chart, and the cuts.} Parametrise the embedded arc by
arclength as $\gamma(s)$ with $\gamma(0)=p$ and unit tangent $\gamma'(0)=u$.
Because the tangent turns strictly positively, after shrinking the chart there
is a continuous strictly increasing lift $\theta$ with
\[
   \gamma'(s)=\cos\theta(s)\,u+\sin\theta(s)\,Ju,\qquad
   \theta(0)=0,\qquad |\theta(s)|<\tfrac\pi2,
\]
where $J$ is the positive quarter turn. Put $v=-Ju$, so that $(v,u)$ is a
positively oriented orthonormal basis, and set
\[
   \xi(s)=\langle\gamma(s)-p,\,v\rangle,\qquad
   \eta(s)=\langle\gamma(s)-p,\,u\rangle,
\]
so that $\xi'(s)=-\sin\theta(s)$ and $\eta'(s)=\cos\theta(s)$. Hence $\xi$
increases strictly for $s<0$ and decreases strictly for $s>0$, with a strict
maximum $\xi(0)=0$, while $\eta$ increases throughout. For every level
sufficiently close to $\xi(0)$ from below, strict monotonicity and the
intermediate value theorem give unique cuts $s_-<0<s_+$ with
$\xi(s_-)=\xi(s_+)$, and both tend to $0$ as the level tends to $\xi(0)$. Write
$p_\pm=\gamma(s_\pm)$. Equal transverse coordinates and monotone $\eta$ give
\[
   p_+-p_-=\ell u,\qquad \ell>0,\qquad \ell\to0 .
\]
Two side facts follow at once and are what the earlier revision lacked: the
central arc $\gamma((s_-,s_+))$ lies strictly in $\xi>\xi(s_-)$, and both
retained tails lie strictly in $\xi<\xi(s_-)$.

\emph{The affine fit, quantitatively.} In the positively oriented model basis
$v_0=(-1,0)$, $u_0=(0,-1)$ one has $q=\tfrac4{11}$ and
\[
   b'(-2)\in\mathbb R_{>0}(qv_0+u_0),\qquad
   b'(0)\in\mathbb R_{>0}u_0,\qquad
   b'(2)\in\mathbb R_{>0}(-qv_0+u_0).
\]
Write the two target endpoint tangents in the basis $(v,u)$ as
$T_-=\gamma'(s_-)=av+bu$ and $T_+=\gamma'(s_+)=-cv+du$; then
$a=-\sin\theta(s_-)>0$, $c=\sin\theta(s_+)>0$ and $b,d=\cos\theta(s_\mp)>0$,
and $a,c\to0$, $b,d\to1$ as the cuts approach $p$. The two outer determinants
are $\det(qv_0+u_0,-qv_0+u_0)=2q>0$ and $\det(T_-,T_+)=ad+bc>0$, so
\eqref{cf:eq-orderedray} holds; it is computed, not inferred. Define
\[
   H=\frac{2}{\tfrac ba+\tfrac dc}=\frac{2ac}{bc+ad},\qquad
   x=\frac Hq,\qquad
   y=\frac{\tfrac ba-\tfrac dc}{q\bigl(\tfrac ba+\tfrac dc\bigr)}
    =\frac{bc-ad}{q(bc+ad)},
\]
and let $A$ be the linear map with $A u_0=u$ and $Av_0=xv+yu$. Then
\[
   A(qv_0+u_0)=\tfrac Ha\,T_-,\qquad
   A(-qv_0+u_0)=\tfrac Hc\,T_+,\qquad
   \det A=x=\frac{2ac}{q(bc+ad)}>0,
\]
each by substitution, and
\[
   1-(qy)^2=\frac{4abcd}{(bc+ad)^2}>0,
\]
so $|y|<1/q$ uniformly, while $x\to0$ as the cuts approach $p$: the map $A$
stays bounded. This is the quantitative statement that the earlier revision
replaced by an unquantified ``fixed unit-normalized part of $A$''.

Since $b(2)-b(-2)=12u_0$, put $B=\tfrac{\ell}{12}A$ and
$\Phi(z)=p_-+B\bigl(z-b(-2)\bigr)$. Then $\Phi(b(-2))=p_-$,
$\Phi(b(2))=p_-+\ell A u_0=p_-+\ell u=p_+$, and the endpoint tangent rays match
with positive scale. Because $\ell\to0$ while $A$ stays bounded, the diameter of
$\Phi(c([-2,2]))$ tends to $0$, so the cuts may be chosen so that the whole
inserted arc lies in a preassigned disc $\Delta$ about $p$ whose intersection
with the diagram is contained in the embedded chart --- the disc is fixed first
and the cuts afterwards.

\emph{Only the model's own double point is created.} In the basis $(v_0,u_0)$
the transverse coordinate of $c(t)$ is $1-t^2$, which equals $-3$ at both
endpoints, so the model's transverse excess over its endpoint chord is
$4-t^2>0$ for $-2<t<2$. Since $Bv_0=\tfrac{\ell}{12}(xv+yu)$ with $x>0$ and
$Bu_0=\tfrac{\ell}{12}u$, the physical transverse excess of the inserted arc
over the chord $[p_-,p_+]$ is exactly
\[
   \frac{\ell x}{12}\bigl(4-t^2\bigr)>0 .
\]
So the inserted interior lies strictly in $\xi>\xi(s_-)$ while both retained
tails lie strictly in $\xi<\xi(s_-)$: they meet only at the two joins. Together
with the shrinking of the previous paragraph, which keeps the inserted arc away
from every nonlocal part of the diagram, this \emph{proves} rather than assumes
that the only double point created is the model's own.

\emph{The replacement, and why the curve stays $C^\infty$ regular.} Replace the
sub-arc of $F$ between $p_-$ and $p_+$ by $\Phi\circ c$, modified near its two
endpoints as follows. At each join the two one-sided unit tangents already
agree, $\Phi c'(\pm2)$ being a positive multiple of $\Phi b'(\pm2)$ and that of
the old tangent, so the concatenation is $C^1$; it need not be $C^\infty$, and
the collar that repairs that is constructed here rather than described. An
earlier revision said only that one may ``interpolate the tangent angle by the
same device'', which does not say what curve results, and interpolating angles
moves the endpoint.

Before choosing the two collars, take disjoint small neighbourhoods of
$p_-$ and $p_+$ meeting no part of the retained old curve or fitted model
outside their respective local endpoint intervals. These neighbourhoods
exist by compactness: each endpoint has a unique preimage in the joined
curve, the only model double point is away from the endpoints, and the
transverse-excess argument has excluded a model--tail intersection.
Delete small endpoint parameter intervals; each remaining compact image
has positive distance from that endpoint. Take a ball smaller than those
distances. As its length tends to zero, the region between the two local
graphs used for a collar shrinks to the endpoint, so both collars can be
chosen inside these balls. This also separates them from the monogon and
from one another.

\emph{The collar, as a graph.} Work at the join $p_-$ first; the join at
$p_+$ is constructed after property (3) below, with its own signs printed
rather than left to "roles exchanged". In the chart coordinates
$(\xi,\eta)$ of the positive-turn chart --- $\xi=\langle\gamma-p,v\rangle$ with
$v=-Ju$, $\eta=\langle\gamma-p,u\rangle$ --- the old arc is a graph
$\xi=f(\eta)$, since $\eta'=\langle\gamma',u\rangle=\cos\theta>0$ there. Its
slope carries a sign that is
load-bearing below, so it is computed:
$\xi'=\langle\gamma',v\rangle=-\sin\theta$, whence
\[
   f'(\eta)=\frac{\xi'}{\eta'}=-\tan\theta .
\]
For $\eta<\eta(p)$ the angle $\theta$ is strictly negative, $\theta$ being
strictly increasing with $\theta(0)=0$, so $f'>0$ there --- not $f'<0$, which is
what the earlier revision's $f'=\tan\theta$ gave. Let
$\eta_-=\eta(p_-)<\eta(p)$. The inserted arc leaves $p_-$ with the same unit
tangent, so near its start it is also a graph $\xi=g(\eta)$ with
$g(\eta_-)=f(\eta_-)$ and $g'(\eta_-)=f'(\eta_-)$; and because the model turns
strictly negatively its tangent angle only decreases from $\theta(s_-)<0$, so by
the same formula
\[
   g'>f'>0
\]
on that stretch. Two consequences are used below: $g>f$ on
$(\eta_-,\eta_-+\lambda]$, the two agreeing at $\eta_-$ with $g'>f'$; and at
every point of the collar the smaller slope is $f'(\eta)$, still positive. The
bound is pointwise of necessity: $f'=-\tan\theta$ \emph{decreases} along the
collar, $\theta$ increasing, so $0<f'(\eta)<f'(\eta_-)$ there and an endpoint
bound from the left end is unavailable. Fix
$\lambda>0$ small enough that $[\eta_-,\eta_-+\lambda]$ lies both below
$\eta(p)$ and inside the stretch on which the inserted arc is a graph, and set,
on $[\eta_-,\eta_-+\lambda]$,
\[
   h(\eta)=\bigl(1-\varphi_\lambda(\eta)\bigr)f(\eta)+\varphi_\lambda(\eta)\,
   g(\eta),
   \qquad
   \varphi_\lambda(\eta)=\varphi\Bigl(\frac{\eta-\eta_-}{\lambda}\Bigr),
\]
with $\varphi$ the transition profile of Lemma~\ref{cf:lem-rounding}. The modified
curve runs along the old arc up to $\eta_-$, along the graph of $h$ across the
collar, and along the inserted arc from $\eta_-+\lambda$ on. Three properties,
each proved from the formula.

\emph{(1) It is $C^\infty$ and regular, with positive speed.} $h$ is a $C^\infty$
function, and every derivative of $\varphi$ vanishes at $0$ and $1$, so $h$
agrees with $f$ to infinite order at $\eta_-$ and with $g$ to infinite order at
$\eta_-+\lambda$: the two joins are $C^\infty$. Parametrized by $\eta$, the
collar has velocity $h'(\eta)\,v+u$ in the frame $(v,u)$, of norm
$\sqrt{1+h'(\eta)^2}\geq1$, so its speed is positive and it is regular; being a
graph, it is embedded.

\emph{(2) No tangent of the collar equals $u$, and no smallness condition is
needed.} The tangent is parallel to $u$ exactly where $h'=0$, and
\[
   h'=(1-\varphi_\lambda)f'+\varphi_\lambda g'
      +\frac{\varphi'\bigl((\eta-\eta_-)/\lambda\bigr)}{\lambda}\,
       \bigl(g-f\bigr).
\]
Every term is nonnegative, and the sum of the first two is positive: it is the
convex combination $(1-\varphi_\lambda)f'+\varphi_\lambda g'$ of two positive
numbers --- each \emph{term} separately can vanish, $\varphi_\lambda$ being $0$
at the left end and $1$ at the right --- while the third term is
$\varphi'\cdot(g-f)/\lambda\geq0$, since $\varphi'\geq0$ and $g-f\geq0$ there,
the two agreeing at $\eta_-$ with $g'>f'$. Hence, pointwise,
\[
   h'(\eta)\geq(1-\varphi_\lambda)f'(\eta)+\varphi_\lambda g'(\eta)
   \geq f'(\eta)>0 ,
\]
the middle inequality because $g'>f'$, so $h'$ never vanishes and the collar
has no tangent equal to $u$; and none equal to $-u$ either, a graph over $\eta$ having
$\eta'>0$ throughout. An earlier revision, working from the wrong slope sign,
had both $f'$ and $g'$ negative, needed a Taylor bound on $g-f$ to control the
third term, and chose $\lambda$ small to make the sum negative. With the frame
computed the third term is on the same side as the other two and no choice of
$\lambda$ is required.

\emph{(3) It creates no double point and changes no turning.} Pointwise
$f\leq h\leq g$ on the collar --- the value is a convex combination and $g\geq f$
there, which is the separation the corrected slope chain supplies --- so the
collar lies in the region between the two graphs; that region is inside the
disc $\Delta$ and inside the embedded chart, where the only strands of the
diagram are the old arc --- whose collar stretch is removed --- and the inserted
arc. The collar is embedded by~(1) and meets those two only at its endpoints,
so no double point is created in it, and the model's own double point, at
$|t|=1$, is interior to the inserted arc and untouched. For the turning: the
collar's tangent angle is continuous, stays in $(-\pi/2,0)$ --- the slope
$h'$ being positive and finite, so the tangent angle $\theta$ has
$-\tan\theta=h'>0$ with $\cos\theta>0$ --- and takes at its ends exactly the
angles of the two arcs it joins, so the lift of the tangent
along the modified curve is the concatenation of the old lifts with an
interpolating stretch of the same endpoint values; the total tangent winding is
unchanged by its insertion. An earlier revision stopped at $C^1$
and the document then carried a lemma to recover smoothness. This sentence is what the downstream consumers of this
lemma read, the front construction below taking tangent directions of the
curve, and an intermediate revision of this proof deleted it while rewriting
the tail; it is restored here.

\emph{The collar at $p_+$.} The same three properties are now proved at the
other join, with its own signs printed: here the \emph{incoming} arc is the
inserted one and the \emph{outgoing} arc is the old one, and both slopes are
negative. Let $\eta_+=\eta(p_+)>\eta(p)$. On a stretch
$[\eta_+-\lambda_+,\eta_+]$ chosen above $\eta(p)$ and inside the stretches
on which both arcs are graphs, write $\xi=g(\eta)$ for the inserted arc and
$\xi=f(\eta)$ for the old arc extended backward from $p_+$; the two agree at
$\eta_+$ with common slope $m_+=-\tan\theta_+<0$, the angle
$\theta_+=\theta(s_+)$ lying in $(0,\tfrac\pi2)$. Just left of the join the
inserted arc turns strictly negatively into its endpoint, so its angle
exceeds $\theta_+$ there, while the old arc turns positively, so its angle
lies below $\theta_+$; by $\text{slope}=-\tan\theta$, decreasing in
$\theta$,
\[
   g'(\eta)<m_+<f'(\eta)<0
   \qquad\text{on }[\eta_+-\lambda_+,\eta_+),
\]
the outer inequality $f'<0$ because $\theta>0$ above $\eta(p)$. Integrating
$g'-f'<0$ backward from the common value at $\eta_+$ gives $g>f$ on
$[\eta_+-\lambda_+,\eta_+)$. Set
\[
   h=\bigl(1-\chi_{\lambda_+}\bigr)g+\chi_{\lambda_+}f,
   \qquad
   \chi_{\lambda_+}(\eta)=\varphi\Bigl(\frac{\eta-(\eta_+-\lambda_+)}{\lambda_+}\Bigr),
\]
with the same transition profile $\varphi$: the modified curve runs along the
inserted arc up to $\eta_+-\lambda_+$, along the graph of $h$ across the
collar, and along the old arc from $\eta_+$ on, the two joins $C^\infty$ by
the flat endpoint jets of $\varphi$ exactly as in (1). For the derivative,
\[
   h'=\bigl(1-\chi_{\lambda_+}\bigr)g'+\chi_{\lambda_+}f'
      +\frac{\varphi'\bigl((\eta-\eta_++\lambda_+)/\lambda_+\bigr)}{\lambda_+}
       \,\bigl(f-g\bigr),
\]
and the sign of $f-g$ in the cutoff term is exactly what the negativity of
the two slopes alone does not control: here $f-g\leq0$ by the separation just
integrated, and $\varphi'\geq0$, so the cutoff term is nonpositive, while the
convex part is a convex combination of two negative numbers. Hence,
pointwise,
\[
   h'(\eta)\leq\bigl(1-\chi_{\lambda_+}\bigr)g'(\eta)+\chi_{\lambda_+}f'(\eta)
   \leq f'(\eta)<0 ,
\]
the middle inequality because $g'<f'$, and the bound pointwise for the
mirrored reason: $f'$ varies along the collar and no endpoint bound is
available. So $h'$ never vanishes, and this collar --- a graph over $\eta$
with $\eta'>0$ throughout --- has no tangent equal to $u$ and none equal to
$-u$; its tangent angle stays in $(0,\tfrac\pi2)$, with
$-\tan\theta=h'<0$ and $\cos\theta>0$, and takes at its ends exactly the
angles of the two arcs it joins, so the lift concatenation of (3) applies to
it verbatim and the total tangent winding is unchanged. Pointwise
$f\leq h\leq g$, the value being a convex combination and $g\geq f$ there, so
the collar lies in the region between the two graphs, inside the disc
$\Delta$ and the embedded chart, and it is embedded and meets the two arcs
only at its endpoints: no double point is created, by the argument of (3)
unchanged.

\emph{The four conclusions, separately.} (ii): $c'(t)=(2t,1-3t^2)$ is parallel to
$u_0$ only at $t=0$, where $c'(0)=-u_0$; since $A$ is invertible with
$Au_0=u$, the inserted arc has no tangent $u$ and exactly one tangent $-u$, and
the retained tails have neither, their angles lying in $(-\tfrac\pi2,\tfrac\pi2)$
away from $0$.

(iii): if $c(s)=c(t)$ then $s=\pm t$, and the second coordinate forces
$2t(1-t^2)=0$, so the only pair of distinct parameters is $\{-1,1\}$: exactly
one double point. With the later branch over the earlier one its sign is that of
$\det(c'(1),c'(-1))=-8<0$, and $\det B>0$ preserves it.

(iv) for $\rot$: Let $R$ be the rotation carrying the positive orthonormal basis
$(v_0,u_0)$ to $(v,u)$ and put $C=R^{-1}A$. Relative to $(v_0,u_0)$,
$C$ has matrix
$\left(\begin{smallmatrix}x&0\\y&1\end{smallmatrix}\right)$ with $x>0$, so
$C_t=\left(\begin{smallmatrix}1-t+tx&0\\ty&1\end{smallmatrix}\right)$
joins $I$ to $C$ inside $\mathrm{GL}^+(2,\mathbb R)$. If $R_t$ rotates
through $t$ times any fixed angle representing $R$, first $C_t$ and then
$R_tC$ give a path $A_t$ from $I$ to $A=RC$ in that group. Apply
$z\mapsto A_tz/|A_tz|$ to the closed direction loop obtained by following
the tangent path of $c$ and the tangent path of $b$ backwards.
Lemma~\ref{cf:lem-turnlift}(iii) therefore preserves their compatible endpoint
lift-increment difference. The tangent of $b$ follows the short positive arc between the
endpoint rays; the tangent of $c$ follows the complementary clockwise arc, since
$\det(c',c'')=-2(1+3t^2)<0$ and it passes through $-u_0$ at $t=0$; so their
relative winding is $-1$. The old central arc also follows the short positive
arc, its lifted sweep $\theta(s_+)-\theta(s_-)$ lying in $(0,\pi)$. Hence Lemma~\ref{cf:lem-turnlift}(iii) says that the replacement changes the
closed curve's rotation by exactly $-1$. For
$w$: the removed arc was embedded and carried no crossing, so the writhe changes
by the one new negative crossing, that is by $-1$.

\emph{The polynomial conclusion, without a knot-category detour.}
The construction just completed also proves that $F'$ remains in the lemma's
input class.  It replaces one oriented arc by one $C^\infty$ regular oriented
arc, leaves the component count unchanged, and creates exactly one transverse
double point in a disc meeting no other strand.  Thus its double points remain
finite and transverse and no triple point is created.

Let $F^{\circ}$ be the oriented diagram obtained from $F'$ by deleting the
unique monogon in $\Delta$ by a Reidemeister-I move.  The move is available:
the preceding separation argument shows that the monogon is disjoint from
all other strands, and the crossing calculation shows it is the unique new
double point.  It replaces that monogon by one crossing-free regular arc in
$\Delta$, so $F^{\circ}$ is again a one-component immersed-circle diagram with
finite transverse double points and no triple points.  Reidemeister-I
invariance of the local polynomial $P$ (Theorem~\ref{lp:core}, from Literature
input~\ref{lp:lm}) gives $P_{F'}=P_{F^{\circ}}$.

Outside $\Delta$, the diagrams $F^{\circ}$ and $F$ coincide.  Inside
$\Delta$, each has one crossing-free oriented arc joining the same two
boundary points in the same traversal direction.  Match every marked
preimage outside that local traversal interval with the identical marked
preimage of the other diagram.  The local intervals contain no marked
preimage, so this bijection preserves cyclic order, crossing pairs,
over/under positions and signs: it is a record isomorphism.  Both actual diagrams have one component and finite transverse crossings.
Lemma~\ref{rp:record-polynomial} therefore gives
$P_{F^{\circ}}=P_F$ directly, and hence $P_{F'}=P_F$.
Equality outside $\Delta$ proves the
remaining clause of~(i).
\end{proof}

\begin{theorem}[geometric carrier floor]\label{cf:thm-carrierfloor}
\begin{enumerate}
\item[(R)] Here $\rot$ is as in Lemma~\ref{lem:rot} for polygons and
Definition~\ref{cf:def-turning} for the underlying plane curve of a
smooth diagram. Let $D$ be an oriented knot diagram and $-D$ the same
diagram with every arrow reversed. Then $P_{-D}=P_D$; consequently an
oriented knot and its reverse have the same polynomial. Moreover $-D$ has
the same crossing signs and the same writhe as $D$, and the rotation of
its underlying plane curve is the negative of that of $D$.
\item[(A)] Let $L$ and $D$ satisfy the hypotheses of Lemma~\ref{cf:lem-rounding}. For every
$\varepsilon\in(0,\varepsilon_0(L))$, the construction in the proof of that
lemma returns \emph{one} curve and \emph{one} diagram at those data: the
junction inserted at the corner $q_i$ is determined by $\varepsilon$, by the
two incident unit directions and by the transition profile, which is fixed once
and for all in that proof and is not chosen per corner or per curve; the arc
length $\ell$ is then determined by the endpoint condition, and the rest of the
curve is $L$ itself. Write
\[
   \mathrm{Round}(L,D,\varepsilon)=(L_\varepsilon,D_\varepsilon)
\]
for that pair, the \emph{rounding record} of $(L,D)$ at $\varepsilon$, and call
$L_\varepsilon$ the rounded curve and $D_\varepsilon$ the rounded diagram.
\item[(B)] Here $\rot$ is as in Lemma~\ref{lem:rot}, applied to the polygon $L$.
Let $L$ be a closed polygon with nonzero edges and nonzero principal
turns, finitely many transverse double points, no triple points, no corner
at a double point and no corner on a non-incident edge, carrying a diagram
$D$. Assume, after reversing orientation if necessary, that either all
principal turns are positive, or exactly one is negative and all others
are positive. Put $R=|\rot(L)|$. Then
there are a direction $u\in S^1$ and an $\varepsilon_1>0$ such that \emph{for
every} $\varepsilon\in(0,\varepsilon_1)$ the rounded curve $L_\varepsilon$ of
the record $\mathrm{Round}(L,D,\varepsilon)$
from clause~(A) has exactly $R$ points at which its unit
tangent equals $u$ and exactly $R$ at which it equals $-u$; at each of them the
tangent crosses that direction in the positive sense.
\item[(C)] Here $\rot$ is as in Lemma~\ref{lem:rot}, applied to the polygon $L$.
Let $D$ be an oriented knot diagram all of whose crossings are positive, with
writhe $w$, whose underlying plane curve is a closed polygon $L$ with all
principal turns existing, \emph{nonzero}, and of magnitude below $\pi$, with
finitely many double points, all transversal, with no triple points, none of
them a corner of $L$, and no corner of $L$ lying on a non-incident edge; and let
$R$ be the absolute value of its Whitney rotation number.  The
finiteness/transversality and the two corner conditions are what
Lemma~\ref{cf:lem-rounding} asks of its input.  The no-triple condition is preserved by rounding and is needed below
when diagram records are compared.  These
conditions are stated here so that no generic parent polygon is assumed. Assume that, after reversing the
orientation if necessary --- which by clause~(R) changes neither
$P_D$, nor the crossings and their signs, nor $w$, nor $R$ --- either every
principal turn is positive, or exactly one is negative and every other is
positive. Then
\begin{equation}
\mindeg_a P_D(a,z)\;\geq\;1-w-R,
\label{cf:eq-floor}
\end{equation}
and the same bound holds for $f_D(a)=[z^0]P_D(a,z)$ whenever $f_D\neq0$.
\end{enumerate}
\status{new (round~3: clause~(R) and the mirror paragraph of the proof consume Lemma~\ref{lp:coefficient-transport}, F-25-54; the statement's commentary moved to Remark \texttt{rem:rounding-record} and the justification of~(R) into its proof, F-25-57, the rounding record and the determinacy clause of~(A) kept; the transverse knot named $K_T$, F-25-84; the $\rot$ scoping of~(R), F-25-88; the remainder is the round-1 transcription held by bench~A, with per-item locators: (R) the reversal and mirror deductions printed here; (A) CV thm:carrierfloor(A); (B) CV thm:carrierfloor(B); (C) CV thm:carrierfloor(C) (C032)); round~6: the spelling of the minimum $a$-degree unified with Definition~\ref{def:adeg} at three sites, F-25-173}
\end{theorem}
\begin{proof}
\emph{Clause (R).}
Reversing all arrows of a diagram preserves every crossing sign: the sign
is $\sgn\det(u_{\mathrm{over}},u_{\mathrm{under}})$, and reversing both
strands replaces both vectors by their negatives, leaving the determinant
unchanged. Hence the writhe is preserved, and reversal sends a skein triple
$(D_+,D_-,D_0)$ of diagrams to the skein triple $(-D_+,-D_-,-D_0)$, the
oriented smoothing of the reversed diagram being the reverse of the
oriented smoothing. Consider the two maps $D\mapsto P_{-D}$ and
$D\mapsto P_D$ on all oriented link diagrams, with values in
$R=\ZZ[a^{\pm1},z^{\pm1}]$. Both take the value $1$ on the crossing-free
circle, both satisfy the campaign skein of Theorem~\ref{lp:core}, the
first by the preceding sentence, and both are invariant under planar
isotopy and the three Reidemeister moves, since $P$ is
(Theorem~\ref{lp:core}) and reversal carries each such move and isotopy
to one of the same kind. Lemma~\ref{lp:coefficient-transport} identifies
the two maps: $P_{-D}=P_D$ for every diagram $D$. Since $P_D$ is the
polynomial $H_D$ of the oriented link presented by $D$
(Theorem~\ref{lp:core}), an oriented knot and its reverse have the same
polynomial. For a polygonal underlying curve every principal turn changes
sign and Lemma~\ref{cf:lem-turnlift}(ii) negates rotation. For a $C^1$
regular curve, if $\theta$ lifts its normalised tangent, then
$\theta(1-s)+\pi$ lifts the reversed tangent with the negative endpoint
increment, as in Lemma~\ref{cf:lem-turnlift}(i).

\emph{Clause (A).} The transition profile is fixed once and for all in the proof of
Lemma~\ref{cf:lem-rounding}.  At each corner, $\varepsilon$ and the two incident
unit directions determine that profile's junction, and the endpoint condition
determines its length $\ell$; outside the rounding discs the curve is $L$, and
the diagram data are inherited.  Thus the named construction returns one
specified pair at every allowed $\varepsilon$.  This asserts uniqueness of the
fixed construction's record, not uniqueness among all possible smoothings.

\emph{Clause (B).}
Reversing the orientation of $L$ negates every principal turn and the signed rotation and
preserves every self-crossing sign, so perform the allowed normalization once. In the
uniform case $\rot(L)=R\geq1$ by Lemma~\ref{lem:uniformrot}(i); in the one-dissent case the same holds by Lemma~\ref{lem:uniformrot}(ii).

Take $\varepsilon_1=\varepsilon_0(L)$, the clearance of Lemma~\ref{cf:lem-rounding}.
Choose $u\in S^1$ so that neither $u$ nor $-u$ is an edge direction and, in the
one-dissent case, neither lies in the negative turn's closed swept arc $A$, of
length $\alpha:=|\vartheta_-|<\pi$. Such $u$ exists: in the one-dissent case $A\cup(-A)$
has length at most $2\alpha<2\pi$, while in the uniform case there is no such arc; the additional
forbidden set of edge directions and their antipodes is finite. This choice depends only on $L$. Fix an argument $\phi$ of $u$.

Fix $\varepsilon\in(0,\varepsilon_1)$, write
$(L_\varepsilon,D_\varepsilon)=\mathrm{Round}(L,D,\varepsilon)$, and let $T$ be
its unit-tangent map. Parameterize the rounded traversal by $s\in[0,1]$, with
$s=0$ in the interior of a straight part; because $u,-u$ are not edge directions, its tangent
is congruent to neither $\phi$ nor $\phi+\pi$. Concatenating straight-part arguments
with the junction lifts supplied by Lemma~\ref{cf:lem-rounding}(b) gives a continuous
$\theta\colon[0,1]\to\mathbb R$, starting at an argument $\theta(0)$ of $T(0)$.
Since $T(1)=T(0)$, write $\theta(1)=\theta(0)+2\pi m$ for an integer $m$.

Each straight portion contributes zero tangent increment, and each
rounding junction contributes its principal angle by
Lemma~\ref{cf:lem-rounding}(b). Thus
$2\pi m=\sum_i\vartheta_i=2\pi\rot(L)=2\pi R$, by
Lemma~\ref{lem:rot}.
Thus
\[
 \theta(1)=\theta(0)+2\pi R.                                  \tag{*}
\]

Fix $\beta=\phi$ or $\phi+\pi$. Because $u,-u$ are not edge directions, no
level $\beta+2\pi k$ is met on a straight part or at a junction endpoint; by
the choice of $u$, none is met on the decreasing junction. Hence floor jumps
occur only in increasing junctions and never elsewhere. Lemma~\ref{cf:lem-rounding}(b)
makes each occurrence unique with positive angular derivative. Put
\[
 x=\frac{\theta(0)-\phi}{2\pi},\qquad
 y=\frac{\theta(0)-(\phi+\pi)}{2\pi}.
\]
For $u$, (*) and the integer floor identity give, separately,
\[
 \#T^{-1}(u)=\lfloor x+R\rfloor-\lfloor x\rfloor=R.
\]
For $-u$, the same argument at the other levels gives
\[
 \#T^{-1}(-u)=\lfloor y+R\rfloor-\lfloor y\rfloor=R.
\]
There are finitely many junctions, so the preimages are finite; their positive derivatives prove the assertion for arbitrary $\varepsilon$.

\emph{Clause (C).} Fix $\varepsilon\in(0,\varepsilon_1)$ with $\varepsilon_1$ the clearance of
clause~(B) and pass to the rounding record
$(L_\varepsilon,D_\varepsilon)=\mathrm{Round}(L,D,\varepsilon)$ of
clause~(A) --- one curve and one diagram, not a member
of an unnamed family. Lemma~\ref{cf:lem-rounding}(c),(d) preserves all crossings and their
oriented decorated data, the writhe and the rotation number. The rounded
diagram and the original polygonal diagram have the same named record,
including the single original component, so
Lemma~\ref{rp:record-polynomial} gives equal $P$ values. By clause~(B) there is a direction $u$ met by the tangent of
$L_\varepsilon$ exactly $R$ times, each time positively, and likewise for
$-u$. Rotate coordinates so
that $u$ is the downward vertical. The rounded diagram then has exactly $R$
downward vertical tangencies and $R$ upward ones.

Switch every crossing. The result is a diagram $\overline D$ of the mirror knot,
with the same underlying plane curve --- hence the same tangencies --- and with
writhe $-w$; every crossing of $\overline D$ is negative. Apply
Lemma~\ref{cf:lem-curl} at each of the $R$ downward vertical tangencies; they are
isolated, lie in the interiors of positively turning rounding arcs
(clause~(B)), and each rounding junction sweeps an angle of magnitude below $\pi$, so it
contains at most one of these downward tangencies. The rounding discs
returned by Lemma~\ref{cf:lem-rounding}(e) are pairwise disjoint and meet no
crossing. Choose each curl disc inside its corresponding rounding disc.
Thus all $R$ replacements have disjoint supports and leave one another
intact. Call the result $T$. By
Lemma~\ref{cf:lem-curl}, $P_T=P_{\overline D}$, $T$ has no downward vertical
tangency, and it has writhe
\[
w(T)=-w-R .
\]
Every crossing of $T$ is negative: the old ones by the switch, the $R$ new ones by
Lemma~\ref{cf:lem-curl}(iii).

We now construct the transverse lift rather than import a
front criterion.  Write the smooth regular parametrisation of $T$ as
$s\mapsto(x(s),z(s))$, put
\[
 v(s)=\sqrt{x'(s)^2+z'(s)^2},\qquad
 y_0(s)=-\frac{x'(s)}{v(s)+z'(s)},
\]
and use the convention that, in the $xz$ projection, the branch with smaller
$y$-coordinate is drawn over.  Since $T$ has no downward vertical tangency,
$v+z'>0$ everywhere: equality would force $x'=0$ and $z'<0$.  Thus $y_0$ is
smooth and periodic, and direct algebra, with no division by $x'$, gives
\[
 z'-y_0x'=z'+\frac{x'^2}{v+z'}
 =z'+\frac{v^2-z'^2}{v+z'}=v>0.                 \tag{*}
\]

It remains only to impose the crossing order.  At a crossing branch there is
no vertical tangency, so put $m=z'/x'$.  The strict inequality in~(*) permits
$y<m$ on a rightward branch and $y>m$ on a leftward branch.  For every crossing
choose admissible constants $c_{\rm O}<c_{\rm U}$ at its over- and underpassing
preimages.  Such a choice is automatic if the overpass points right: after
choosing any admissible $c_{\rm U}$, take
$c_{\rm O}<\min\{m_{\rm O},c_{\rm U}\}$.  If both point left, choose any
$c_{\rm O}>m_{\rm O}$ and then
$c_{\rm U}>\max\{m_{\rm U},c_{\rm O}\}$.  The only remaining case has
$x'_{\rm O}<0<x'_{\rm U}$.  Every crossing of $T$ is negative, and here
\[
 0>\det(t_{\rm O},t_{\rm U})
   =x'_{\rm O}x'_{\rm U}(m_{\rm U}-m_{\rm O}).
\]
Since $x'_{\rm O}x'_{\rm U}<0$, this says $m_{\rm O}<m_{\rm U}$; choose
$m_{\rm O}<c_{\rm O}<c_{\rm U}<m_{\rm U}$.  These are exactly the two allowed
half-lines and the required over/under order.

There are finitely many crossing preimages.  Around each one $s_j$, choose
pairwise disjoint cyclic parameter intervals $(\alpha_j,\beta_j)$ so small that
$z'-c_jx'>0$ throughout.  Using the explicit flat transition profile
$\varphi$ displayed in the proof of Lemma~\ref{cf:lem-rounding}, define a smooth
bump on the circle by
\[
 \psi_j(s)=
 \begin{cases}
 \varphi\!\left((s-\alpha_j)/(s_j-\alpha_j)\right),&
       \alpha_j\le s\le s_j,\\
 \varphi\!\left((\beta_j-s)/(\beta_j-s_j)\right),&
       s_j\le s\le\beta_j,\\
 0,&\text{otherwise}.
 \end{cases}
\]
All endpoint derivatives are zero, so this is smooth, equals one at $s_j$ and
has support in that interval.  Set
\[
 y=y_0+\sum_j\psi_j(c_j-y_0).
\]
The supports are disjoint.  On the $j$-th one,
\[
 z'-yx'=(1-\psi_j)(z'-y_0x')+\psi_j(z'-c_jx')>0,
\]
and outside them this is~(*).  Hence $y$ is smooth and periodic and the lift
$s\mapsto(x(s),y(s),z(s))$ is positively transverse to $dz-y\,dx$.

Finally the lift is regular because its $xz$ projection is immersed.  If two
lifted points agreed, their projected points would agree; the only distinct
parameters with that property are the two preimages of a listed double point,
and there $c_{\rm O}<c_{\rm U}$.  Thus the lift is injective, hence, by
compactness of the parameter circle, an embedded transverse knot $K_T$.
Its front and its smaller-$y$ over/under assignments are exactly the
diagram $T$.

By Theorem~\ref{fd:contact}, $\operatorname{sl}(K_T)=w(T)=-w-R$. Since
$P_T=P_{\overline D}$, Theorem~\ref{fd:contact} applied to $K_T$ with its
actual finite regular generic diagram $T$ gives
\[
\max\deg_a P_{\overline D}\;\leq\;-\operatorname{sl}(K_T)-1\;=\;w+R-1 .
\]
For every oriented link diagram $D$ let $\overline D$ be its mirror, the
diagram with every crossing switched, and put
$\mathcal M(D)=P_{\overline D}(a,z)$. Mirroring exchanges $D_+$ with $D_-$
and fixes $D_0$, so the campaign skein of Theorem~\ref{lp:core} on the
mirrored triple, multiplied by $-1$, is
\[
 a^{-1}\mathcal M(D_+)-a\mathcal M(D_-)=-z\,\mathcal M(D_0),
 \qquad \mathcal M(\bigcirc)=1.
\]
Let $\iota$ be the involutive automorphism $(a,z)\mapsto(a^{-1},-z)$ of
the Laurent ring. If a map $Q$ on oriented link diagrams satisfies these
two conditions and is invariant under planar isotopy and the three
Reidemeister moves, then $\iota\circ Q$ satisfies the campaign skein and
the unknot value with the same invariance, because $\iota$ sends the
coefficients $a^{-1},-a,-z$ to $a,-a^{-1},z$. Two such maps $Q,Q'$
therefore have $\iota\circ Q=\iota\circ Q'$ by
Lemma~\ref{lp:coefficient-transport}, hence $Q=Q'$. Both $\mathcal M$
and $D\mapsto\iota(P_D)=P_D(a^{-1},-z)$ satisfy the two conditions and are
invariant --- the first because $P$ is and mirroring carries a planar
isotopy or a Reidemeister move to one of the same kind, the second because
$P$ is --- hence $P_{\overline D}(a,z)=P_D(a^{-1},-z)$ for every diagram
$D$. As a Laurent-ring automorphism, the
substitution negates every $a$-exponent and cannot cancel a nonzero extreme
coefficient. Hence
$\max\deg_a P_{\overline D}=-\mindeg_a P_D$, and the displayed bound becomes
$\mindeg_a P_D=-\max\deg_a P_{\overline D}\geq1-w-R$, which is
\eqref{cf:eq-floor}. Finally, the support of $f_D$ is contained in the support of
$P_D$, so the same bound holds for it when it is nonzero.

\end{proof}

\begin{remark}[the rounding record]\label{rem:rounding-record}
The record $\mathrm{Round}(L,D,\varepsilon)$ of clause~(A) exists to be
quantified over. An earlier revision of this floor construction, written
when the transition profile was a per-use choice, had every consumer say
``every curve the lemma produces'' --- a quantifier over an unnamed family
that read differently in each consumer. The profile is now fixed once and
for all in the proof of Lemma~\ref{cf:lem-rounding}, so the record is one
curve and one diagram at each $\varepsilon$; and clause~(B) states its
hypotheses and produces its clearance itself rather than importing them by
``as in Lemma~\ref{cf:lem-rounding}''. This is commentary on the form of
the statement; the theorem's content is clauses~(R)--(C) (F-25-57).
\end{remark}


\begin{theorem}[corner theorem; Theorem~2 of the main text]\label{thm:floor}
Let $Q$ be a subpolygon of a decomposition of a generic polygon, and suppose
that after possibly reversing its orientation either all turns are left, or
exactly one turn is right. Then
\[
\mindeg_a H^+_Q\geq 1-m_Q-|r_Q|=d_Q,\qquad
H^+_Q\in\ZZ[a^{\pm1},z^2],\ \text{so }\mindeg_zH^+_Q\geq0.
\]
\status{new (round~5: the general-polygon branch of the hypothesis removed --- its conclusion was stated in the symbols $H^+_Q$, $m_Q$, $d_Q$, defined only for a subpolygon of a decomposition, and no proof consumed it; the broader class stays available as clause~(C) of Theorem~\ref{cf:thm-carrierfloor} in diagram symbols, F-25-164; round~4: the non-incidence hypothesis of clause~(C) supplied at the site, F-25-132; the remainder is the round-1 transcription of CV thm:carrierfloor(A),(B),(C) via Theorem~\ref{cf:thm-carrierfloor}, lem:rounding, lem:curl (C032; F-25-47), held by bench~A; the $z$-parity clause is the knot clause of registry L-1, support half printed in Theorem~\ref{lp:core}, identification half consuming Literature input~\ref{lit:homfly} (F-25-48); lit: Registry~\ref{reg:etnyre}, \ref{reg:slbound}, \ref{reg:homfly}}
\end{theorem}
\begin{proof}
For a carrier, Lemma~\ref{lem:carriers} supplies every geometric
hypothesis of clause~(C) of Theorem~\ref{cf:thm-carrierfloor} except that
no corner lies on a non-incident edge, which holds directly: a carrier's
corners are original vertices and selected crossing points; an original
vertex lies on no non-incident edge segment of $P$ by (G1), and the only
subsegments containing it are its two incident ones; a selected crossing
point lies in exactly two edge interiors by (G2) and is an endpoint, not an
interior point, of the subsegments into which the crossing sites divide
those two edges. Its actual positive
lift has exactly $m_Q$ crossings, each positive, so its writhe is $m_Q$.
The rotation in that clause is $R=|r_Q|$. Thus its lower degree bound is
$1-m_Q-|r_Q|=d_Q$ for the polynomial of that actual diagram.
Theorem~\ref{lp:core} identifies this polynomial with $H_Q^+$ and proves
its nonvanishing and its support in $\ZZ[a^{\pm1},z^2]$.
Every supported $z$ exponent is therefore nonnegative, which gives the
second bound. A zero coefficient at the corner is permitted.
\end{proof}


\begin{corollary}[the corner is the south-west slot]\label{cor:corner}
For $Q$ as in Theorem~\ref{thm:floor}, $(d_Q,0)$ is the south-west corner of
the Newton polygon of $H^+_Q$ in the sense that no monomial $a^iz^j$ of
$H^+_Q$ has $i<d_Q$ or $j<0$; $c(Q)$ is the coefficient at that slot, which
may be $0$. \status{new (printed in round 1, C013)}
\end{corollary}
\begin{proof}
Theorem~\ref{thm:floor} gives both that every nonzero monomial has
$a$-exponent at least $d_Q$ and that all $z$-exponents are nonnegative.
Consequently no exponent lies strictly below or to the left of the stated
slot. By Definition~\ref{def:C}, $c(Q)$ is exactly the coefficient
$[a^{d_Q}z^0]H^+_Q$; neither bound requires that this coefficient be nonzero.
\end{proof}


\subsection{Actual carrier blocks and elementary corner values}

\begin{definition}[Undominated blocks and their owners]\label{cb:blocks}
Fix a generic polygon $P$, its interlacement graph $G_P$, and an
independent support $S$. For a vertex set $T$ write $N_G(T)$ for the
union of its graph neighbourhoods, and put
\begin{equation}\label{cb:undominated}
 U(S)=V(G_P)\setminus\bigl(S\cup N_{G_P}(S)\bigr).
\end{equation}
The connected components of $G_P[U(S)]$ are called its \emph{blocks}.
An undominated crossing has both visits on one carrier by
Lemma~\ref{lem:carriers}; that carrier is its owner. Write $D_A$ for
the actual positive diagram of a carrier $A$, and $P_A=P_{D_A}$.
These polynomials equal the corresponding $H_A^+$ by
Theorem~\ref{lp:core}.
\end{definition}

\begin{lemma}[Actual block diagrams and the carrier product]
\label{cb:products}
Every block $H$ has one owner and admits an actual positive carrier
diagram $D_H$ with exactly its restricted named cyclic record. Its
polynomial $P_H$ is independent of the further smoothings used to
produce it. For every original carrier $A$,
\begin{equation}\label{cb:product}
 P_A=\prod_{H\text{ owned by }A}P_H,
 \qquad m_A=\sum_{H\text{ owned by }A}|H|.
\end{equation}
With no owned blocks the actual diagram has value $1$ and $m_A=0$.
\status{proved (refereed: bench A, 2026-09-06T00:35:51Z and 2026-09-06T02:03:45Z; transcribed from RC d6\_vertexedge, carrier products (C035); round~3, per clause (F-25-105): (a) one owner per block --- printed here from Lemma~\ref{lem:carriers}(iii),(iv); in RC the assertion sits with its canonical definition (def:canonical, d1\_setup: ``each such component has a unique owner carrier'', ``each undominated crossing has one owner'') and the splitting argument of def:canonical-total, and is applied in the unlabelled carrier-realization paragraph immediately after the proof of mp:blocks (d2c\_polynomial\_products, lines 471--489; the proof ends at line 469), not inside it (bench B, P-25-36/37 F105-POSITION); the round-3 locators s7:owner-restriction (an owner-restricted-set display) and s7:singleton (the singleton-cut lemma) do not carry it (F-25-105, re-checked in round~5); (b) an actual block diagram with the restricted record --- the greedy support printed here (RC mp:blocks takes such diagrams as supplied); (c) independence of $P_H$ from the further smoothings --- printed here from Lemma~\ref{rp:record-polynomial}; the product identity --- Lemma~\ref{mp:blocks} (RC d2c\_polynomial\_products mp:blocks))}
\end{lemma}
\begin{proof}
The self-crossings of a carrier are exactly its owned undominated
labels, by Lemma~\ref{lem:carriers}(iii). If two such labels interlace,
their alternating four visits cannot lie on two different carriers,
by that lemma's noncrossing assertion. Applying this along an
interlacement path puts every connected block on one owner.

Fix one nonempty block $H$. Starting at $S$, choose any label
$c\in U(S)\setminus H$ and adjoin it to the support. It is adjacent
to no selected label, so the enlarged support is independent. At any
such step, with current support $T$, the new residual set is exactly
\begin{equation}\label{cb:greedy-step}
 U(T\cup\{c\})=U(T)\setminus\bigl(\{c\}\cup N_{G_P}(c)\bigr).
\end{equation}
No original label outside $H$ in $U(S)$ is adjacent to a label of $H$,
since $H$ is a connected component of the induced graph. Thus every
step preserves all labels of $H$ and removes at least the chosen label.
The sets only shrink. After at most $|U(S)\setminus H|$ steps there is
an independent support $S_H\supseteq S$ with $U(S_H)=H$.

Lemma~\ref{lem:carriers} supplies actual geometric carriers for $S_H$.
The preceding connectedness argument puts all labels of $H$ on one
of them. No other self-crossing label survives, so its actual positive
diagram is a required $D_H$. Successor splitting retains the cyclic
order inherited from the original traversal. The surviving crossing
germs have not changed. Therefore the complete record of $D_H$ is
exactly the original record restricted to $H$, with its signs, pairing
and over/under choices. This statement concerns an actually constructed
diagram; it does not assume arbitrary Gauss subwords are realizable.
Any other choices produce the same named record, so
Lemma~\ref{rp:record-polynomial} proves independence of $P_H$.

For the original owner $A$, the connected components of its own
interlacement graph are precisely the blocks owned by it: its
self-crossing labels are their union and there are no graph edges
between different blocks. The preceding construction has supplied an
actual diagram for every one of these restricted records.
Lemma~\ref{mp:blocks} now constructs clean marked joins with the full
record of $D_A$ and proves the product in~\eqref{cb:product}.
Every retained crossing is positive and appears once, proving the
crossing-count identity. If there are no owned blocks, $D_A$ is an
actual crossing-free one-circle diagram. Its local polynomial is $1$
by the initialization in Theorem~\ref{lp:core}, and its crossing count
is zero. No empty link is evaluated.
\end{proof}

\begin{lemma}[An isolated crossing forces a zero corner coefficient]
\label{cb:singleton}
Let $A$ be a uniform carrier of $S$, and suppose one of its self-crossing
labels $c$ interlaces no other self-crossing of $A$. Then $c(A)=0$.
\status{proved (refereed: bench A, 2026-09-05T15:51:04Z, 2026-09-06T01:50:42Z and 2026-09-06T05:12:19Z; transcribed from RC d6\_vertexedge, isolated crossing (C035))}
\end{lemma}
\begin{proof}
An undominated label interlacing $c$ would have the same owner $A$ by
Lemma~\ref{cb:products}. The hypothesis therefore makes $\{c\}$ a
singleton block of $G_P[U(S)]$. Since $c\in U(S)$,
$S'=S\cup\{c\}$ is independent. Equation~\eqref{cb:greedy-step} and
$N_{G_P}(c)\cap U(S)=\varnothing$ give $U(S')=U(S)\setminus\{c\}$.

Smoothing $c$ splits $A$ into two actual daughter carriers
$\Lambda_1,\Lambda_2$. Each other block owned by $A$ lies wholly in
one of the two intervals cut by $c$. Indeed none of its chords
interlaces $c$, and a connected interlacement path cannot connect
chords internal to different intervals. Successor splitting therefore
assigns that whole block to exactly one daughter. Blocks owned by
other carriers remain there. The singleton block's actual one-crossing
diagram has polynomial $1$ by Lemma~\ref{lc:single-crossing} and
positive writhe $1$. Applying Lemma~\ref{cb:products} before and after
the split consequently gives
\begin{equation}\label{cb:singleton-products}
 P_A=P_{\Lambda_1}P_{\Lambda_2},\qquad
 m_A=m_{\Lambda_1}+m_{\Lambda_2}+1.
\end{equation}
The two new smoothing turns are nonzero and opposite as real principal
angles. Every old turn of $A$ lies on exactly one daughter. Each
daughter has at least three corners by Lemma~\ref{lem:carriers},
including its one new corner, so it has at least two inherited turns.
One daughter is uniform with the old sign and the other has exactly
one dissent. Summing principal angles gives
$r_A=r_{\Lambda_1}+r_{\Lambda_2}$. Lemma~\ref{lem:uniformrot}, with
orientation reversal if the old sign is negative, puts all three
rotation integers on that old sign's ray. Hence
\begin{equation}\label{cb:singleton-rotations}
 |r_A|=|r_{\Lambda_1}|+|r_{\Lambda_2}|.
\end{equation}

Put $f_A=[z^0]P_A$ and $g_i=[z^0]P_{\Lambda_i}$. The knot support in
Theorem~\ref{lp:core} has no negative $z$ exponents, so the first
identity in~\eqref{cb:singleton-products} gives $f_A=g_1g_2$.
If $f_A=0$, the desired coefficient is zero without assigning a degree
to it. Otherwise both $g_i$ are nonzero in the Laurent integral domain.
Theorem~\ref{thm:floor} applies to the actual uniform and one-dissent
daughters: they are subpolygons of the decomposition $S'$ of $P$
(Lemma~\ref{lem:carriers}(i)), the theorem's domain, and their turn
patterns are the ones it requires. It gives
$\mindeg_a g_i\geq 1-m_{\Lambda_i}-|r_{\Lambda_i}|$ separately.
Every exponent of the product is at least the sum of these bounds;
coefficient cancellation cannot create a smaller exponent. Thus
\begin{align}
 \mindeg_a f_A
 &\geq 2-(m_{\Lambda_1}+m_{\Lambda_2})
                -(|r_{\Lambda_1}|+|r_{\Lambda_2}|)\notag\\
 &=2-(m_A-1)-|r_A|\notag\\
 &=(1-m_A-|r_A|)+2=d_A+2.
 \label{cb:singleton-gap}
\end{align}
The coefficient at degree $d_A$ is zero, which is exactly $c(A)$.
\end{proof}

\begin{lemma}[Exterior-angle count for an embedded polygon]
\label{cb:embedded-rotation}
Every embedded regular polygon has rotation $+1$ or $-1$, according
to its traversal orientation around the bounded complementary region.
\status{proved (refereed: bench A, 2026-09-05T19:32:30Z; transcribed from the exterior-angle count as printed (C035; F-25-40))}
\end{lemma}
\begin{proof}
Lemma~\ref{lem:gauss-two-discs} supplies a finite straight triangulation
of the bounded complementary region, a disc with Euler characteristic
one. Its boundary is the polygon with possible straight subdivisions.
Write $I$ for the number of interior vertices, $B$ for boundary vertices,
$E$ for all edges and $F$ for triangular faces. A simple boundary cycle
has $B$ boundary edges. Euler's count is
\begin{equation}\label{cb:euler-disc}
 I+B-E+F=1.
\end{equation}
Every interior edge belongs to two triangles and every boundary edge
to one. Counting triangle-edge incidences gives
\begin{equation}\label{cb:edge-incidences}
 3F=2(E-B)+B=2E-B.
\end{equation}
Equation~\eqref{cb:euler-disc} first gives $E=I+B+F-1$.
Substituting this into~\eqref{cb:edge-incidences} gives
$3F=2I+B+2F-2$, and hence $F=2I+B-2$.

The sum of all Euclidean triangle angles is $\pi F$. At each
interior vertex the angles fill a full turn $2\pi$. If $\alpha_j$
is the angle of the bounded region at boundary vertex $j$, it follows
that $\sum_j\alpha_j=\pi F-2\pi I$. Substituting the preceding face
count yields $\sum_j\alpha_j=\pi(B-2)$.
Traverse the boundary with this disc on the left. Its principal turn
at $j$ is $\pi-\alpha_j$, including zero at a straight subdivision.
The region is an embedded disc, so $0<\alpha_j<2\pi$ and these are
the principal turns. Their sum is therefore
$\pi B-\pi(B-2)=2\pi$. Subdivision contributes zero and changes no
other turn. Lemma~\ref{lem:rot} gives rotation $1$ for the original
polygon. The opposite traversal negates every turn and gives $-1$.
This counts exterior angles directly; it makes no use of Whitney's
formula or of a bound depending on the number of polygon vertices.
\end{proof}

\begin{lemma}[elementary corner values]\label{lem:corner-values}
\begin{enumerate}
\item[(i)] If $Q$ is uniform and embedded ($m_Q=0$) then $|r_Q|=1$, $d_Q=0$
and $c(Q)=1$.
\item[(ii)] If $Q$ is uniform and $\{y\}$ is a crossing of $Q$ interlacing no
other crossing of $Q$, then $c(Q)=0$.
\end{enumerate}
\status{proved (refereed: bench A, 2026-09-05T18:22:57Z; transcribed from (i) via Lemma~\ref{cb:embedded-rotation}; (ii) CV thm:s7universal(D)(i) via Lemma~\ref{cb:singleton} (C035))}
\end{lemma}
\begin{proof}
(i) Lemma~\ref{cb:embedded-rotation} gives $|r_Q|=1$. Thus
$d_Q=1-0-1=0$. The actual positive diagram is crossing-free and has
polynomial $1$ by Theorem~\ref{lp:core}; its $a^0z^0$ coefficient is $1$.
(ii) The carrier $Q$ and its isolated self-crossing satisfy
Lemma~\ref{cb:singleton}, which gives the conclusion directly. The
condition is interlacement isolation; no visible planar monogon is
assumed or needed.
\end{proof}

